% Standalone research note for the hybrid-local-solver project.

\documentclass[11pt]{article}

% Common class-neutral preamble for standalone notes under manuscript/notes/.
% Note entry points all live exactly two directories below manuscript/.

\usepackage[margin=1in]{geometry}
\usepackage{amsmath,amssymb,amsthm,mathtools,bm}
\usepackage{algorithm}
\usepackage{algorithmic}
\usepackage{booktabs,tabularx,multirow,array,longtable}
\usepackage{enumitem}
\usepackage{microtype}
\usepackage[numbers,sort&compress]{natbib}
\usepackage{xcolor}
\usepackage[colorlinks=true,citecolor=blue,linkcolor=black,urlcolor=blue]{hyperref}
\usepackage[capitalize,nameinlink,noabbrev]{cleveref}

\newtheorem{theorem}{Theorem}[section]
\newtheorem{lemma}[theorem]{Lemma}
\newtheorem{proposition}[theorem]{Proposition}
\newtheorem{corollary}[theorem]{Corollary}
\newtheorem{conjecture}[theorem]{Conjecture}
\theoremstyle{definition}
\newtheorem{definition}[theorem]{Definition}
\newtheorem{assumption}[theorem]{Assumption}
\newtheorem{problem}[theorem]{Problem}
\newtheorem{observation}[theorem]{Observation}
\theoremstyle{remark}
\newtheorem{remark}[theorem]{Remark}
\theoremstyle{plain}

% Canonical mathematical notation for the active manuscript.
%
% This file preserves the recurring notation style used in the author's
% NeurIPS 2024 and 2025 papers while excluding unrelated tensor and
% machine-learning boilerplate from those archived command collections.
%
% Prefix convention:
%   \v...  bold vectors
%   \m...  bold matrices
%   \g...  calligraphic sets, graphs, and families
%   \s...  blackboard-bold spaces and number systems

\newcommand{\method}{\textsc{HybridLocalSolver}}

% Font and scalar shorthands retained from the archived manuscripts.
\newcommand{\mc}{\mathcal}
\newcommand{\R}{\mathbb{R}}
\newcommand{\E}{\mathbb{E}}
\newcommand{\G}{\mathbb{G}}
\newcommand{\N}{\mathcal{N}}
% Accuracy namespaces are semantic: do not add a bare \eps alias.
% Degree-normalized residual tolerance of classical APPR is kept distinct from
% the objective-gap and proximal fixed-point tolerances of the RPPR sections.
\newcommand{\epsappr}{\varepsilon_{\mathrm{appr}}}
\newcommand{\epsobj}{\varepsilon_{\mathrm{obj}}}
\newcommand{\epspg}{\varepsilon_{\mathrm{pg}}}
\newcommand{\epsppr}{\varepsilon_{\mathrm{ppr}}}
\newcommand{\epsin}{\varepsilon_{\mathrm{in}}}
\newcommand{\epskkt}{\varepsilon_{\mathrm{kkt}}}
\newcommand{\epsburn}{\varepsilon_{\mathrm{burn}}}
\newcommand{\epssol}{\varepsilon_{\mathrm{sol}}}
\newcommand{\epspprinst}[1]{\varepsilon_{\mathrm{ppr},#1}}
\newcommand{\trans}{\mathsf{T}}
\newcommand{\grad}{\nabla}

% Delimiters and generic helpers.
% Norms and inner products use fixed-size delimiters by default.
% Use an optional size (e.g. \norm[\big]{...}) or * for automatic sizing.
\DeclarePairedDelimiter{\norm}{\lVert}{\rVert}
\newcommand{\abs}[1]{\left\lvert #1 \right\rvert}
\DeclarePairedDelimiterX{\inner}[2]{\langle}{\rangle}{#1, #2}
\newcommand{\ip}[2]{\inner{#1}{#2}}
\newcommand{\card}[1]{\left\lvert #1 \right\rvert}
\newcommand{\set}[1]{\left\{#1\right\}}
\newcommand{\ceil}[1]{\left\lceil #1 \right\rceil}
\newcommand{\floor}[1]{\left\lfloor #1 \right\rfloor}
\newcommand{\comp}[1]{\overline{#1}}
\newcommand{\conj}[1]{\overline{#1}}
\newcommand{\mean}[1]{\overline{#1}}
\newcommand{\bigO}{\mathcal{O}}
\newcommand{\softO}{\widetilde{\mathcal{O}}}
\newcommand{\softOmega}{\widetilde{\Omega}}
\newcommand{\pos}[1]{\left[#1\right]_{+}}

% Bold lowercase vectors.
\newcommand{\va}{\bm{a}}
\newcommand{\vb}{\bm{b}}
\newcommand{\vc}{\bm{c}}
\newcommand{\vd}{\bm{d}}
\newcommand{\ve}{\bm{e}}
\newcommand{\vf}{\bm{f}}
\newcommand{\vg}{\bm{g}}
\newcommand{\vh}{\bm{h}}
\newcommand{\vi}{\bm{i}}
\newcommand{\vj}{\bm{j}}
\newcommand{\vk}{\bm{k}}
\newcommand{\vl}{\bm{l}}
\newcommand{\vm}{\bm{m}}
\newcommand{\vn}{\bm{n}}
\newcommand{\vo}{\bm{o}}
\newcommand{\vp}{\bm{p}}
\newcommand{\vq}{\bm{q}}
\newcommand{\vr}{\bm{r}}
\newcommand{\vs}{\bm{s}}
\newcommand{\vt}{\bm{t}}
\newcommand{\vu}{\bm{u}}
\newcommand{\vv}{\bm{v}}
\newcommand{\vw}{\bm{w}}
\newcommand{\vx}{\bm{x}}
\newcommand{\vy}{\bm{y}}
\newcommand{\vz}{\bm{z}}

% Bold Greek vectors and special vectors.
\newcommand{\valpha}{\bm{\alpha}}
\newcommand{\vbeta}{\bm{\beta}}
\newcommand{\vdelta}{\bm{\delta}}
\newcommand{\vepsilon}{\bm{\epsilon}}
\newcommand{\veta}{\bm{\eta}}
\newcommand{\vgamma}{\bm{\gamma}}
\newcommand{\vlambda}{\bm{\lambda}}
\newcommand{\vmu}{\bm{\mu}}
\newcommand{\vomega}{\bm{\omega}}
\newcommand{\vphi}{\bm{\phi}}
\newcommand{\vpi}{\bm{\pi}}
\newcommand{\vpsi}{\bm{\psi}}
\newcommand{\vrho}{\bm{\rho}}
\newcommand{\vsigma}{\bm{\sigma}}
\newcommand{\vtheta}{\bm{\theta}}
\newcommand{\vxi}{\bm{\xi}}
\newcommand{\vell}{\bm{\ell}}
\newcommand{\vDelta}{\bm{\Delta}}
\newcommand{\vGamma}{\bm{\Gamma}}
\newcommand{\vLambda}{\bm{\Lambda}}
\newcommand{\vPhi}{\bm{\Phi}}
\newcommand{\vSigma}{\bm{\Sigma}}
\newcommand{\vzero}{\bm{0}}
\newcommand{\vone}{\bm{1}}
\newcommand{\one}{\mathbf{1}}
\newcommand{\zero}{\vzero}
\newcommand{\eunit}[1]{\bm{e}_{#1}}
\newcommand{\vDeltax}{\bm{\Delta}\bm{x}}

% Bold uppercase matrices.
\newcommand{\mA}{\bm{A}}
\newcommand{\mB}{\bm{B}}
\newcommand{\mC}{\bm{C}}
\newcommand{\mD}{\bm{D}}
\newcommand{\mE}{\bm{E}}
\newcommand{\mF}{\bm{F}}
\newcommand{\mG}{\bm{G}}
\newcommand{\mH}{\bm{H}}
\newcommand{\mI}{\bm{I}}
\newcommand{\mJ}{\bm{J}}
\newcommand{\mK}{\bm{K}}
\newcommand{\mL}{\bm{L}}
\newcommand{\mM}{\bm{M}}
\newcommand{\mN}{\bm{N}}
\newcommand{\mO}{\bm{O}}
\newcommand{\mP}{\bm{P}}
\newcommand{\mQ}{\bm{Q}}
\newcommand{\mR}{\bm{R}}
\newcommand{\mS}{\bm{S}}
\newcommand{\mT}{\bm{T}}
\newcommand{\mU}{\bm{U}}
\newcommand{\mV}{\bm{V}}
\newcommand{\mW}{\bm{W}}
\newcommand{\mX}{\bm{X}}
\newcommand{\mY}{\bm{Y}}
\newcommand{\mZ}{\bm{Z}}

% Bold Greek matrices retained from the graph papers.
\newcommand{\mBeta}{\bm{\beta}}
\newcommand{\mDelta}{\bm{\Delta}}
\newcommand{\mGamma}{\bm{\Gamma}}
\newcommand{\mLambda}{\bm{\Lambda}}
\newcommand{\mOmega}{\bm{\Omega}}
\newcommand{\mPhi}{\bm{\Phi}}
\newcommand{\mPi}{\bm{\Pi}}
\newcommand{\mSigma}{\bm{\Sigma}}
\newcommand{\mTheta}{\bm{\Theta}}

% Calligraphic symbols.
\newcommand{\gA}{\mathcal{A}}
\newcommand{\gB}{\mathcal{B}}
\newcommand{\gC}{\mathcal{C}}
\newcommand{\gD}{\mathcal{D}}
\newcommand{\gE}{\mathcal{E}}
\newcommand{\gF}{\mathcal{F}}
\newcommand{\gG}{\mathcal{G}}
\newcommand{\gH}{\mathcal{H}}
\newcommand{\gI}{\mathcal{I}}
\newcommand{\gJ}{\mathcal{J}}
\newcommand{\gK}{\mathcal{K}}
\newcommand{\gL}{\mathcal{L}}
\newcommand{\gM}{\mathcal{M}}
\newcommand{\gN}{\mathcal{N}}
\newcommand{\gO}{\mathcal{O}}
\newcommand{\gP}{\mathcal{P}}
\newcommand{\gQ}{\mathcal{Q}}
\newcommand{\gR}{\mathcal{R}}
\newcommand{\gS}{\mathcal{S}}
\newcommand{\gT}{\mathcal{T}}
\newcommand{\gU}{\mathcal{U}}
\newcommand{\gV}{\mathcal{V}}
\newcommand{\gW}{\mathcal{W}}
\newcommand{\gX}{\mathcal{X}}
\newcommand{\gY}{\mathcal{Y}}
\newcommand{\gZ}{\mathcal{Z}}

% Blackboard-bold symbols.
\newcommand{\sA}{\mathbb{A}}
\newcommand{\sB}{\mathbb{B}}
\newcommand{\sC}{\mathbb{C}}
\newcommand{\sD}{\mathbb{D}}
\newcommand{\sE}{\mathbb{E}}
\newcommand{\sF}{\mathbb{F}}
\newcommand{\sG}{\mathbb{G}}
\newcommand{\sH}{\mathbb{H}}
\newcommand{\sI}{\mathbb{I}}
\newcommand{\sJ}{\mathbb{J}}
\newcommand{\sK}{\mathbb{K}}
\newcommand{\sL}{\mathbb{L}}
\newcommand{\sM}{\mathbb{M}}
\newcommand{\sN}{\mathbb{N}}
\newcommand{\sO}{\mathbb{O}}
\newcommand{\sP}{\mathbb{P}}
\newcommand{\sQ}{\mathbb{Q}}
\newcommand{\sR}{\mathbb{R}}
\newcommand{\sS}{\mathbb{S}}
\newcommand{\sT}{\mathbb{T}}
\newcommand{\sU}{\mathbb{U}}
\newcommand{\sV}{\mathbb{V}}
\newcommand{\sW}{\mathbb{W}}
\newcommand{\sX}{\mathbb{X}}
\newcommand{\sY}{\mathbb{Y}}
\newcommand{\sZ}{\mathbb{Z}}

% Frequently used operators.
\DeclareMathOperator{\Cov}{Cov}
\DeclareMathOperator{\Var}{Var}
\DeclareMathOperator{\diag}{diag}
\DeclareMathOperator{\nei}{Nei}
\DeclareMathOperator{\nnz}{nnz}
\DeclareMathOperator{\pr}{pr}
\DeclareMathOperator{\prox}{prox}
\DeclareMathOperator{\push}{push}
\DeclareMathOperator{\res}{res}
\DeclareMathOperator{\rel}{Rel}
\DeclareMathOperator{\relu}{ReLU}
\DeclareMathOperator{\sign}{sign}
\DeclareMathOperator{\supp}{supp}
\DeclareMathOperator{\tr}{tr}
\DeclareMathOperator{\Tr}{Tr}
\DeclareMathOperator{\vol}{vol}
\DeclareMathOperator{\work}{work}
\DeclareMathOperator{\Work}{Work}
\DeclareMathOperator*{\argmax}{arg\,max}
\DeclareMathOperator*{\argmin}{arg\,min}

% Project-wide names and claim-status labels used by standalone research notes.
% Mathematical notation belongs in math_commands.tex.

% Algorithms and solver families.
\newcommand{\AESP}{\textsc{AESP}}
\newcommand{\APPR}{\textsc{APPR}}
\newcommand{\ASPR}{\textsc{ASPR}}
\newcommand{\FISTA}{\textsc{FISTA}}
\newcommand{\ISTA}{\textsc{ISTA}}
\newcommand{\LOCAPPR}{\textsc{LocAPPR}}
\newcommand{\LOCGD}{\textsc{LocGD}}
\newcommand{\LOCSOR}{\textsc{LocSOR}}
\newcommand{\LOCCH}{\textsc{LocCH}}
\newcommand{\LOCHB}{\textsc{LocHB}}
\newcommand{\LocProxGS}{\textsc{LocProxGS}}
\newcommand{\Hybrid}{\textsc{AESP--LocSOR}}

% Claim classifications required by the repository research-note protocol.
\newcommand{\statussource}{\textbf{Source result}}
\newcommand{\statusproved}{\textbf{Proved here}}
\newcommand{\statusconditional}{\textbf{Conditional}}
\newcommand{\statusconjecture}{\textbf{Open / conjectural}}
\newcommand{\statusopen}{\textbf{Open}}
\newcommand{\statusempirical}{\textbf{Empirical}}
\newcommand{\statusobservation}{\textbf{Empirical observation}}
\newcommand{\statusrefuted}{\textbf{Refuted route}}

% Compatibility names used by the older complete-note presentation layer.
\newcommand{\Status}[2]{\textbf{#1:} #2}
\newcommand{\SourceStatus}{\textnormal{\textsc{Source result}}}
\newcommand{\ProvedStatus}{\textnormal{\textsc{Proved here}}}
\newcommand{\ConditionalStatus}{\textnormal{\textsc{Conditional}}}
\newcommand{\ComputedStatus}{\textnormal{\textsc{Computed}}}
\newcommand{\RefutedStatus}{\textnormal{\textsc{Refuted route}}}
\newcommand{\OpenStatus}{\textnormal{\textsc{Open}}}



\title{Removing the Repeated-Active-Set Factor:\\
Nested SDD Corrections and Persistent Local Responses}
\author{Baojian Zhou\\
\small Working research note for \texttt{hybrid-local-solver}}
\date{August 2026; revised September 7, 2026}

\begin{document}
\maketitle

\begin{abstract}
Wei and Yang's growing-active-set PageRank method solves a nearly-linear
SDD system on every reached face, giving repeated-support work
$\widetilde O(|S^\star|\vol(S^\star))$. This note studies what is needed
to reuse the nested systems. Exact Schur corrections and a telescoping
energy identity supply warm-start information, but they do not pay for
repeatedly writing old coordinates or searching all boundary conditions.
Endpoint paths admit a linear-work realization using scalar elimination
records and one final reverse pass.

A geometric-publication reduction separates response production from
boundary delivery. Original-coordinate publications have near-linear total
delivery work and admit a certified ACL completion; finding all due
publications remains an explicit producer obligation on general graphs.
Several concrete constructions narrow that obligation. Persistent affine
response curves give an exact $O(n\log^2 n)$ solve on a supplied tree.
A local solver on cycles with arbitrary pendant-leaf counts reaches
$\widetilde O(1/\epsappr)$ ACL work. For a cycle seed and tree attachments
of size at most $q$, a locally discovered changing-response window gives
$\widetilde O(q^3/\epsappr)$ work. A separate local construction uses an
implemented planar reporter to avoid inactive attachment scans when the
positive obstacle support contains at most one nonseed branching vertex.

Finally, scalar physical-flux queues support a growing retained branching
core and give $\widetilde O((1+r^2)/\epsappr)$ ACL work, with all searches
and dense inverse updates charged. Here $r$ counts retained vertices in the
returned support. Constant-relative rounding of accumulated Schur loads
fails the original residual guarantee; physical-flux bands avoid that
failure. The explicit inverse still incurs cubic writes on easy balanced
trees, identifying a remaining representation cost. These constructions
can also peel settled core vertices, replacing the permanent count by the
maximum live core. Completed and interior tree supports show respectively
why admission order can help and why it does not guarantee a small core.
An ordered path-cluster construction now removes the growing-core factor on
arbitrary trees. Persistent projective upper chains provide exact cluster
callbacks; a checked top-tree source interface pays for online discovery,
all hierarchy changes and final recovery, giving
$O(\operatorname{cvol}(U)\log^4(2+\operatorname{cvol}(U)))$ local word work.
The balancing algorithm is imported, while its application callbacks are
implemented and exactly audited. A uniform negative-load shift then allows
one locally detected cycle to be restored at two fixed ports, with one paid
re-rooting. This gives the same local work bound on arbitrary unicyclic
graphs, for any seed and without a bound on attachment sizes.
On arbitrary graphs, a partition into two-port tree pieces gives the scoped
bound $\widetilde O((1+r^2+q)/\epsappr)$, where $r$ is active cycle rank
and $q$ is locally revealed cycle rank. Selected Green queries permit
old-port promotion, a paid inverse border and temporary-port elimination;
the current coarse inverse is updated without repeated factorization.
Using these component reporters to produce geometric value publications
then removes the revealed-rank term for ACL approximation, giving
$\widetilde O((1+r^2)/\epsappr)$ work. This policy may stop before the
exact obstacle optimum; its final original residual is certified directly.
The results are proof drafts awaiting
independent review. Arbitrary-graph OP3 remains open.
\end{abstract}

\clearpage
\tableofcontents

\section{Scope and claim ledger}
\label{sec:scope}

% Canonical source-aligned PageRank/RPPR reference formulation.
%
% This fragment is intentionally heading-free at the section level so that the
% active paper and every standalone note can input it. It is a manuscript
% reference convention, not an implementation-wide residual decision.
% Primary source: Fountoulakis and Martinez-Rubio, arXiv:2602.21138v2,
% PDF p. 3, Section 3 and equation (RPPR).

\subsection{Common source-aligned PageRank and RPPR model}
\label{subsec:shared-source-aligned-problem}

All documents in this manuscript workspace use the following reference
language. It follows the plain italic notation of the comparison paper:
\(x,s,A,D,Q,I\) denote column vectors or matrices as determined by their
displayed domains. This is the scoped exception to the project's usual
bold-vector and bold-matrix typography. A note may impose a stronger
assumption after this block, but it must not redefine these objects.

Let \(G=(V,E)\) be a finite, simple, undirected, unweighted graph without
isolated vertices, let \(V=[n]:=\{1,\ldots,n\}\), and let
\(A\in\{0,1\}^{n\times n}\) be its symmetric adjacency matrix. For
\(i\in V\), write
\[
    \mathcal N(i):=\{j\in V:\{i,j\}\in E\},
    \qquad
    d_i:=|\mathcal N(i)|,
    \qquad
    D:=\diag(d_1,\ldots,d_n).
\]
For \(S\subseteq V\), define
\[
    \mathcal N(S):=\bigcup_{i\in S}\mathcal N(i),
    \qquad
    \partial S:=\mathcal N(S)\setminus S,
    \qquad
    \operatorname{Ext}(S):=V\setminus(S\cup\partial S),
\]
\begin{equation}
    \vol(S):=\sum_{i\in S}d_i.
    \label{eq:shared-volume}
\end{equation}
The support of a vector is
\(\supp(x):=\{i\in[n]:x_i\neq0\}\). Matrix inequalities use the Loewner
order, and all unqualified vector inequalities are coordinatewise.

The seed is a column distribution
\begin{equation}
    s\in\R_+^n,
    \qquad
    \inner{\one}{s}=1.
    \label{eq:shared-seed}
\end{equation}
The single-seed specialization is always written \(s=e_v\), where \(v\in V\)
is the seed vertex. Thus \(s\) is never used as a vertex label. For
\(\alpha\in(0,1]\), define
\begin{equation}
\begin{aligned}
    \mathcal L&:=I-D^{-1/2}AD^{-1/2},\\
    Q&:=\alpha I+\frac{1-\alpha}{2}\mathcal L
      =\frac{1+\alpha}{2}I
       -\frac{1-\alpha}{2}D^{-1/2}AD^{-1/2},\\
    b&:=\alpha D^{-1/2}s.
\end{aligned}
    \label{eq:shared-pagerank-matrices}
\end{equation}
Since \(0\preceq\mathcal L\preceq2I\),
\(\alpha I\preceq Q\preceq I\). The shared smooth PageRank quadratic is
\begin{equation}
    f(x):=\frac12\inner{x}{Qx}-\inner{b}{x},
    \qquad
    \nabla f(x)=Qx-b.
    \label{eq:shared-pagerank-objective}
\end{equation}
Its unique unregularized minimizer and associated lazy personalized PageRank
vector are
\begin{equation}
    x^0:=Q^{-1}b,
    \qquad
    \pi:=D^{1/2}x^0.
    \label{eq:shared-ppr-solution}
\end{equation}

For a regularization parameter \(\rho>0\), reserve \(g_\rho\) for the
nonsmooth weighted \(\ell_1\) term and define
\begin{equation}
    g_\rho(x):=\alpha\rho\norm{D^{1/2}x}_1,
    \qquad
    F_\rho(x):=f(x)+g_\rho(x),
    \qquad
    x^\star(\rho):=\argmin_{x\in\R^n}F_\rho(x).
    \label{eq:shared-rppr-objective}
\end{equation}
When \(\rho\) is fixed, \(x^\star\) abbreviates \(x^\star(\rho)\), never the
unregularized point \(x^0\). Write
\[
    S^\star(\rho):=\supp(x^\star(\rho)),
    \qquad
    I^\star(\rho):=[n]\setminus S^\star(\rho).
\]
The source records \(x^\star(\rho)\geq0\),
\(\vol(S^\star(\rho))\leq1/\rho\), and the coordinatewise conditions
\begin{equation}
    \nabla_i f(x^\star(\rho))
    \in
    \begin{cases}
        \{-\alpha\rho\sqrt{d_i}\}, & x_i^\star(\rho)>0,\\
        [-\alpha\rho\sqrt{d_i},0], & x_i^\star(\rho)=0.
    \end{cases}
    \label{eq:shared-rppr-kkt}
\end{equation}

The common computational unit is one repeated adjacency-list scan: scanning
coordinate \(i\) costs \(d_i\), and scanning a set \(S\) costs \(\vol(S)\).
Iteration count and wall-clock time do not replace this degree-weighted work
measure.

Finally, accuracy symbols remain distinct. The APPR threshold is
\(\epsappr\); \(\epsobj\) denotes an objective-gap target;
\(\epspg\) denotes the source experiment's proximal fixed-point tolerance; and
\(\epsppr\) may denote a document-scoped degree-normalized PPR error target.
No equality or conversion among these quantities is assumed. In particular,
this shared model does not resolve the repository-wide residual decision.


The shared project system uses the lazy symmetric matrix $Q$.  Wei and Yang
instead use a degree-form, nonlazy PageRank matrix.  We keep the two
teleportation parameters distinct until \cref{sec:mapping}.  All source claims
below refer to arXiv \texttt{2608.16339v1}; all complexity comparisons charge
at least one word operation per explicitly written vector coordinate and one
degree unit when an adjacency list is scanned.

For every exposed vertex set $U$, define the note-scoped charged volume
\begin{equation}
 \operatorname{cvol}(U):=\sum_{u\in U}(1+d_u).
 \label{eq:charged-volume}
\end{equation}
Because the shared graph is connected and has at least two vertices,
$\vol(U)\leq\operatorname{cvol}(U)\leq2\vol(U)$.  This quantity charges one
stored/output word per active coordinate in addition to its adjacency scan.

\begin{center}
\small
\begin{tabularx}{\textwidth}{@{}p{0.19\textwidth}X@{}}
\toprule
Status & Statement \\
\midrule
\statussource & The growing-active-set method returns an ACL
$\varepsilon$-approximation in $\widetilde O(1/\varepsilon^2)$ time and an
additive RPPR approximation in
$\widetilde O(|S^\star|\vol(S^\star))$ time.  It solves every nested SDD
system from scratch and explicitly names incremental reuse as open. \\
\statusproved & Successive exact restricted solutions satisfy an exact
block-Schur correction formula.  Their squared correction energies telescope;
along a source-valid PageRank trajectory their total is at most the source
teleportation parameter. \\
\statusrefuted & A conventional warm start, followed by a full active-set
matrix pass or explicit active-vector output at every expansion, does not
remove the repeated factor.  Endpoint paths give quadratic repeated-prefix
work. \\
\statusproved & On endpoint paths, append-only tridiagonal elimination
computes every exact boundary gate in constant arithmetic per new vertex and
recovers the final vector in one reverse pass.  Total charged work is linear
in the final active volume. \\
\statusconditional & On arbitrary graphs, the Wei--Yang correctness proof
immediately accepts an implicit incremental solve-and-boundary oracle.  If its
total update, query, and final-output work is near-linear in the final active
volume, the repeated factor disappears. \\
\statusopen & No such graph-uniform local oracle is proved here.  A block
Schur correction can be dense, while boundary residuals depend on the changed
old coordinates.  Energy telescoping alone does not bound the work needed to
represent or query those changes. \\
\bottomrule
\end{tabularx}
\end{center}

The conclusions do not assert a new lower bound for all local graph-access
algorithms.  In fact, the path theorem proves the opposite: the quadratic
path cost belongs to a materialized repeated-solve representation, not to the
PageRank problem itself.

\section{Mapping the source system to the shared RPPR system}
\label{sec:mapping}

Let $\alpha\in(0,1)$ denote the shared lazy PageRank parameter and put
\begin{equation}
 a:=\frac{1+\alpha}{2},
 \qquad
 \beta:=\frac{1-\alpha}{2},
 \qquad
 \bar\alpha:=\frac{\alpha}{a}=\frac{2\alpha}{1+\alpha}.
 \label{eq:alpha-map}
\end{equation}
For $x=D^{1/2}z$,
\begin{equation}
 D^{1/2}QD^{1/2}=aD-\beta A
 =a\bigl(D-(1-\bar\alpha)A\bigr).
 \label{eq:matrix-map}
\end{equation}
Dividing the shared RPPR objective by $a$ therefore gives
\begin{equation}
 \frac12 z^\top L_{\bar\alpha}z
 -\bar\alpha s^\top z
 +\bar\alpha\rho\,d^\top |z|,
 \qquad
 L_{\bar\alpha}:=D-(1-\bar\alpha)A.
 \label{eq:source-degree-form}
\end{equation}
Thus the source regularization parameter $\rho$ is unchanged, its variable is
the degree-normalized variable $z=D^{-1/2}x$, and its teleportation parameter
is $\bar\alpha$.  In particular,
$\bar\alpha=\Theta(\alpha)$ uniformly on $(0,1)$, so polynomial comparisons
in $1/\alpha$ are preserved.

For the rest of the note, write $M:=L_{\bar\alpha}$ and define the source
residue
\begin{equation}
 r(z):=s-\bar\alpha^{-1}Mz.
 \label{eq:source-residue}
\end{equation}
For an active set $U$ containing the seed support and an internal residue
level $\lambda>0$, the exact restricted state is
\begin{equation}
 z[U]_U=M_{UU}^{-1}c_U,
 \qquad
 z[U]_{V\setminus U}=0,
 \qquad
 c:=\bar\alpha(s-\lambda d).
 \label{eq:restricted-state}
\end{equation}
It has $r(z[U])_U=\lambda d_U$.  The exact version of the source gate admits
\begin{equation}
 T(U):=\{v\in\partial U:r(z[U])(v)>(\lambda+\kappa)d_v\},
 \label{eq:exact-gate}
\end{equation}
where $\kappa>0$ is the activation gap.

\section{What the 2026 source proves}
\label{sec:source}

Wei and Yang solve \cref{eq:restricted-state} approximately at every round,
scan the boundary, add all coordinates whose approximate residue exceeds the
gate, and repeat.  Their energy tolerance is chosen small enough to make
active coordinates and boundary residues accurate coordinatewise.  The
M-matrix comparison then gives monotone support growth; in the RPPR setting,
every admitted coordinate lies in the exact optimal support $S^\star$.

Their Theorem~1.2 gives an ACL $\varepsilon$-approximate PageRank vector of
support volume at most $2/\varepsilon$ in
$\widetilde O(1/\varepsilon^2)$ time.  Their Theorem~1.3 gives a
$\xi$-additive RPPR minimizer whose output is also an ACL
$2\rho$-approximation in
\begin{equation}
 O\!\left(
 |S^\star|\,\widetilde{\vol}(S^\star)
 \operatorname{polylog}\frac{1}{\bar\alpha\rho\xi\delta}
 +|S^\star|\vol(S^\star)
 \right)
 =\widetilde O(|S^\star|\vol(S^\star)).
 \label{eq:source-rppr-work}
\end{equation}
The factor $|S^\star|$ comes from at most one fresh solve and active-boundary
scan per admitted batch, with only the guarantee that a nonterminal batch is
nonempty.  The source conclusion explicitly proposes reusing the nested SDD
systems to approach $\widetilde O(1/\varepsilon)$ for ACL PageRank.  The next
sections determine what nestedness gives for free and what it does not.

\section{Exact nested correction and energy telescoping}
\label{sec:telescope}

Define the fixed quadratic
\begin{equation}
 \phi(z):=\frac12 z^\top Mz-c^\top z.
 \label{eq:fixed-quadratic}
\end{equation}
The vector $z[U]$ is its minimizer over vectors supported in $U$.

\begin{theorem}[Block-Schur correction]
\label{thm:block-schur}
Let $U^+=U\cup T$, where $T\cap U=\emptyset$, and put
$z^-=z[U]$, $z^+=z[U^+]$, and $\Delta=z^+-z^-$.  Order the coordinates as
$U,T$ and define
\begin{equation}
 K_T:=M_{TT}-M_{TU}M_{UU}^{-1}M_{UT},
 \qquad
 h_T:=r(z^-)_T-\lambda d_T.
 \label{eq:schur-data}
\end{equation}
Then $K_T$ is positive definite and
\begin{align}
 \Delta_T&=\bar\alpha K_T^{-1}h_T,
 &
 \Delta_U&=-M_{UU}^{-1}M_{UT}\Delta_T,
 \label{eq:block-correction}\\
 \phi(z^-)-\phi(z^+)
 &=\frac12\|\Delta\|_{M_{U^+U^+}}^2
 =\frac{\bar\alpha^2}{2}h_T^\top K_T^{-1}h_T.
 \label{eq:block-energy}
\end{align}
If $h_T>0$ entrywise, then $\Delta>0$ on every coordinate connected through
the active component to $T$.
\end{theorem}

\begin{proof}
The old restricted optimality equations give
$(Mz^--c)_U=0$.  On $T$, \cref{eq:source-residue} gives
$(Mz^--c)_T=-\bar\alpha h_T$.  Subtracting the optimality equation for
$z^+$ therefore yields
\begin{equation}
 \begin{bmatrix}M_{UU}&M_{UT}\\M_{TU}&M_{TT}\end{bmatrix}
 \begin{bmatrix}\Delta_U\\\Delta_T\end{bmatrix}
 =\bar\alpha\begin{bmatrix}0\\h_T\end{bmatrix}.
 \label{eq:block-system}
\end{equation}
Eliminating the first block proves \cref{eq:block-correction}.  Positive
definiteness of $M_{U^+U^+}$ implies positive definiteness of its Schur
complement $K_T$.  Completing the square around the minimizer $z^+$ gives
the first equality in \cref{eq:block-energy}; substituting
\cref{eq:block-correction} gives the second.  Finally, $M_{U^+U^+}$ is a
nonsingular Stieltjes matrix, so its inverse is entrywise nonnegative and is
strictly positive on each connected irreducible block.  Apply this fact to
\cref{eq:block-system}.
\end{proof}

\begin{corollary}[Correction energies telescope]
\label{cor:energy-telescope}
For nested active sets $U_0\subset U_1\subset\cdots\subset U_K$,
\begin{equation}
 \sum_{k=0}^{K-1}
 \|z[U_{k+1}]-z[U_k]\|_{M_{U_{k+1}U_{k+1}}}^2
 =2\bigl(\phi(z[U_0])-\phi(z[U_K])\bigr).
 \label{eq:telescope}
\end{equation}
If the trajectory has $z[U_K]\geq0$ and nonnegative PageRank residue, then
\begin{equation}
 \sum_{k=0}^{K-1}
 \|z[U_{k+1}]-z[U_k]\|_{M_{U_{k+1}U_{k+1}}}^2
 \leq \bar\alpha.
 \label{eq:telescope-bound}
\end{equation}
\end{corollary}

\begin{proof}
Summing \cref{eq:block-energy} proves \cref{eq:telescope}.  The zero vector
is feasible for every restricted problem, hence
$\phi(z[U_0])\leq0$.  Restricted optimality gives
\begin{equation}
 -\phi(z[U_K])=\frac12c^\top z[U_K]
 \leq\frac{\bar\alpha}{2}z[U_K](s).
 \end{equation}
The PageRank mass identity
$d^\top z[U_K]+\mathbf 1^\top r(z[U_K])=1$ and nonnegativity imply
$z[U_K](s)\leq1/d_s\leq1$.  Consequently
$\phi(z[U_0])-\phi(z[U_K])\leq\bar\alpha/2$.
\end{proof}

This is the strongest unconditional warm-start fact obtained from nested
optimality alone.  It controls the sum of squared correction energies, not
the number of coordinates that a solver writes or the edges touched by a
matrix-vector product.  In particular, it cannot charge a fixed positive
full-support pass when the next correction has very small energy.

\section{Why a literal warm start cannot remove the factor}
\label{sec:materialization}

Call an implementation \emph{round-materialized} if, at each active set $U$,
it writes an explicit word for every coordinate of the approximate restricted
solution before evaluating the next gate.  This includes a literal black-box
SDD call returning a dense length-$|U|$ vector.  The next result charges only
those writes and is therefore independent of the internal SDD algorithm.

\begin{theorem}[Endpoint-path materialization barrier]
\label{thm:path-materialization}
For every $J\geq1$, there are positive $\lambda=\kappa$ for which the
successful execution of the source active-set framework on the endpoint-seeded
path $P_{J+1}=(v_0,\ldots,v_J)$ visits every prefix
\begin{equation}
 U_k=\{v_0,\ldots,v_k\},\qquad 0\leq k\leq J,
 \label{eq:path-prefixes}
\end{equation}
and adds exactly one vertex per nonterminal round.  Every round-materialized
implementation consequently performs at least
\begin{equation}
 \sum_{k=0}^{J}|U_k|=\frac{(J+1)(J+2)}2
 =\Omega(|S^\star|\vol(S^\star))
 \label{eq:write-lower}
\end{equation}
word writes, since $S^\star=V(P_{J+1})$ and
$\vol(S^\star)=\Theta(J)$.
\end{theorem}

\begin{proof}
At $\lambda=0$, the exact restricted system on every prefix has a positive
right-hand side at $v_0$.  Its matrix is an irreducible Stieltjes matrix, so
$z[U_k](v_k)>0$.  For $k<J$, the only boundary vertex is $v_{k+1}$ and
\begin{equation}
 r(z[U_k])(v_{k+1})
 =\frac{1-\bar\alpha}{\bar\alpha}z[U_k](v_k)>0.
 \label{eq:positive-path-gate}
\end{equation}
There are finitely many prefixes, and every restricted state and boundary
residue is continuous in $\lambda$.  Hence one may choose
$\lambda=\kappa=\eta>0$ so small that, for every $k<J$, the exact boundary
residue is greater than $2.2\eta d_{v_{k+1}}$ and
$2\eta d_{v_0}<1$, while every prefix state remains positive.

Condition on all SDD calls satisfying the source energy guarantee.  The
source boundary-error lemma then makes the approximate boundary residue
differ from the exact one by at most $0.1\kappa d_{v_{k+1}}$.  It is therefore
strictly larger than
$(\lambda+\kappa)d_{v_{k+1}}=2\eta d_{v_{k+1}}$.  The unique boundary vertex
is admitted.  Induction proves \cref{eq:path-prefixes}; the full path then has
empty boundary and terminates.  The source RPPR support-containment lemma
implies that every admitted vertex lies in $S^\star$, hence $S^\star=V$.
Writing every coordinate on every prefix costs the sum in
\cref{eq:write-lower}.
\end{proof}

The theorem also applies to any iterative warm start that performs one full
$M_{U_kU_k}$ matrix pass in every nonterminal round: on a bounded-degree path
each pass costs $\Theta(\vol(U_k))$, giving the same quadratic sum.  Thus a
successful general improvement must be implicit and incremental; changing
only the initial vector of a conventional SDD solver is insufficient.

\section{The repeated factor disappears exactly on endpoint paths}
\label{sec:path-solver}

We now use the same path family constructively.  The statement is written for
an arbitrary symmetric positive definite tridiagonal Stieltjes matrix, so it
applies both to the source degree form and to the shared symmetric system.

Let the diagonal of $M$ on $(v_0,\ldots,v_N)$ be $a_i>0$ and write the edge
entry as $M_{i-1,i}=M_{i,i-1}=-b_i$ with $b_i>0$.  Let $c_i$ be the fixed
right-hand side.  Define append-only forward records
\begin{align}
 \delta_0&:=a_0,
 &\eta_0&:=c_0,
 \label{eq:path-forward-base}\\
 \delta_i&:=a_i-\frac{b_i^2}{\delta_{i-1}},
 &\eta_i&:=c_i+\frac{b_i}{\delta_{i-1}}\eta_{i-1}.
 \label{eq:path-forward}
\end{align}

\begin{theorem}[Activation-once path solver]
\label{thm:path-linear}
Consider the exact active-set gate \cref{eq:exact-gate} on an endpoint-seeded
path.  At prefix $U_k$, the last coordinate of the exact restricted solution
is
\begin{equation}
 z[U_k](v_k)=\frac{\eta_k}{\delta_k}.
 \label{eq:path-endpoint-value}
\end{equation}
For the source degree-form system, the unique next vertex $v_{k+1}$ violates
the gate exactly when
\begin{equation}
 c_{k+1}+b_{k+1}\frac{\eta_k}{\delta_k}
 >\bar\alpha\kappa d_{v_{k+1}}.
 \label{eq:path-gate-scalar}
\end{equation}
Appending \cref{eq:path-forward} after each successful gate and, at the
terminal prefix $U_J$, performing the single reverse recurrence
\begin{align}
 z(v_J)&=\eta_J/\delta_J,
 \nonumber\\
 z(v_i)&=\frac{\eta_i+b_{i+1}z(v_{i+1})}{\delta_i},
 \qquad i=J-1,\ldots,0,
 \label{eq:path-reverse}
\end{align}
returns the exact terminal restricted solution.  Each admitted vertex and
edge is processed $O(1)$ times.  The total arithmetic, storage, and charged
adjacency work is
\begin{equation}
 O(\operatorname{cvol}(U_J)+d_{v_{J+1}}),
 \label{eq:path-linear-work}
\end{equation}
where the last term is absent when $J=N$.  In particular, the factor
$|U_J|$ multiplying final active volume is eliminated.
\end{theorem}

\begin{proof}
\Cref{eq:path-forward-base,eq:path-forward} are ordinary scalar
$LDL^\top$ elimination on a tridiagonal matrix.  Every pivot $\delta_i$ is
positive because every leading principal submatrix is positive definite.
At the last row of the eliminated prefix there is no later variable, so its
solution is \cref{eq:path-endpoint-value}.

For $v_{k+1}\notin U_k$, extend the prefix state by zero.  From
$c=\bar\alpha(s-\lambda d)$ and \cref{eq:source-residue},
\begin{equation}
 \bar\alpha\bigl(r(z[U_k])(v_{k+1})-\lambda d_{v_{k+1}}\bigr)
 =c_{k+1}-M_{k+1,k}z[U_k](v_k).
\end{equation}
Substituting $M_{k+1,k}=-b_{k+1}$ and
\cref{eq:path-endpoint-value} proves \cref{eq:path-gate-scalar}.  Once the
gate fails, standard triangular back substitution is exactly
\cref{eq:path-reverse}.  The forward pass stores two scalars per vertex, the
reverse pass visits each record once, and path adjacency is exposed once,
which proves \cref{eq:path-linear-work}.
\end{proof}

\begin{corollary}[Source guarantees with linear path work]
\label{cor:path-source-guarantees}
On endpoint paths, replace every approximate SDD solve in the source active-set
framework by \cref{thm:path-linear} and evaluate the exact gate.
\begin{enumerate}[label=(\roman*)]
\item With $\lambda=\kappa=\varepsilon/2$, the output is a deterministic ACL
$\varepsilon$-approximate PageRank vector of support volume at most
$2/\varepsilon$.
\item With $\lambda=\rho$ and
$0<\kappa\leq\min\{\rho,\xi/\bar\alpha\}$, every active prefix is contained
in $S^\star$, the returned degree-form vector has RPPR objective gap at most
$\bar\alpha\kappa\leq\xi$, and its PageRank output is an ACL
$(\rho+\kappa)$-approximation.
\end{enumerate}
Both conclusions use the linear work in \cref{eq:path-linear-work} and have
no polynomial dependence on $1/\bar\alpha$.
\end{corollary}

\begin{proof}
At termination, exact restricted optimality gives residue $\lambda d$ on the
active prefix and the failed gate gives residue at most
$(\lambda+\kappa)d$ outside it.  Residue nonnegativity follows from the
Stieltjes comparison and the boundary formula.  The PageRank mass identity
then gives active volume at most $1/\lambda$, proving (i).

For (ii), the exact support comparison used by Wei and Yang shows inductively
that a boundary vertex with residue greater than $\rho d_v$ lies in
$S^\star$; the stronger gate therefore keeps every prefix inside
$S^\star$.  For the RPPR objective $\psi$, a valid subgradient at the terminal
point is $\bar\alpha(\rho d-r)$.  It is zero on the active prefix and bounded
below by $-\bar\alpha\kappa d$ outside.  Convexity, zero exterior coordinates,
and $d^\top z^\star\leq1$ give
\begin{equation}
 \psi(z)-\psi(z^\star)
 \leq -\langle\bar\alpha(\rho d-r),z^\star\rangle
 \leq\bar\alpha\kappa.
\end{equation}
The remaining claims follow directly from the terminal residue bounds and
\cref{thm:path-linear}.
\end{proof}

This theorem is stronger than an accelerated
$\widetilde O(\vol(S^\star)/\sqrt\alpha)$ path bound: it is output-linear.
It is compatible with the activation-token and exact-Schur path results in
the project's adaptive-revisit and propagate--settle notes.  It also explains
why the materialization lower bound is not an information lower bound.

\section{The exact arbitrary-graph target}
\label{sec:interface}

The path proof suggests the correct abstraction.  It must include boundary
queries; an incremental linear-system solver alone does not implement the
outer gate.

\begin{definition}[Implicit nested solve-and-boundary interface]
\label{def:interface}
For fixed $(M,c,\lambda,\kappa)$, an implicit nested interface stores a shared
state for the current $U$ and supports:
\begin{enumerate}[label=(\roman*)]
\item \textsc{Violations}: return exactly, or with the source-safe error
margin, all $v\in\partial U$ satisfying the activation gate;
\item \textsc{Expand}$(T)$: update the representation from $z[U]$ to
$z[U\cup T]$ without materializing every old coordinate;
\item \textsc{Finalize}: materialize the terminal vector once and return the
required residue or objective certificate.
\end{enumerate}
Its total work includes all adjacency exposure, numerical updates, violation
queries, and output writes.
\end{definition}

\begin{theorem}[Conditional removal of the repeated factor]
\label{thm:conditional-interface}
Suppose \cref{def:interface} realizes the source error margins and has total
work $\mathcal I(U_K)$ over a complete growing-active-set trace.  Then the
source ACL and RPPR correctness proofs remain valid and their total work is
\begin{equation}
 O\bigl(\mathcal I(U_K)+\operatorname{cvol}(U_K)\bigr),
 \label{eq:interface-work}
\end{equation}
with no additional factor equal to the number of active-set rounds.  In
particular, a graph-uniform
$\mathcal I(U_K)=\widetilde O(\operatorname{cvol}(U_K))$ interface would give
$\widetilde O(1/\rho)$ RPPR work, while
$\mathcal I(U_K)=\widetilde O(\operatorname{cvol}(U_K)/\sqrt\alpha)$ would
give the project's conjectured product scale.
\end{theorem}

\begin{proof}
The source induction uses only the restricted-solution error bounds, the
boundary-residue error bound, and the returned violation set.  It is
independent of how those quantities are represented.  Support containment
then gives $U_K\subseteq S^\star$ and
$\operatorname{cvol}(U_K)=O(\operatorname{cvol}(S^\star))$.  Charging the
interface once over the whole trace and the terminal output once proves
\cref{eq:interface-work}.  The source mass bound
$\vol(S^\star)\leq1/\rho$ gives the two stated consequences.
\end{proof}

The theorem is deliberately conditional: writing an oracle definition does
not construct the oracle.  It does, however, prevent a hidden repeated scan
from being mislabeled as incremental reuse.

\section{Why the general construction remains open}
\label{sec:open}

The block identity pinpoints three distinct difficulties.
\begin{enumerate}[label=(\arabic*)]
\item \emph{Dense transport.}  Even when $T$ is one boundary vertex,
$\Delta_U=-M_{UU}^{-1}M_{UT}\Delta_T$ can be nonzero throughout the old
active component.  Explicitly applying this correction recreates the
repeated-prefix cost.
\item \emph{Boundary queries.}  A changed old coordinate changes every
outside residue incident to it.  A dynamic SDD solve that answers energy or
single-coordinate queries does not automatically maintain the complete set
of threshold violations required by the active-set gate.
\item \emph{Energy is not work.}  The telescope
\cref{eq:telescope-bound} permits many small, spatially broad corrections.
It bounds a numerical quantity but does not pay a fixed full-support pass,
factor rebuild, or vector write.
\end{enumerate}

The next proof target is therefore not ``warm-start the nearly-linear SDD
solver.''  It is the following concrete data-structural statement.

\begin{conjecture}[Aggregate implicit continuation]
\label{conj:aggregate}
For the nested support-safe PageRank/RPPR trace, there is a local data
structure implementing \cref{def:interface} with high probability in total
work
\begin{equation}
 \widetilde O\!\left(
 \frac{\operatorname{cvol}(S^\star)}{\sqrt\alpha}
 \operatorname{polylog}\frac{1}{\rho\xi\delta}
 \right).
 \label{eq:aggregate-target}
\end{equation}
The stronger $\widetilde O(\operatorname{cvol}(S^\star))$ bound suggested by
the 2026 source is also open on arbitrary graphs.
\end{conjecture}

Paths prove \cref{conj:aggregate} with the stronger output-linear bound.
Equitable shells, bounded response blocks, and several tree-like classes in
the companion notes also admit activation-once compressed transport.  The
unresolved case is a large nonequitable cyclic core in which Schur transport
is broad and the boundary has no known low-dimensional kinetic description.

\section{OP3 exploration: geometric value events}
\label{sec:op3-geometric-events}

\paragraph{Status and scope (6 September 2026).}
The statements marked \statusproved{} in this section are new proof drafts,
awaiting independent review. They do not resolve OP3. Their purpose is to
replace exact reproduction of every source gate by a weaker, explicitly
quantified sufficient interface. The investigation and source comparison are
recorded in \texttt{OP3\_DIRECTIONS\_20260906.md}. The exact audit uses dense
reference solves to supply the missing interface and does not measure an
output-linear solver.

Throughout this section, $0<\alpha<1$ and $0<\varepsilon<1$, with
$\varepsilon=\epsappr$ in the source-native ACL namespace. Retain the
mapping of \cref{sec:mapping}, and put
\[
 \gamma:=1-\bar\alpha,\qquad
 \bm M=\bm D-\gamma \bm A,\qquad
 \lambda:=\varepsilon/2.
\]
For a reached face $U$ containing the seed, define the auxiliary response
and its degree-coordinate solution by
\begin{equation}
 \bm u^U_U=\bm M_{UU}^{-1}(\bm e_v-\lambda\bm d)_U,
 \qquad \bm u^U_{\bar U}=\bm0,
 \qquad \bm z^U=\bar\alpha\bm u^U.
 \label{eq:op3-scaled-response}
\end{equation}
Thus the canonical RPPR point is $\bm D^{1/2}\bm z^U$, and the PageRank
output is $\bm D\bm z^U$. No optimum symbol is renamed. The source residual is
$\bm r^U=\bm e_v-\bm M\bm u^U$.

\begin{lemma}[Response size and monotonicity; proof draft]
\label{lem:op3-response-size}
Start with $U=\{v\}$ when $\lambda d_v<1$, and enlarge $U$ by any
nonempty subset of exterior vertices satisfying $r_i^U>\lambda d_i$.
Every restricted solution is positive on its face, increases on old
coordinates, and is bounded above by the scaled RPPR optimum. Moreover,
\begin{equation}
 \bm r^U\geq\bm0,\qquad r_i^U=\lambda d_i\ (i\in U),\qquad
 \lambda\vol(U)+\bar\alpha\sum_{i\in U}d_i u_i^U
       +\sum_{i\notin U}r_i^U=1.
 \label{eq:op3-mass-identity}
\end{equation}
In particular, $\vol(U)\leq2/\varepsilon$ and
$0\leq u_i^U\leq1/(\bar\alpha d_i)$.
\end{lemma}
\begin{proof}
The singleton has positive load. On an admitted block the padded old point
has strictly positive residual relative to the new restricted load. Inverse
positivity of the principal Stieltjes matrix proves monotone, positive
correction, as in \cref{thm:block-schur}. Comparing the restricted equation
with the full obstacle supersolution proves the upper bound by the RPPR
optimum. An exterior violation cannot lie outside that optimum's support:
at a zero optimum coordinate, its negative off-diagonal entries and the
comparison on $U$ would otherwise contradict obstacle dual feasibility.
Residuals outside the face are $\gamma \bm A_{\bar U,U}\bm u_U^U\geq\bm0$;
inside they equal $\lambda\bm d_U$. Finally
$\one^\top \bm M=\bar\alpha\bm d^\top$. Summing the residual identity proves
\cref{eq:op3-mass-identity} and its consequences.
\end{proof}

\begin{definition}[Geometric value publications]
\label{def:op3-publications}
Set
\begin{equation}
 h:=\frac{\varepsilon}{16\gamma},\qquad
 q:=\frac54,\qquad a_0:=\frac{11}{10},\qquad
 \vartheta:=\frac{11\varepsilon}{20}.
 \label{eq:op3-event-constants}
\end{equation}
For each active coordinate maintain a published value $\ell_i\geq0$,
initially zero. Every publication is a lower bound on the current $u_i^U$.
The first positive publication is at least $h/2$; every later change
multiplies the preceding positive value by at least $a_0$. Before declaring
a face ready for a boundary decision, the producer certifies
\begin{equation}
 \ell_i\leq u_i^U\leq q\ell_i+h\qquad(i\in U).
 \label{eq:op3-publication-band}
\end{equation}
The producer must find changes and certify readiness without assuming a
free scan. Its entire update and search work is denoted by $\mathcal P$.
The actual coordinatewise bound is required; a spectral approximation alone
does not certify it.
\end{definition}

For an exterior vertex maintain the lower boundary signal
\begin{equation}
 L_j:=\gamma\sum_{i\in U\cap\mathcal N(j)}\ell_i.
 \label{eq:op3-lower-signal}
\end{equation}
When $\ell_i$ increases by $\Delta\ell_i$, scan its cached adjacency list
and add $\gamma\Delta\ell_i$ to each surviving exterior neighbor's key.
The incidence reads are charged even when the list was exposed earlier.
Admit any nonempty subset with $L_j>\vartheta d_j$, and repeat. A terminal
decision requires both a ready producer and an empty violation dictionary.

\begin{theorem}[An output-linear recipient of value events; proof draft]
\label{thm:op3-geometric-recipient}
Suppose a producer implements \cref{def:op3-publications} along this
execution. Let $V=\vol(U_{\rm fin})$. The execution is support-safe and
terminates after at most $|U_{\rm fin}|$ nonempty faces. Its exact terminal
state satisfies
\begin{equation}
 r_i^{U_{\rm fin}}=\tfrac12\varepsilon d_i\quad(i\in U_{\rm fin}),
 \qquad
 0\leq r_j^{U_{\rm fin}}\leq\tfrac34\varepsilon d_j
       \quad(j\notin U_{\rm fin}).
 \label{eq:op3-quiet-certificate}
\end{equation}
All graph exposure, repeated incidence delivery, dictionary maintenance,
and publication records outside the producer cost
\begin{equation}
 \widetilde O\!\left(
  V\left[1+\log_+\frac{32\gamma}{\bar\alpha\varepsilon}\right]+1
 \right).
 \label{eq:op3-recipient-work}
\end{equation}
Here $\log_+ t=\max\{0,\log t\}$. The full work is the displayed cost
plus $\mathcal P$ and terminal numerical completion. In particular,
\emph{this theorem does not bound $\mathcal P$}.
\end{theorem}
\begin{proof}
For any exterior $j$, positivity and \cref{eq:op3-publication-band} imply
\begin{equation}
 L_j\leq r_j^U
 \leq qL_j+\gamma h\,|\mathcal N(j)\cap U|
 \leq qL_j+\tfrac1{16}\varepsilon d_j.
 \label{eq:op3-signal-sandwich}
\end{equation}
Since $\vartheta>\lambda$, each reported admission satisfies the exact
support-safe condition of \cref{lem:op3-response-size}. If no key exceeds
its threshold, then
$q\vartheta+\varepsilon/16=3\varepsilon/4$, proving
\cref{eq:op3-quiet-certificate}.

Each coordinate starts positive publication at $h/2$ or more, grows by a
factor at least $a_0$, and never exceeds $1/(\bar\alpha d_i)$. It therefore
has at most
$1+\lceil\log_{a_0,+}(32\gamma/(\bar\alpha\varepsilon))\rceil$
positive publications. Multiplying by $d_i$ and summing gives the repeated
incidence bound. A dictionary with logarithmic word cost stores the
exposed neighbor labels, degrees, scores, and active flags. The number of
distinct labels and incidences is $O(V+1)$. Candidate keys are monotone
until their vertex is admitted; they can be queued once and removed at
admission. New-face exposure costs $V$, and there are at most $|U_{\rm fin}|$
nonempty faces. All event-location and ready-certification costs remain in
$\mathcal P$, including empty searches.
\end{proof}

\paragraph{Finite bands do not require exact bucket decisions.}
An exact geometric floor with ratio $5/4$ and minimum value $h$ is one
valid publication policy. It is not necessary to decide such floors near
equality. For example, suppose a queried coordinate comes with an interval
$[L,H]$ containing its exact value and of width at most $h/4$. Retain the
old publication if $H\leq q\ell+h$. Otherwise $L>q\ell+3h/4$; publish
the largest value $(h/2)a_0^k\leq L$. If $\ell>0$, this is a strict
increase by at least $a_0$, because $q>a_0^2$; if $\ell=0$, it is at least
$h/2$. The new value is at most the truth, and
$q\ell_{\rm new}+h\geq H$. Thus there is an overlap band and no
trajectory minimum-complementarity-margin parameter. Producing these
intervals and choosing which coordinates to query are still charged to
$\mathcal P$.

\begin{lemma}[One final certified SDD solve suffices; proof draft]
\label{lem:op3-terminal-repair}
Let $U$ satisfy \cref{eq:op3-quiet-certificate}. Compute a numerical
solution $\widetilde{\bm z}$ on $U$ such that
\begin{equation}
 \|\widetilde{\bm z}-\bm z^U\|_\infty\leq\delta,
 \qquad \delta:=\bar\alpha\varepsilon/16.
 \label{eq:op3-terminal-accuracy}
\end{equation}
Return $\widehat z_i=\max\{0,\widetilde z_i-\delta\}$ on $U$ and zero
outside. Then $\bm D\widehat{\bm z}$ is an ACL $\varepsilon$-approximation
of support volume at most $2/\varepsilon$. A supplied-face SDD solver with
certified retries performs this completion in expected
$\widetilde O(V+1)$ word work, hiding only logarithmic parameter factors.
\end{lemma}
\begin{proof}
Write $\bm e=\bm z^U-\widehat{\bm z}$; then $0\leq e_i\leq2\delta$.
Outside $U$ the residual can only decrease and remains nonnegative.
On $U$ its upper bound is
\[
 r_i(\widehat{\bm z})
 =\lambda d_i+\bar\alpha^{-1}(\bm M\bm e)_i
 \leq(\lambda+2\delta/\bar\alpha)d_i
 =\tfrac58\varepsilon d_i.
\]
If $\widehat z_i>0$, the same formula gives a lower bound
$(\lambda-2\gamma\delta/\bar\alpha)d_i\geq0$.
If $\widehat z_i=0$, its residual is directly
$(\bm e_v)_i+\gamma\bar\alpha^{-1}(\bm A\widehat{\bm z})_i\geq0$.
This proves the ACL witness and the support bound.

The matrix $\bm M_{UU}$ is SDD with $O(V)$ nonzeros. The exact test
\[
 \|\bm D_U^{-1}(\bm M_{UU}\widetilde{\bm z}_U
       -\bar\alpha(\bm e_v-\lambda\bm d)_U)\|_\infty
 \leq\bar\alpha\delta
\]
implies \cref{eq:op3-terminal-accuracy}, since
$\bm M_{UU}\one\geq\bar\alpha\bm d_U$ and the inverse is nonnegative.
It costs one charged pass over the face. Also
$\|\bm z^U\|_{\bm M_{UU}}^2\leq\bar\alpha$, so relative energy error at
most $\bar\alpha^2\varepsilon/16$ suffices for this test by rowwise
energy Cauchy--Schwarz. The tolerance is polynomial in the named parameters.
The supplied SDD result of \citet{koutis2011nearly}, applied only to the
exposed face, followed by independent certified retries gives the claimed
expected cost. Capped retries alternatively give explicit high-probability
completion or a reported failure. There is no ambient scan.
\end{proof}

The cases $\lambda d_v\geq1$, $\alpha=1$, and $\varepsilon\geq1$ are
handled directly by the zero or one-coordinate solution. Combining
\cref{thm:op3-geometric-recipient,lem:op3-terminal-repair} leaves a precise
\statusopen{} target: construct the producer with
$\mathcal P=\widetilde O(V+1)$ along this finite-band, one-source trace.
This permits a different sequence of faces from the literal Wei--Yang
algorithm. It neither maintains arbitrary changing right-hand sides nor
answers arbitrary inverse queries. The older additive-quantization and
exact-full-boundary obstructions must not be silently transferred to this
weaker interface.

\subsection{A separate teleportation-continuation probe}
\label{subsec:op3-teleportation-probe}

Put $t=\bar\alpha/(1-\bar\alpha)$ and $\bm L=\bm D-\bm A$. For fixed $\lambda$,
define, only in this subsection,
\[
 \bm w_t:=\argmin_{\bm w\geq0}
 \left\{\tfrac12\bm w^\top(\bm L+tD)\bm w
                  -(\bm e_v-\lambda\bm d)^\top\bm w\right\}.
\]
The degree-coordinate RPPR optimum is $t\bm w_t$. If $0<t'<t$, then
$\bm w_{t'}\geq\bm w_t$: the smaller-$t$ solution is a nonnegative
supersolution for the larger-$t$ obstacle, and the obstacle solution is the
least such supersolution. Thus supports are nested in this rescaling.
For $t'=t/2$, the formed matrices obey
\begin{equation}
 \tfrac12(\bm L+tD)\preceq \bm L+(t/2)\bm D\preceq \bm L+tD.
 \label{eq:op3-teleportation-spectral-sandwich}
\end{equation}
Both are useful \statusproved{} proof-draft facts, but they do not give a
constant-cost local continuation stage.

On the three-vertex path seeded at its first endpoint, take $\lambda=1/5$.
The two exact solutions are
\[
 \bm w_{3/2}=(8/25,0,0)^\top,\qquad
 \bm w_{3/4}=(96/205,4/205,0)^\top.
\]
Direct substitution verifies all obstacle inequalities. A new coordinate
becomes positive even though the matrices satisfy
\cref{eq:op3-teleportation-spectral-sandwich}. Therefore no uniform
coordinatewise multiplicative bracket by the old solution follows. A
preconditioned obstacle step must still discover new support and apply its
non-Euclidean constrained map; a good preconditioner for an unconstrained
linear solve does not perform that step.

\section{Supplied trees: persistent affine response curves}
\label{sec:op3-persistent-tree}

\paragraph{Status and model.}
\statusproved{} below means a new proof draft, awaiting independent review.
The input is an explicitly supplied tree $T$ on $n\geq2$ vertices, rooted at
the seed $v$, with diagonal data $d_i\geq\deg_T(i)>0$. Put
$\bm M=\bm D-\gamma\bm A_T$, $0<\gamma<1$, $\bar\alpha=1-\gamma$, and
$\bm b=\bm e_v-\lambda\bm d$ for $\lambda>0$. We solve exactly
\begin{equation}
 \min_{\bm u\geq\bm0}
 \frac12\bm u^\top\bm M\bm u-\bm b^\top\bm u.
 \label{eq:op3-tree-obstacle}
\end{equation}
All input reads, arithmetic, comparisons, stored words, copied nodes and
output writes are charged in the exact-real word model. Coefficient bit
length and floating-point stability are outside this claim. The whole tree
is supplied: the theorem does not charge its discovery to optimal support.

\paragraph{Source and adaptation.}
The affine composition tree of \citet[Section 3, Theorem 3.1]{agarwal2010act}
represents a piecewise-linear curve that is monotone in both coordinates.
Its whole-curve affine operation costs constant time, and evaluations,
insertions and interval transforms cost logarithmic time. The response
embedding and persistence argument below are our adaptation. We spell out
the needed operations rather than invoking the source's more general
tree-set representation: transforming separate summands by an inverse-type
map would not transform their sum correctly.

\begin{lemma}[A strictly increasing response embedding]
\label{lem:op3-tree-embedding}
For a nonroot vertex $i$, condition its parent's potential to be $t\geq0$,
minimize the descendant objective, and denote the resulting root coordinate
by $u_i(t)$. The curve $\nu_i(t):=t+\gamma u_i(t)$ is convex, continuous,
strictly increasing and piecewise affine. Extend its initial branch to
negative inputs by $\nu_i(t)=t$. It has at most $|T_i|$ knots, where $T_i$
is the descendant subtree. It can be constructed from its children's
curves by addition, two whole-curve affine maps and one new knot.
\end{lemma}
\begin{proof}
For $k_i$ children define
\[
 S_i(x)=\sum_{j\text{ child of }i}\nu_j(x),\qquad
 H_i(x)=(d_i+k_i)x-b_i-S_i(x).
\]
For $x\geq0$, $H_i$ is the derivative of the minimum descendant energy
conditioned on $u_i=x$. On each active-set interval its slope is the Schur
diagonal of a principal submatrix of $\bm M$. Since
$\bm M\succeq\bar\alpha\bm D$, minimizing the quadratic form with root
coordinate fixed to one gives $H_i'\geq\bar\alpha d_i>0$.
Inductively the child curves are convex; hence $H_i$ is concave and strictly
increasing. Its inverse is convex. Since $b_i=-\lambda d_i<0$, the response
is zero until $t_i^0=\lambda d_i/\gamma$ and thereafter satisfies
$H_i(u_i(t))=\gamma t$. This proves convexity and strict increase of
$\nu_i$ across the join with the line $t$.

The graph of $H_i$ is obtained from that of $S_i$ by
\[
 (x,y)\longmapsto\bigl(x,(d_i+k_i)x-b_i-y\bigr).
\]
The positive branch of $\nu_i$ follows by the second map
\[
 (x,y)\longmapsto\bigl(y/\gamma,\gamma x+y/\gamma\bigr).
\]
Each map has nonzero determinant and preserves the order of points along
the applicable curve. Every inherited knot has $x>0$ before the second
map, and therefore new abscissa greater than $H_i(0)/\gamma=t_i^0$.
Insert $(t_i^0,t_i^0)$ and replace the left ray by the line $t$. There are
at most the sum of the children's knot counts plus one knots. Leaves
start the induction with one knot.
\end{proof}

At the seed construct $H_v$ without the inverse-type map, and set
$u_v=\max\{0,H_v^{-1}(0)\}$. Recover a child coordinate from
\begin{equation}
 u_i=\frac{\nu_i(u_{\operatorname{parent}(i)})
                 -u_{\operatorname{parent}(i)}}{\gamma}.
 \label{eq:op3-tree-reconstruct}
\end{equation}
Thus the children's original curve versions must remain available after
the upward pass. Destructive reuse alone would not implement this step.

\begin{lemma}[Persistent curve primitives]
\label{lem:op3-persistent-primitives}
The operations used in \cref{lem:op3-tree-embedding}, together with adding
one convex child curve containing $k$ knots to another, admit a persistent
implementation. Whole-curve affine maps take $O(1)$ work and new nodes;
evaluation, inverse evaluation, knot insertion and a positive hinge
addition take $O(\log(n+1))$ work and at most that many new nodes. Adding
the child takes $O((k+1)\log(n+1))$ work and new nodes. Old curve versions
remain unchanged.
\end{lemma}
\begin{proof}
Store knots in an immutable AVL tree, in their common increasing order in
both coordinates. Each node contains its point, two child references,
height, subtree size, and a constant-size pending affine map. Store the two
unbounded-ray slopes with the curve. A whole-curve map copies the root,
transforms its point and ray slopes, and composes its pending map.
Monotonicity makes the old search-tree order valid after the transform.

An evaluation or inverse evaluation follows one search path, carrying
composed affine maps, then interpolates between the bracketing knots or
on an end ray. Insertion copies one search path, pushing each pending map
to copies of at most two child roots. AVL rotations involve a constant
number of such pushes and copies per visited level. Thus both work and
allocations are logarithmic; old roots and descendants are immutable.

For a hinge $c(x-t)_+$, $c\geq0$, first insert the point at $t$ if needed.
The suffix shear $(x,y)\mapsto(x,y+c(x-t))$ visits one path to the suffix
boundary and tags whole subtrees branching to its right. This copies only
$O(\log(n+1))$ nodes. The shear is continuous at $t$ and preserves strict
increase. Every child has a representation
$\nu(x)=x+\sum_a c_a(x-t_a)_+$ with $c_a\geq0$.
Traverse only the smaller child's $k$ knots to obtain successive slope
increments, add its initial $x$ by a whole-curve shear, then add the hinges.
This proves the asserted merge cost without traversing the larger curve.
\end{proof}

\begin{theorem}[Exact supplied-tree solve; proof draft]
\label{thm:op3-persistent-tree}
Problem \eqref{eq:op3-tree-obstacle} on the supplied tree can be solved
deterministically in $O(n\log^2(n+1))$ exact-real word operations and
allocated words, including one terminal output of all coordinates. There
is no polynomial dependence on $1/\alpha$ or $1/\lambda$ in this word bound.
\end{theorem}
\begin{proof}
Read the supplied tree and degrees, root it, and compute subtree sizes in
$O(n)$ work. At each vertex retain the largest-subtree child's curve as
the initial sum and merge all smaller children using
\cref{lem:op3-persistent-primitives}. The sum of the sizes of smaller child
subtrees, over all vertices, is $O(n\log(n+1))$: a vertex can belong to a
smaller child at most logarithmically many times, because the containing
subtree at least doubles on each such ascent. Each child curve contains
at most its subtree size in knots. The logarithmic cost per visited knot,
plus the $O(n)$ affine maps and $O(n)$ knot insertions, gives the stated
upward work and allocation bound. Coincident knots only reduce this count.
Retain every child's immutable root. The root inverse query and the
downward evaluations in \cref{eq:op3-tree-reconstruct} cost
$O(n\log(n+1))$, including all accesses and output writes. Conditional
minimization at every parent proves that the resulting vector is the
unique optimum of \cref{eq:op3-tree-obstacle}.
\end{proof}

\paragraph{Measured audit and remaining gap.}
The persistent rational implementation passed
1,152 supplied-tree comparisons, 43,614 value/inverse round trips and
19,584 checks of retained child versions against separate exact obstacle
solves. Twenty-four path, star, binary-tree and comb diagnostics count
actual allocations and visits. These are falsification checks, not an
independent proof review. A 512-vertex path still constructs all 512
vertices despite an optimum supported on 66. A demand-driven construction
with near-linear work in discovered support volume remains \statusopen{}.
The companion \texttt{delayed\_reflection\_ladder} note already has a lazy
local tree iterator; its additional ancestor-walk cost is precisely what
an OP3-scale local extension would have to remove.

\section{A local exact solver on cycles with unequal leaf attachments}
\label{sec:op3-local-cycle}

\paragraph{Status and scope.}
This is a \statusproved{} proof draft with an exact implementation and
falsification audit, awaiting independent review. The promised graph is a
simple cycle of unknown length $m\geq3$, with any number $k_i\geq0$ of
degree-one leaves attached to cycle vertex $i$. The $k_i$ may all differ.
The seed is arbitrary. Input consists only of the seed and degree/adjacency
queries; there is no supplied cycle ordering, length, decomposition,
equitable partition or future attachment data. Class verification is not
part of the promise algorithm. General unicyclic graphs with nontrivial
attached trees, and arbitrary cyclic graphs, remain outside this theorem.

\paragraph{Context.}
Root-conditioned scalar homotopy and a symmetry-based pure-cycle solver
already appear in the companion \texttt{delayed\_reflection\_ladder} note.
Here we give the two-tip implementation for unequal leaf attachments,
including its closing-block exception. The proof below is self-contained.

Use the exact-real word model and the degree obstacle problem
\eqref{eq:op3-tree-obstacle}, now with the actual graph degrees. Here
$\lambda>0$ is free; the ACL corollary will set $\lambda=\epsappr$.
All degree replies, adjacency entries, retained records, event searches,
arithmetic and output writes are charged. We first take a cycle seed and
extend to a leaf seed at the end.

\subsection{The continuation parameter and the two exposed ends}

Fix the seed potential $u_v=t\geq0$ and minimize over all other
coordinates. Denote the derivative of this reduced objective by $g(t)$.
It is continuous, strictly increasing and piecewise affine. On each
active face its slope is a positive Schur diagonal. The conditional
coordinates are nondecreasing in $t$, by inverse positivity and the
obstacle comparison principle. The optimum is zero if
$1-\lambda d_v\leq0$; otherwise it occurs at the unique positive root of
$g$. Start at $t=0$, with only the seed exposed.

The two cycle directions from the seed are called arms. Before they meet,
each has an admitted consecutive prefix and one unadmitted candidate. An
unadmitted cycle vertex becomes positive when the sum of its incident
active potentials times $\gamma$ reaches $\lambda d_i$. No other cycle
vertex can activate before one of these candidates. The adjacency row of
a candidate is read only after this event is known to precede the root of
$g$. Its two nonleaf neighbors identify the continuation ports using
charged degree queries. Its other neighbors form a leaf group.

A negative-load leaf attached to a center of potential $x$ has value
\begin{equation}
 u_{\rm leaf}=(\gamma x-\lambda)_+.
 \label{eq:op3-cycle-leaf}
\end{equation}
All $k$ leaves at one center activate at $x=\lambda/\gamma$. Once active,
eliminating them changes its diagonal and load by
\begin{equation}
 d\longleftarrow d-k\gamma^2,\qquad
 b\longleftarrow b-k\gamma\lambda.
 \label{eq:op3-cycle-leaf-fold}
\end{equation}
Reading and storing their labels is charged to the center's degree even
when some leaves are zero in the final answer.

\begin{lemma}[Settled interior prefixes]
\label{lem:op3-cycle-settled}
Before cycle closure, a cycle vertex leaves the active end of an arm only
after its leaf group has activated, or when that group is empty. Therefore
every committed interior record has a fixed quadratic response for the
rest of the continuation.
\end{lemma}
\begin{proof}
A next candidate reached from just this arm requires
$\gamma u_{\rm tip}=\lambda d_{\rm next}\geq2\lambda$.
The tip's leaves activate already at $u_{\rm tip}=\lambda/\gamma$.
Monotonicity prevents later deactivation. The only exception to the
one-sided candidate rule is the last inactive cycle vertex, which is
handled separately below.
\end{proof}

\subsection{Constant-size end records and exact event location}

Eliminate a settled arm prefix outward from the seed. Its current tip has
diagonal $\delta_i>0$ and affine load $\eta_i(t)=a_i t+b_i$, so its
potential is $(a_i t+b_i)/\delta_i$. Initially
$(\delta_i,a_i,b_i)=(d_i,\gamma,-\lambda d_i)$. When committing this tip
and admitting the next cycle vertex $j$, its new record is
\begin{equation}
 \delta_j=d_j-\frac{\gamma^2}{\delta_i},\qquad
 a_j=\frac{\gamma a_i}{\delta_i},\qquad
 b_j=-\lambda d_j+\frac{\gamma b_i}{\delta_i}.
 \label{eq:op3-cycle-append}
\end{equation}
Apply \cref{eq:op3-cycle-leaf-fold} to a tip when its leaf group activates.
All Schur pivots stay positive because they eliminate principal submatrices
of the positive definite Stieltjes matrix.

Maintain two accumulated coefficients over committed records,
\[
 P=\sum_i a_i^2/\delta_i,\qquad
 R=\sum_i a_i b_i/\delta_i.
\]
Let $(\delta_v,b_v)$ denote the seed diagonal/load with any activated seed
leaves folded. While the arms are distinct, the current face derivative is
\begin{equation}
 g(t)=\left(\delta_v-P-\sum_{\rm tips}\frac{a_i^2}{\delta_i}\right)t
       -b_v-R-\sum_{\rm tips}\frac{a_i b_i}{\delta_i}.
 \label{eq:op3-cycle-root-derivative}
\end{equation}
This follows by differentiating the eliminated contributions
$-\eta_i(t)^2/(2\delta_i)$. Each committed record is inserted into $P,R$
once; old records are not refreshed when $t$ changes.

At a state, compute the zero $t_\star$ of the current affine expression
for $g$ and the earliest of at most five prospective events: the seed leaf
group, up to two tip leaf groups, and up to two next cycle candidates.
Every event time follows from a single affine equality using the stored
tip responses. If $t_\star$ is no later than the first event, stop.
Otherwise process that event and repeat. Equal-time events are processed
in any fixed order; a zero at the same time causes stopping before row
exposure. Continuity of $g$ makes this tie rule correct, with zero
coordinates omitted from the output.

\subsection{Closing the cycle without folding inactive leaves}

If the candidate labels of the two arms coincide at $j$, the candidate
receives two-sided load. Write each tip response as $p_a t+q_a$; an empty
arm contributes the seed response $t$. Its activation time is now
\begin{equation}
 t_j=\frac{\lambda d_j-\gamma(q_1+q_2)}
                 {\gamma(p_1+p_2)}.
 \label{eq:op3-cycle-close-time}
\end{equation}
Using the two separate one-sided events here would miss the first
activation. Furthermore, $t_j$ can precede activation of a tip's leaves.
Such a tip cannot yet become a permanently folded interior record.

Keep both current tips and the newly admitted $j$ in a Schur block of
dimension at most three. The rest of each arm already satisfies
\cref{lem:op3-cycle-settled}. The block matrix $\bm S$ has the stored tip
diagonals, diagonal $d_j$, and off-diagonal $-\gamma$ between $j$ and
each nonempty arm tip. Its load is $\bm a t+\bm b$: old tip coefficients
are retained and $a_j=\gamma$ times the number of empty arms. The block
response is
\[
 \bm x(t)=\bm S^{-1}\bm a\,t+\bm S^{-1}\bm b.
\]
Replace the two tip sums in \cref{eq:op3-cycle-root-derivative} by
$\bm a^\top\bm S^{-1}\bm a$ and
$\bm a^\top\bm S^{-1}\bm b$. Any remaining leaf events are obtained from
$x_i(t)=\lambda/\gamma$. There are at most three such groups. Each is
folded by a diagonal/load change inside this same block, and both solves
have dimension at most three. All seed leaves were already settled before
closure by the one-sided first admission. Thus no larger dynamic solve or
search over old vertices is hidden in this step.

At the final $t_\star$, evaluate the tip or closing-block responses. Recover
each arm's committed prefix in reverse order using its stored records:
\begin{equation}
 u_i=\frac{a_i t_\star+b_i+\gamma u_{\rm next}}{\delta_i}.
 \label{eq:op3-cycle-reverse}
\end{equation}
Read each exposed center's leaf list once and use
\cref{eq:op3-cycle-leaf}. Materialize positive coordinates only.

\subsection{Correctness, locality and the approximation guarantee}

\begin{theorem}[Local exact obstacle solve on the promised cyclic family]
\label{thm:op3-local-cycle}
On a simple cycle with arbitrary pendant-leaf counts and an arbitrary
seed, the above algorithm computes the unique obstacle optimum with
$\widetilde O(1+\vol(S))$ charged exact-real word work and storage, where
$S$ is its positive support. Its structural arithmetic, record and query
counts are $O(1+\vol(S))$; the tilde allows deterministic label-dictionary
overhead. It exposes adjacency rows only at vertices in $S$, and
$\lambda\vol(S)\leq1$. It requires no unknown-support preprocessing and
no polynomial factor in $1/\alpha$.
\end{theorem}
\begin{proof}
For a cycle seed, the conditional obstacle has monotone support as $t$
increases. A face starts valid at its activation time, with newly admitted
coordinates zero. Its equations remain valid until the first nonactive
cycle or leaf gate reaches zero. The two-arm description enumerates all
possible cycle gates; disconnected inactive vertices have strictly positive
dual slack. The leaf formula enumerates every remaining type of gate.
\Cref{lem:op3-cycle-settled} justifies committing interiors, while the
small closing block retains all possible unsettled exceptions. The Schur
formulas give the exact reduced derivative and event times on this face.
Its positive slope and continuity certify either the exact optimum before
the next event or the next valid face. At any processed event, $g<0$, so
the eventual root lies strictly later. Monotonicity and strict positive
boundary influence imply that every exposed center is positive finally.
The reverse equations and leaf formulas recover the conditional minimizer
at the root of $g$; all obstacle KKT conditions therefore hold.

There is one admission per exposed cycle vertex and at most one leaf event
per such vertex. Only a constant number of event keys is inspected each
time. Each prefix record is written once and read once in reconstruction;
at most three constant-dimensional leaf updates occur after closure.
Scanning a center and classifying its neighbors costs $O(1+d_i)$, including
all neighbor degree replies. Outputting its leaves and reaccessing stored
labels has the same bound. All exposed centers are in $S$, so these costs
are bounded by $O(1+\vol(S))$, even if some inspected leaf or candidate
labels are outside $S$. A deterministic balanced label dictionary adds at
most logarithmic overhead.

For a leaf seed $s$, first return zero if $\lambda\geq1$. Otherwise its
conditional value at its neighbor $v$ of potential $t$ is
$u_s=1-\lambda+\gamma t>0$. Read the seed's one-entry row and $d_v$.
Eliminating $s$ changes the root data to
\[
 \delta_v=d_v-\gamma^2,\qquad
 b_v=\gamma(1-\lambda)-\lambda d_v.
\]
If $b_v\leq0$, only $s$ is positive, with value $1-\lambda$, and the
root row is not read. Otherwise apply the same cycle-root algorithm with
these modified data, excluding $s$ from its negative leaf group; recover
$s$ by its displayed formula. The additional work is constant and all
exposed rows still belong to positive output coordinates.

Finally sum the optimal equations on $S$. For nonempty support the seed
lies in $S$ and
\[
 1-\lambda\vol(S)
 =\bar\alpha\sum_{i\in S}d_i u_i
   +\gamma\sum_{\substack{\{i,j\}\in E\\i\in S,\ j\notin S}}u_i
 \geq0.
\]
This proves the support-volume bound and completes the accounting.
\end{proof}

\begin{corollary}[The OP3 approximation scale on this graph class]
\label{cor:op3-local-cycle-acl}
Let $\lambda=\epsappr\in(0,1)$ and let
$\widehat{\bm\pi}=\bar\alpha\bm D\bm u$ be the output of
\cref{thm:op3-local-cycle}. For the shared lazy PageRank vector $\bm\pi$,
\[
 \bm0\leq\bm\pi-\widehat{\bm\pi}\leq\epsappr\bm d.
\]
The charged work is $\widetilde O(1/\epsappr)$ in the exact-real word
model on this promised graph family.
\end{corollary}
\begin{proof}
Put $\bm r=\bm e_v-\bm M\bm u$. On positive coordinates
$r_i=\lambda d_i$; on zero coordinates the obstacle conditions and
nonpositive off-diagonals give $0\leq r_i\leq\lambda d_i$.
The empty-support case has the same inequality. Since
$\bm M\one=\bar\alpha\bm d$ and $\bm M^{-1}\geq0$,
\[
 \bm0\leq\bm\pi-\bar\alpha\bm D\bm u
 =\bar\alpha\bm D\bm M^{-1}\bm r
 \leq\lambda\bm d.
\]
Apply the preceding work and volume bounds.
\end{proof}

\paragraph{Measured evidence and limits.}
The exact local routine is \texttt{local\_sun\_solver.py}. The audit
enumerates ordered leaf counts in $\{0,1,2\}$ on cycles of length three
through five, cycle and leaf seeds, and nine parameter pairs: 5,265
comparisons with a separate exact obstacle routine pass. Full-graph KKT
checks are audit-only and excluded from solver counters. Larger asymmetric
cycles demonstrate local stopping without scanning the ambient graph.
This is a family theorem, not arbitrary-graph OP3 and not a claim to realize
every intermediate Wei--Yang active-set gate. Its decisive special property
is that an interior attachment settles before its center leaves the frontier.
Arbitrary attached trees can violate that ordering; identifying a paid
replacement for their changing responses is the next question.

\begin{proposition}[A length-two attachment need not settle; proof draft]
\label{prop:op3-cycle-unsettled-attachment}
The automatic-settling claim fails if pendant leaves are replaced by
arbitrary pendant trees. This refutes that extension of the record rule,
not the existence of an efficient solver on the larger graph class.
\end{proposition}
\begin{proof}
Attach a path $c-a-b$ to a cycle center $c$, where $d_a=2$ and $d_b=1$.
Take $\gamma=3/4$. While $b$ is inactive, the branch response is
$u_a=(\gamma u_c-2\lambda)_+/2$. Thus $a$ activates at
$u_c=8\lambda/3$, and $b$ activates only at $u_c=56\lambda/9$.
A next one-sided cycle neighbor of degree three activates at
$u_c=4\lambda$, strictly between those thresholds. At that point
$u_a=\lambda/2$ and $u_b=0$. Eliminating this branch subtracts $9/32$
from the center diagonal before $b$ activates, but $9/23$ after it
activates. A single fixed quadratic record cannot describe both regimes.

This ordering occurs on an actual source continuation. Use the cycle
$0,1,\ldots,11,0$, add edges $\{1,12\},\{12,13\},\{2,14\}$, seed at
$0$, and $\lambda=1/1000$. The exact conditional solve at
$u_0=39/2000$ gives $u_1=1/250$, $u_2=0$, $u_{12}=1/2000$, $u_{13}=0$
and reduced source derivative $-10639/11000<0$. At $u_0=1/20$ it gives
$u_2=307/118500>0$, $u_{13}=89/31600>0$ and source derivative
$-2147709/2325760<0$. The full vectors and exact conditional
solves are recorded in the audit. Both events therefore precede
the final root, and the old branch changes after the next cycle admission.
\end{proof}

\section{A bounded region of changing tree attachments}
\label{sec:op3-bounded-attachments}

\paragraph{Status and model.}
The following are \statusproved{} proof drafts, awaiting independent review.
The graph is a simple cycle with disjoint finite trees attached to its
centers through one edge each. The seed is on the cycle. Neither the cycle
ordering nor any attachment size is supplied. Use the degree obstacle
problem with $\bm M=\bm D-\gamma\bm A$,
$\bm b=\bm e_v-\lambda\bm d$, $0<\gamma<1$, and $\lambda>0$.
Let $q\geq1$ be the maximum size of an attachment incident to a cycle
center that is positive in the final answer, taking $q=1$ if there are no
such attachments. This is an analysis parameter, not advice to the solver.
The exact-real word model charges all bounded probes, even into inactive
vertices, all curve construction, stored-state work, event searches and
output. No coefficient-bit or floating-point claim is made.

\subsection{When a tree attachment has finished changing}

\begin{lemma}[A discounted-hitting saturation bound]
\label{lem:op3-attachment-saturation}
Let a tree $T$ of at most $q$ vertices be attached through one edge to a
center whose potential is fixed at $x\geq0$. All loads in $T$ are
$-\lambda d_i$, using actual graph degrees. Its conditional obstacle
solution has every coordinate strictly positive, and its response has
reached the final affine piece, whenever
\begin{equation}
 x\geq X_q:=\frac{\lambda}{1-\gamma}
                    \bigl(\gamma^{-2q^2}-1\bigr).
 \label{eq:op3-attachment-saturation}
\end{equation}
The algorithm need not evaluate this threshold.
\end{lemma}
\begin{proof}
Run simple random walk until it first hits the center, with hitting time
$\tau$, and put $p_i=\mathbb E_i[\gamma^\tau]$. Harmonic conditioning on
the first step gives $p_i=\gamma\sum_{j\sim i}p_j/d_i$, with center value
one. Direct substitution therefore shows that the unconstrained solution is
\begin{equation}
 u_i=p_i\left(x+\frac{\lambda}{1-\gamma}\right)
                -\frac{\lambda}{1-\gamma}.
 \label{eq:op3-discounted-attachment}
\end{equation}

Root the attachment at the center. The expected time from a child $j$ to
its parent is $2|T_j|-1$. Indeed, first-step decomposition gives
$H_j=d_j+\sum_{k\text{ child of }j}H_k$; induction over the descendant
tree gives the displayed expression. The walk from any $i$ to the center
crosses the successive parent edges in order. The strong Markov property
then gives
$\mathbb E_i\tau\leq q(2q-1)<2q^2$.
The function $t\mapsto\gamma^t$ is convex, so
$p_i\geq\gamma^{\mathbb E_i\tau}>\gamma^{2q^2}$.
Under \cref{eq:op3-attachment-saturation},
\cref{eq:op3-discounted-attachment} is strictly positive for every $i$.
It is therefore the obstacle optimum. It remains positive at larger $x$,
so the response thereafter is this single affine formula.
\end{proof}

\begin{lemma}[Potential growth behind the active end]
\label{lem:op3-cycle-interior-growth}
Before cycle closure, let $y_0\geq0$ be the potential at an active end,
$y_{-1}=0$ the next inactive cycle coordinate, and $y_k$ the consecutive
active-center potentials toward the seed. Up to and including the seed,
\begin{equation}
 y_k\geq\frac{\lambda}{1-\gamma}(\gamma^{-k}-1).
 \label{eq:op3-cycle-interior-growth}
\end{equation}
Consequently all attachments at a center at least $2q^2$ edges behind an
active end have finished changing.
\end{lemma}
\begin{proof}
A negative-load attachment has no coordinate greater than its center:
a larger positive maximum would contradict its own obstacle equality and
$\gamma<1$. At a nonseed cycle center of degree $d_i=2+k_i$, the sum of
the $k_i$ attached-root potentials is at most $k_i y_i$. Its active
equation therefore implies
\[
 y_{i+1}+y_{i-1}
 =\frac{2+k_i}{\gamma}(y_i+\lambda)
     -\sum_{j\text{ attached at }i}u_j
 \geq\frac2\gamma(y_i+\lambda).
\]
Put $w_i=y_i+\lambda/(1-\gamma)$. Then
$w_{i+1}\geq2w_i/\gamma-w_{i-1}$, and
$w_0\geq w_{-1}=\lambda/(1-\gamma)$.
The first step gives $w_1\geq w_0/\gamma$. If
$w_i\geq w_{i-1}/\gamma$, the recurrence gives
$w_{i+1}\geq(2/\gamma-\gamma)w_i\geq w_i/\gamma$.
Induction proves \cref{eq:op3-cycle-interior-growth}. This uses only
nonseed equations even when the last derived value is the seed value.
Apply \cref{lem:op3-attachment-saturation} at distance $2q^2$.
\end{proof}

The length-two counterexample in
\cref{prop:op3-cycle-unsettled-attachment} is consistent with these bounds.
An attachment can still change after one advance; it cannot remain
unsettled arbitrarily far behind the active end when its size is bounded.

\subsection{Finding ports and building only bounded components}

At an admitted cycle center $c$, read its adjacency row and try exploration
limits $L=1,2,4,\ldots$. For each neighbor $j$, explore its component in
$G-c$ until it is exhausted or one of the following certificates occurs:
more than $L$ vertices are found, a vertex has degree greater than $L$, or
another neighbor of $c$ is reached. In the degree test, stop before reading
that vertex's adjacency row. An exhausted component is an attached tree of
size at most $L$. The other cases provisionally mark a port. Reaching
another neighbor marks both of the corresponding ports.

The actual two cycle ports are always marked: their common component in
$G-c$ either exceeds the limit or connects the two neighbors. An attached
tree can be marked only when it exceeds $L$. Thus exactly two marked
ports certify that all attached components have been fully discovered.
Otherwise double $L$ and retry, caching graph replies. If the largest
attachment at $c$ has size $q_c\geq1$, this terminates with $L<2q_c$,
or with $L=1$ when there is no attachment. Each trial explores at most
$L$ rows of degree at most $L$ per incident neighbor, plus a constant
number of stopping replies. Its charged work is $O(L^2d_c)$; all trials
sum to $O(q_c^2d_c)$. No ambient graph scan or structural advice is used.

Build each discovered attachment's response once using
\cref{thm:op3-persistent-tree}, retaining its child curves for terminal
reconstruction. It is permissible to build already exhausted small
components again during failed classification trials, provided this work
is charged; the implementation does so. Each attachment response is
convex and has at most its size in knots. For a center, sort the union of
its attachment knots in the center's own potential and maintain one cursor.
If its current summed response is $C_i x+E_i$, the effective center
diagonal and load are $d_i-C_i$ and $-\lambda d_i+E_i$; the seed also has
its source load one. Crossing a hinge $c(x-\tau)_+$ subtracts $c$ from
that diagonal and $c\tau$ from that load. Equal knots are consumed together.

\subsection{Eagerly store settled prefixes and retain the changing region}

Run the exact source-potential continuation of
\cref{sec:op3-local-cycle}. On each arm retain a consecutive window of
centers adjacent to the active end. After any event, repeatedly eliminate
the first center of the window if its attachment cursor is exhausted,
always leaving at least the last center. Eliminate from the seed outward,
store one fixed reverse record, and accumulate its source-derivative
contribution as in \cref{eq:op3-cycle-append,eq:op3-cycle-root-derivative}.
The rule depends on the exhausted cursor, not on knowing $q$.

If a window has more than one center, its first center has an unsettled
attachment. By \cref{lem:op3-cycle-interior-growth}, its distance to the
active end is less than $2q^2$. Therefore each retained arm has at most
$2q^2$ centers. The attached components have already been compressed to
their current scalar affine pieces. Their remaining cycle system is
tridiagonal, with affine loads in the source potential $t$. Two forward
and reverse passes compute all retained response coefficients in $O(q^2)$
work. The exact root derivative follows from their Schur contributions.

Inspect only the next attachment knot at each retained center, the seed's
next knot, and the one or two cycle candidates. Every knot time is the
zero of a known affine equality. The earliest event or earlier root of
the reduced derivative is therefore found in $O(q^2)$ work, including
the no-event certificate. A fully settled eliminated center has no future
attachment event. No old prefix coordinate is queried or rewritten.

If the two candidates coincide, use their combined load for the closing
gate. After its admission, the two windows and the last center form one
tridiagonal path, ordered from one seed-side prefix to the other, of size
at most $4q^2+1$. Retain this path for the rest of the continuation.
It contains every remaining nonseed attachment event. The root's own
knots are still checked directly. Compute the two affine right-hand sides
and the first remaining event by linear passes in this bounded path.
Finish with one reverse pass through all stored prefixes and one traversal
of the retained attachment curves.

\begin{theorem}[Local solve with bounded tree attachments; proof draft]
\label{thm:op3-bounded-attachments}
For a cycle-seeded graph in the stated family, the described algorithm
returns the exact obstacle optimum with
\begin{equation}
 \widetilde O\bigl(q^3(1+\vol(S))\bigr)
 \quad\text{word work and}\quad
 \widetilde O\bigl(q^2(1+\vol(S))\bigr)
 \quad\text{working storage},
 \label{eq:op3-bounded-attachment-work}
\end{equation}
where $S$ is its positive support. The quantity $q$ and the graph structure
need not be supplied. Admitted center rows and all bounded probes are
charged. Taking $\lambda=\epsappr$ gives an ACL approximation with work
$\widetilde O(q^3/\epsappr)$ and no polynomial dependence on $1/\alpha$.
For every fixed attachment-size bound this attains the OP3 scale on the
promised family. It is not an arbitrary-graph or an arbitrary-seed theorem.
\end{theorem}
\begin{proof}
Port classification is exact under the graph promise. At each admitted
center the supplied bounded components therefore give its entire exact
conditional response. The two-arm description enumerates all cycle gates
before closure, including combined load at the last inactive center.
All other possible events are attachment knots. Eliminated centers have
exhausted their lists; the retained windows and the seed contain the rest.
The source derivative, response equations and knot times are exact Schur
computations on these current pieces. The positive Schur derivative and
continuity justify advancing precisely when its zero lies after the next
event. Thus every admitted center is positive at the final root; equality
causes stopping before a new row is admitted. Final reconstruction solves
all active equations and all other obstacle inequalities.

Summing center degrees is bounded by $\vol(S)$. Classification, including
failed bounded probes and reaccesses, costs $O(q^2\vol(S))$. The number
of attachment vertices built is at most $q$ times the number of attached
components, hence $O(q\vol(S))$. Repeated bounded builds and sorting cost
$\widetilde O(q\vol(S))$. The number of consumed knots and center
admissions is $O(q\vol(S))$. Each costs $O(q^2)$ for retained factors,
responses and event searches, including unsuccessful searches. Each
committed record is created once and used once at the end; even array
shifts of retained windows cost only $O(q^2)$ per committed record.
Terminal attached-curve evaluations cost
$\widetilde O(q\vol(S))$ and positive output writes cost $O(\vol(S))$.
These terms give the work bound. Cached probe incidences, persistent
bounded-component curves, current knot lists, center records and temporary
tridiagonal arrays fit the stated working storage bound; discarded
construction and solve allocations are already charged to work.

Finally the full-graph obstacle equations give
$\lambda\vol(S)\leq1$, and the inverse-positive residual argument of
\cref{cor:op3-local-cycle-acl} gives the ACL error
$\bm0\leq\bm\pi-\bar\alpha\bm D\bm u\leq\lambda\bm d$.
\end{proof}

\paragraph{Measured evidence and remaining question.}
The exact bounded-attachment implementation passes 1,909
comparisons against a separate exact obstacle solver, including branching
attachments and the preceding length-two witness. Eight larger instances
pass full-graph KKT checks, performed only by the audit. There are also
195 exact discounted-hitting checks on rooted trees through size six.
For sizes two and three the reported larger runs retain at most five cycle
centers at once; the theorem uses the conservative bound $4q^2+1$.
The implementation counts failed classification probes, repeated curve
construction and every retained-window factor row. These finite checks do
not replace independent proof review. The remaining OP3 question is to
remove the growing-attachment factor or replace it by a paid approximate
response certificate, and then to handle more than one cyclic core.

\section{Support-local elimination with two retained vertices}
\label{sec:op3-two-port-frontier}

\paragraph{Status, scope and provenance.}
The results below are \statusproved{} proof drafts awaiting independent
review. Use the degree obstacle with
$\bm M=\bm D-\gamma\bm A$ and
$\bm b=\bm e_v-\lambda\bm d$, and let $S$ be its positive support.
Assume
\begin{equation}
 \#\{i\in S\setminus\{v\}:d_i\geq3\}\leq1.
 \label{eq:op3-two-port-promise}
\end{equation}
This is a condition on the actual positive support. No separator, cycle
ordering, path lengths, attachment sizes or ambient graph size is supplied.
The rest of the graph may have arbitrary degrees and cycles.

Earlier notes already prove the stable fixed-port reporter principle
and exact rooted-spider elimination. The companion probe identifies their
precise source files and theorem labels. The new contribution here is a locally
discovered graph condition that guarantees only individually charged rekeys,
including shared inactive candidates and multiple cycles. We give the
algebra and an implemented, conservative planar reporter explicitly. These
are contextual pointers, not imported proof dependencies.

\subsection{The retained matrix and stable frontier rows}

Retain the seed and, when it is first admitted, the first nonseed vertex
of degree at least three. Call this set $C$, of size at most two. Eliminate
every other admitted vertex immediately. Write $E$ for the eliminated set
and $B$ for the exposed, inactive frontier. Cache a frontier record
\begin{equation}
 \delta_j>0,\qquad h_j<0,\qquad \bm a_j\in\mathbb R^2_{\geq0},
 \qquad g_j=h_j+\bm a_j^{\mathsf T}\bm u_C.
 \label{eq:op3-two-port-row}
\end{equation}
All retained quantities use two fixed slots. Before the second vertex
exists, its diagonal is one, its load and couplings zero; it is merely
padding. The stored Schur matrix $\bm H_C$ and load $\bm f_C$ then give the
exact retained values $\bm u_C=\bm H_C^{-1}\bm f_C$ in constant work.

\begin{lemma}[Only one outside successor at a low-degree admission]
\label{lem:op3-two-port-stability}
Every legal admitted vertex of degree at most two has at most one original
neighbor still inactive. Eliminating it changes $\bm H_C,\bm f_C$ and at most one
surviving frontier record. Every other record in
\cref{eq:op3-two-port-row} remains exactly unchanged.
\end{lemma}
\begin{proof}
A positive gate requires a positive original active neighbor, because the
nonseed load is negative. At most one of the two original neighbors is
therefore still inactive. Inductively, no eliminated path creates a fill
edge between two inactive vertices: just before elimination the pivot is
coupled only to retained vertices and that possible single inactive
successor. Original edges between two inactive vertices are irrelevant to
the restricted solve and are revealed only when one endpoint is admitted.

For pivot $i$, abbreviate its record by $(\delta,h,\bm a)$, and let $j$
be its possible inactive successor, with original conductance $w=\gamma$.
Scalar Schur complementation gives
\begin{align}
 \bm H_C&\leftarrow \bm H_C-\bm a\bm a^{\mathsf T}/\delta,
 & \bm f_C&\leftarrow \bm f_C+\bm a h/\delta,
 \label{eq:op3-two-port-core-update}\\
 \delta_j&\leftarrow\delta_j-w^2/\delta,
 & h_j&\leftarrow h_j+wh/\delta,
 & \bm a_j&\leftarrow\bm a_j+w\bm a/\delta.
 \label{eq:op3-two-port-frontier-update}
\end{align}
The second line is absent without a successor. It includes any prior
contribution to $j$ from another path; such a shared candidate is never
treated as independent copies. Positive Schur pivots follow from positive
definiteness. Initially $h_j=-\lambda d_j<0$, and the displayed recurrence
preserves its sign and nonnegative couplings. All changes are confined to
the displayed constant-size records.
\end{proof}

When a degree-at-least-three candidate $i$ is first admitted, place its
existing diagonal and load in the second retained slot, with coupling
$-a_{i,1}$ to the seed slot. Scan its original row, adding new original
conductances to the appropriate second coupling of each inactive neighbor.
Its old neighbors already admitted have had their edges accounted for.
An edge dictionary prevents counting an original edge twice. Old frontier
records unrelated to this row keep second coupling zero. There is no bulk
change of their coordinates. If another nonseed branching candidate has a
positive gate later, stop with a certificate that
\cref{eq:op3-two-port-promise} fails, before reading that candidate's row.

\subsection{An implemented planar reporter with a conservative bound}

Normalize each frontier record to the point
\begin{equation}
 \bm p_j:=\bm a_j/(-h_j)\in\mathbb R^2_{\geq0}.
 \label{eq:op3-two-port-point}
\end{equation}
Then $g_j>0$ exactly when
$\bm p_j^{\mathsf T}\bm u_C>1$. One extreme-point query either produces
a legal singleton admission or certifies all frontier gates nonpositive.
Only insertions, deletions and individual replacements of points are used.
The retained solve changes the query direction, not the stored points.

\begin{lemma}[A fully charged dynamic upper-hull implementation]
\label{lem:op3-avl-upper-hull}
An initially empty planar point dictionary can support a labeled point
insertion, deletion or replacement in $O(\log^4(2+N))$ word work and an
extreme query in any nonnegative direction in $O(\log^2(2+N))$ work,
where $N$ is the largest dictionary size so far. The working storage is
$O(N\log(2+N))$. Over $K$ operations, allocations and discarded-state
work are bounded by $O(K\log^4(2+N))$.
\end{lemma}
\begin{proof}
Maintain an AVL search tree on lexicographic keys $(x,y,\text{label})$.
Each node stores its subtree's upper convex chain as an immutable AVL
sequence, with subtree sizes and heights. Joining, splitting and indexed
selection of such sequences use standard height recursion; even allowing
$O(\log^2(2+N))$ for split/join is sufficient below. These operations copy
only their recursion paths, never the whole sequence.

It remains to charge merging the upper chains $A,B$ of horizontally
separated point sets. Both chains have strictly increasing $x$ and strictly
decreasing edge slopes. If their boundary $x$ coordinates agree, discard
the lower endpoint there, with a fixed tie rule, using a sequence split.
After this the remaining nonempty chains are strictly separated.
For $A_i$, put
\[
 F_i=\max_j\operatorname{slope}(A_i,B_j).
\]
The slopes to successive $B_j$ first increase and then decrease, by
concavity of $B$. A binary search gives the last maximizing $j$ in
$O(\log^2(2+N))$ work, counting every indexed sequence access.

The bridge endpoint in $A$ is the first $i$ for which the next edge slope
of $A$ is at most $F_i$, taking the last index if needed. This predicate
is monotone: if that edge slope is at most $F_i$, the same maximizing
point of $B$ has slope at least $F_i$ from $A_{i+1}$, while later edge
slopes in $A$ only decrease. Conversely, when an edge slope exceeds
$F_i$, every slope to $B$ decreases after moving across that edge.
At the first passing index, the preceding edge slope is therefore at least
$F_i$. The resulting line supports both chains from above. Keep the prefix
through $A_i$ and the suffix from the last maximizing $B_j$; their
concatenation is the merged upper hull. The outer binary search costs
$O(\log^3(2+N))$, including the inner search and indexed accesses.

An AVL point update changes $O(\log(2+N))$ nodes and rotations. Rebuilding
each node's chain by two separated merges gives the stated fourth-power
bound. Along a stored upper chain, dot product with a nonnegative direction
first increases and then decreases, so indexed binary search gives the
query bound. Equal points and vertical coincidences follow the stated
endpoint rule; every label still exists in the point dictionary for later
deletion.

At one depth of the point search tree the represented subtrees are
disjoint, so their total chain lengths are at most $N$. Summing over its
logarithmic height gives the working-storage bound, even without credit
for shared immutable sequence nodes. Each update allocates and discards
only the charged recursion records. Reclamation costs can be charged once
to each allocation. No old version is retained by the solver.
\end{proof}

The exact implementation is \texttt{dynamic\_upper\_hull.py}; its simpler
fourth-power logarithm is adequate for OP3's soft-$O$ target. As a
\statussource{} alternative, Jacob--Brodal's dynamic planar convex hull
supports individual insertions and deletions in amortized logarithmic time
and extreme queries in logarithmic time; see \citet{jacob2019dynamic}, Theorem~1 and
computational model, Sections~1--2, arXiv:1902.11169v1, PDF pages~2 and~5--6.
The implementation here does not import their more elaborate construction.
Neither primitive supplies the unresolved bulk-pullback-and-meld reporter
for a changing general separator hierarchy.

\subsection{The resulting local theorem}

\begin{theorem}[Two retained vertices; exact local proof draft]
\label{thm:op3-two-port-frontier}
Under \cref{eq:op3-two-port-promise}, the described algorithm returns the
exact obstacle optimum with
\begin{equation}
 O\bigl((1+\vol(S))\log^4(2+\vol(S))\bigr)
 \label{eq:op3-two-port-work}
\end{equation}
exact-real word work and
$O((1+\vol(S))\log(2+\vol(S)))$ working storage. It reads precisely the
original adjacency rows of positive output vertices and may query degrees
of their inactive neighbors. There is no polynomial dependence on
inverse teleportation or on inactive component sizes. With
$\lambda=\epsappr$, the output $\widehat{\bm\pi}=\bar\alpha\bm D\bm u$
is an ACL approximation and the work is
$\widetilde O(1/\epsappr)$ on this support-defined graph class.
\end{theorem}
\begin{proof}
If $\lambda d_v\geq1$, all loads are nonpositive, and zero is the optimum;
only the seed degree is read. Otherwise the singleton seed solution is
strictly positive. At every stage the retained matrix and saved
eliminations represent the exact restricted solution on the admitted set.
Its nonseed boundary gates are exactly
\cref{eq:op3-two-port-row}, by Schur elimination.

A strict positive gate admits a vertex with a positive Schur value. Inverse
positivity makes every old admitted coordinate increase. Moreover every
restricted positive coordinate is bounded above by the full optimum:
on an admitted set $U$, the full optimum satisfies
$\bm M_{UU}\bm u^*_U\geq\bm b_U$, hence
$\bm u^*_U\geq\bm M_{UU}^{-1}\bm b_U$.
If a boundary vertex were zero at the full optimum, its full neighbor flux
would be at most $\lambda d_j$, so its restricted gate could not be
positive. Induction therefore puts every admission in $S$. The promised
second branching admission cannot occur.

Every successful query deletes one admitted label, and the final failed
query certifies every exposed boundary inequality, including shared
candidates. An unexposed vertex has no admitted neighbor and has negative
load, so its inequality holds too. One reverse pass through the records
\[
 u_i=\frac{h_i+\bm a_i^{\mathsf T}\bm u_C+w u_j}{\delta_i}
\]
recovers the exact positive coordinates, taking $u_j=0$ if the optional
successor never activated. This gives all KKT conditions and hence the
unique optimum.

Each admitted row is scanned once. Degree replies, original-edge dictionary
operations and boundary exposures total $O(1+\vol(S))$, with logarithmic
dictionary overhead allowed. The seed and at most one other retained row
contribute at most their degrees in point updates. Every other admission
has constant-size Schur work and changes at most one surviving point.
There are $O(1+\vol(S))$ point operations, retained solves, queries and
reverse records. Apply \cref{lem:op3-avl-upper-hull}, charging all
allocations, failed queries and output words. Inactive adjacency rows are
never scanned. Finally, the obstacle mass identity gives
$\lambda\vol(S)\leq1$, and its nonnegative residual satisfies
$\bm0\leq\bm e_v-\bm M\bm u\leq\lambda\bm d$.
The ACL conclusion follows as in \cref{cor:op3-local-cycle-acl}.
\end{proof}

\paragraph{Measured evidence and the remaining limitation.}
The graph implementation passes 8,329 exact comparisons on connected atlas
graphs through seven vertices, every seed and four parameter pairs. Another
18,791 cases correctly detect that the support condition fails. Nine larger
instances pass separate full-graph KKT audits. With a growing inactive
descendant tree, ambient sizes 73, 265 and 1,033 all give the same eight
row scans, 17 adjacency entries and nine degree replies, while an attachment
really activates. A shared-report graph with 66,162 vertices requires only
98 positive row scans and 258 adjacency entries. Unequal path bundles with
up to 128 paths exercise multiple cycles and shared candidates. Reference
solves and global validation are excluded from local operation counters.

The new structural result does not handle a growing number of admitted
branching vertices. Keeping more vertices raises reporter dimension; hiding
them behind a changing hierarchy recreates bulk transformation and event
location obligations. Those remain separate from the successfully paid
two-vertex construction.

\section{A growing branching core with scalar flux events}
\label{sec:op3-branch-core-flux}

\paragraph{Status and purpose.}
This section is a \statusproved{} proof draft awaiting independent review.
It combines the original-residual publication bands of
\cref{sec:op3-geometric-events} with low-degree path elimination.
The number of retained vertices may now grow. A fully specified dense
inverse pays for their response changes, leaving an explicit quadratic
core-size factor. No high-dimensional geometric reporter is required.
General OP3 remains \statusopen.

Set $\lambda=\epsappr/2$, $\bm M=\bm D-\gamma\bm A$ and
$\bm b=\bm e_v-\lambda\bm d$. At each stage, $U$ is the admitted set
and $\bm u^U_U=\bm M_{UU}^{-1}\bm b_U>0$, with zero outside $U$.
Retain the seed and every admitted vertex of original degree at least three
in a set $C$. Immediately eliminate all other admitted vertices into saved
scalar equations. Let $r=|C|$ at termination, $n_U=|U|$, and $V=\vol(U)$.
These are measured properties of the returned support, not input advice.

\subsection{One retained endpoint per unfinished path}

\begin{lemma}[Stable physical-flux arms]
\label{lem:op3-scalar-flux-arms}
For every original edge $ij$ with $i\in U$ and $j\notin U$, the physical
flux $F_{ij}=\gamma u_i^U$ has a representation
\begin{equation}
 F_{ij}=a_{ij}u_{c(ij)}^U+b_{ij},\qquad
 a_{ij}>0,\quad b_{ij}\leq0,\quad c(ij)\in C.
 \label{eq:op3-physical-flux-arm}
\end{equation}
The coefficients and core label are fixed during this edge's lifetime on
the boundary. An admission creates at most one such record per newly scanned
original edge, and retires the records ending at its vertex. The total number
of arm creations and retirements is $O(V)$.
\end{lemma}
\begin{proof}
An eliminated vertex has original degree at most two. Every connected
component of eliminated vertices is consequently a path, and it touches the
retained set: the admitted set is connected to the retained seed. If a path
also touches an inactive vertex, it has exactly one retained endpoint.
It cannot have an additional inactive or retained exit. Components joining
two retained endpoints have no inactive exit and produce no frontier arm.

A direct retained-to-inactive edge has $(a,b)=(\gamma,0)$.
Suppose an admitted degree-two vertex $i$ has an inactive successor $j$.
It has exactly one original active neighbor, so its old frontier record
contains exactly one arm $a u_c+b$. With current Schur diagonal $\delta_i$,
its equation, keeping $j$ as a formal outside variable, is
\[
 \delta_i u_i=-\lambda d_i+b+a u_c+\gamma u_j.
\]
Eliminating $i$ creates the new physical-flux arm, evaluated at $u_j=0$,
\[
 F_{ij}=\frac{\gamma a}{\delta_i}u_c
       +\frac{\gamma(-\lambda d_i+b)}{\delta_i}.
\]
Its diagonal contribution at $j$ is $-\gamma^2/\delta_i$.
Positive Schur pivots preserve the coefficient signs. Nothing changes an
existing arm's internal path until its target is admitted, when that arm
is retired. Other admissions change only its retained endpoint potential.
Each original edge is first seen when one endpoint's row is scanned;
an edge dictionary avoids counting its reverse appearance again.
\end{proof}

Several arms may end at the same inactive vertex, including arms with
different retained endpoints. Their contributions are summed in one
candidate record. In particular, the exact gate is
\begin{equation}
 g_j=-\lambda d_j+\sum_{ij:\,i\in U}F_{ij}.
 \label{eq:op3-arm-original-gate}
\end{equation}
It is not a collection of independent-copy gates. A candidate's stored
Schur diagonal includes all prior path eliminations into that candidate.

\subsection{Scalar queues certify the original residual scale}

For each arm keep a lower published flux $\ell_{ij}$, initially zero.
Use $h=\epsappr/16$ and $\kappa=5/4$. The next publication is due when
\begin{equation}
 u_{c(ij)}^U>
 \tau_{ij}:=\frac{\kappa\ell_{ij}+h-b_{ij}}{a_{ij}}.
 \label{eq:op3-arm-next-event}
\end{equation}
A minimum heap for each retained coordinate stores these thresholds.
After updating the retained potentials, visit each heap and process all
due entries. Publish the actual flux $\ell_{ij}\leftarrow F_{ij}$, update
the candidate's sum $L_j=\sum_i\ell_{ij}$, and insert the new threshold.
Retired arms may be removed lazily: every discarded entry is charged once.

When all heaps are settled, every live arm satisfies
\[
 \ell_{ij}\leq F_{ij}\leq\tfrac54\ell_{ij}+\epsappr/16.
\]
Queue a candidate as soon as $L_j>(11/20)\epsappr d_j$. This is a legal
admission because its exact gate in \cref{eq:op3-arm-original-gate} is
strictly positive for $\lambda=\epsappr/2$. A queued candidate stays legal:
all restricted solutions and physical fluxes increase under legal admissions.
When the ready queue is empty after settling all heaps,
\begin{equation}
 0\leq r_j^U=\sum_i F_{ij}
 \leq\tfrac54L_j+\tfrac1{16}\epsappr d_j
 \leq\tfrac34\epsappr d_j
 \quad (j\notin U).
 \label{eq:op3-arm-terminal-residual}
\end{equation}
An unexposed vertex has no active neighbor and residual zero.
Active residuals are exactly $\lambda d_i$. Thus the exact restricted
solution already supplies an ACL output; a further solve is unnecessary
in this exact-arithmetic implementation.

\begin{lemma}[All scalar reporting work is paid]
\label{lem:op3-scalar-flux-work}
Write $L=1+\log_+(1/(\bar\alpha\epsappr))$. Over the entire run the arm
queues perform $O(VL)$ publications, heap insertions and heap removals,
plus $O(n_Ur+VL)$ top inspections. Their work, including candidate sums,
ready entries and failed event searches, is
$O(n_Ur+VL\log(2+V))$, with $O(V+r)$ stored words.
\end{lemma}
\begin{proof}
Conservation on any positive restricted face gives
$\bar\alpha\sum_{i\in U}d_i u_i^U\leq1$, hence
$F_{ij}\leq\gamma/\bar\alpha$. A first publication exceeds $h$;
each later one exceeds $5/4$ times its predecessor. Each fixed arm therefore
publishes $O(L)$ times. There are $O(V)$ arms by
\cref{lem:op3-scalar-flux-arms}. A publication pops its current entry and
pushes exactly one replacement. Retiring an arm leaves at most one stale
entry. Thus all stale removals and allocated heap records are paid.
Each round scans at most $r$ heaps, and every round except the last admits
one vertex. One failed top comparison per nonempty heap is included.
Every candidate enters the ready queue at most once. All arm records,
candidate records and possibly stale heap entries occupy $O(V+r)$ words.
\end{proof}

This argument uses individual physical fluxes, not relative error after
division by a candidate's accumulated negative Schur load. The latter
contract fails on the source-valid path examples in
\texttt{SCHUR\_RELATIVE\_REPORTER\_PROBE.md}. Large cancellation in an
arm's affine formula is allowed in the exact-real word model; coefficient
bit cost and finite-precision certification remain outside this theorem.

\subsection{An explicit core backend and the resulting bound}

Maintain the inverse $\bm P=\bm H_C^{-1}$ and retained values $\bm u_C$.
At an admission the candidate has a Schur diagonal $\delta$, load $b<0$
and nonnegative sparse coupling vector $\bm a$. Put
\begin{equation}
 \bm z=\bm P\bm a,\qquad s=\delta-\bm a^{\mathsf T}\bm z>0,
 \qquad t=(b+\bm a^{\mathsf T}\bm u_C)/s>0.
 \label{eq:op3-core-rank-update}
\end{equation}
All old retained values change by $\bm u_C\leftarrow\bm u_C+t\bm z$.
If the candidate has degree at most two, eliminate it and update
$\bm P\leftarrow\bm P+\bm z\bm z^{\mathsf T}/s$.
If it is a branching vertex, retain it with value $t$ and bordered inverse
\[
 \begin{pmatrix}
 \bm P+\bm z\bm z^{\mathsf T}/s&\bm z/s\\
 \bm z^{\mathsf T}/s&1/s
 \end{pmatrix}.
\]
These are ordinary scalar block-inverse identities. A low-degree candidate
has at most two arm coordinates; at a branching admission, forming its
couplings costs at most its original degree. Every dense inverse entry
updated, every retained value written and every inverse allocation is
charged. A reverse pass through the saved low-degree equations recovers
all eliminated values in $O(V)$ work.

\begin{theorem}[Local ACL with a growing branching-core parameter]
\label{thm:op3-branch-core-flux}
On every graph in the canonical local access model, the described algorithm
returns a nonnegative ACL approximation at tolerance $\epsappr$. It scans
precisely the rows of its positive output set $U$ and queries only the seed
degree and degrees of exposed neighbors. For a nonempty output, its exact-real
word work is bounded by
\begin{equation}
 O\left(n_Ur^2+rV+VL\log(2+V)\right),\qquad
 L=1+\log_+\frac1{\bar\alpha\epsappr},
 \label{eq:op3-branch-core-work}
\end{equation}
and its working storage is $O(V+r^2)$. Empty output costs $O(1)$.
Moreover $V\leq2/\epsappr$, so the work is
$\widetilde O((1+r^2)/\epsappr)$. In particular, this attains the OP3
soft-$O(1/\epsappr)$ scale when the returned support has at most a
polylogarithmic number of original branching vertices.
\end{theorem}
\begin{proof}
If $\lambda d_v\geq1$, zero is feasible with the required residual bound.
Otherwise the positive singleton seed initializes the construction. The
Schur identities and \cref{lem:op3-scalar-flux-arms} preserve exact restricted
solutions. Strict positive gates give positive new coordinates and monotone
old ones by inverse positivity. As in \cref{thm:op3-two-port-frontier}, every
admission belongs to the positive support of the full obstacle at $\lambda$.
Thus the process terminates and the finite-band gate rule is safe.
Conservation gives $\lambda V\leq1$ on every reached face.

At termination, \cref{eq:op3-arm-terminal-residual} and active-face
equalities give $0\leq\bm e_v-\bm M\bm u^U\leq\epsappr\bm d$.
Consequently $\widehat{\bm\pi}=\bar\alpha\bm D\bm u^U$ has the required
ACL witness. The final reverse pass evaluates every eliminated equation
once and writes only positive output coordinates.

Each of $n_U-1$ admissions changes at most $r^2$ inverse entries. Sparse
inverse-vector products cost at most $O(rV)$ altogether; the sum of
candidate arm counts is $O(V)$. Dense inverse growth allocates $O(r^2)$
entries. Original row entries, cached degrees, edge dictionaries and stored
reverse records cost $O(V\log(2+V))$ using comparison dictionaries.
\Cref{lem:op3-scalar-flux-work} pays for every event and failed queue test.
These bounds include initialization and combine to
\cref{eq:op3-branch-core-work}. Only the inverse is quadratic in storage.
\end{proof}

\paragraph{Implementation and remaining lemma.}
\texttt{branch\_core\_flux.py} implements this construction with exact
fractions, scalar binary heaps and explicitly updated dense inverse entries.
Its Python dictionaries are an implementation convenience; comparison
dictionaries give the conservative word bound above. Full-graph solves,
inverse checks and residual checks occur only in the separate audit wrapper.
This implementation does not reproduce the full obstacle support: its output
is an exact restricted solution chosen to meet the ACL residual tolerance.

The remaining general-graph obligation is now concentrated in the retained
core: replace dense inverse changes and full scans of the growing core's
queues by an implicit response primitive that reports all due scalar
thresholds and certifies quietness. A fast named coordinate query or a fresh
near-linear core solve alone does not pay for that sequence. The polynomial
$r^2$ factor is an upper bound for this backend, not a lower bound for OP3.

\begin{proposition}[Dense-core cost on an easy tree family]
\label{prop:op3-dense-core-binary-tree}
On the complete binary tree of height $h\geq2$, take seed at the root,
$\alpha=1/1009$ and $\epsappr=\alpha/(10\cdot6^h)$.
Every ACL output has full support. The explicit inverse backend above
therefore performs $\Omega(n^3)$ word writes on this $n=2^{h+1}-1$ vertex
family, while its total number of flux publications is near-linear up to
logarithmic factors. Scanning all retained queues is separately charged.
This is a statement about the explicit retained-inverse representation.
\end{proposition}
\begin{proof}
The Neumann expansion of $\bm M^{-1}$ contains the unique source-to-vertex
walk of length at most $h$. Since all degrees are at most three and
$\gamma>1/2$, its contribution gives
\[
 \frac{\pi_i}{d_i}
 =\bar\alpha(\bm M^{-1})_{iv}
 \geq\frac{\bar\alpha}{3\cdot6^h}
 >\epsappr.
\]
ACL implies degree-normalized semantic error at most $\epsappr$, so no
coordinate can be zero in an ACL output. The retained set therefore has
$r=2^h-1$ vertices. At each successive retained birth, the previous inverse
has order $k=1,\ldots,r-1$; its top-left update changes every one of its
$k^2$ entries. Indeed the current restricted graph is connected, its
retained Schur inverse is entrywise positive, and the new coupling is
nonzero and nonnegative, so the update vector $\bm z$ is strictly positive.
The mandatory writes total
$\sum_{k=1}^{r-1}k^2=r(r-1)(2r-1)/6=\Omega(n^3)$.
The same reasoning forces changes to all explicitly stored old retained
values, separately from inverse maintenance. Balanced trees already have
short rooted paths; this is not a graph-access hardness example.
\end{proof}

\section{Peeling the live core while preserving physical-flux groups}
\label{sec:op3-live-core-peeling}

\paragraph{Status and scope.}
The following \statusproved{} construction is a proof draft awaiting
independent review. It strengthens \cref{thm:op3-branch-core-flux} by
allowing retained vertices to disappear after their side branches settle.
The work depends on the maximum simultaneous core size, including a
temporarily admitted vertex, rather than on the total number of original
branching vertices admitted. No graph-uniform bound on that maximum is
claimed. General OP3 remains \statusopen.

Use $\lambda=\epsappr/2$, $\bm M=\bm D-\gamma\bm A$, and
$\bm b=\bm e_v-\lambda\bm d$. Every strictly positive admission is
temporarily retained. The seed always remains retained. Other core vertices
are eliminated whenever their current reduced degree is at most two,
counting distinct retained and inactive neighbors after parallel couplings
have been coalesced. Denote the current retained set by $C$ and exposed
inactive set by $B$. The admitted original set $U$ includes both retained
and eliminated vertices.

\subsection{The reduced matrix and grouped physical signals}

Maintain the exact Schur matrix on $C\cup B$, retaining its core block,
core--boundary couplings, boundary diagonals and transformed loads.
Original inactive--inactive edges are revealed only on admission of an
endpoint. No fill edge between two inactive vertices is created.
For a pair $(c,j)\in C\times B$, collect a disjoint group of original
boundary edges whose total physical flux has the form
\begin{equation}
 F_{cj}=a_{cj}u_c+b_{cj},\qquad
 a_{cj}>0,\quad b_{cj}\leq0.
 \label{eq:op3-grouped-physical-flux}
\end{equation}
Store the original incidence count $k_{cj}$ and a lower publication
$\ell_{cj}$. At each inactive $j$ the groups partition all original edges
from $U$ into $j$, so its exact gate is
\[
 g_j=-\lambda d_j+\sum_c F_{cj}.
\]
The group slope is also the reduced coupling between $c$ and $j$.
Its intercept contributes to the transformed load of $j$.

\begin{lemma}[One boundary group per eligible core removal]
\label{lem:op3-live-core-removal}
If a nonseed core vertex $i$ has reduced degree at most two, it has at
least one core neighbor and at most one inactive neighbor. Eliminating
$i$ changes at most two remaining neighbors and moves at most one
physical-flux group. The representation
\cref{eq:op3-grouped-physical-flux}, its signs, and the absence of
inactive--inactive fill are preserved.
\end{lemma}
\begin{proof}
Every admitted vertex has an original active neighbor at admission, hence
the admitted graph is connected to the seed. Positive Schur elimination
preserves connectedness among the remaining retained vertices. Thus a
nonseed retained vertex has a retained neighbor, proving the degree claim.

Suppose the two neighbors are retained $c$ and inactive $j$. Write
\[
 \delta_i u_i=b_i+w_{ic}u_c+a_{ij}u_j,
 \qquad F_{ij}=a_{ij}u_i+b_{ij}.
\]
At the pinned boundary value $u_j=0$, elimination moves this physical group
to retained coordinate $c$ with
\begin{equation}
 a'_{cj}=a_{ij}w_{ic}/\delta_i,\qquad
 b'_{cj}=b_{ij}+a_{ij}b_i/\delta_i,
 \qquad \delta_j\leftarrow\delta_j-a_{ij}^2/\delta_i.
 \label{eq:op3-live-core-group-move}
\end{equation}
All nonseed transformed core loads remain negative: initially they are
$-\lambda d_i$, and eliminating another negative-load vertex adds a
negative multiple to them. Schur pivots remain positive. Thus the group
signs are preserved. Keep its original incidence count and lower
publication unchanged. If a group $(c,j)$ already exists, add the slopes,
intercepts, incidence counts and lower publications. The physical edge
sets being combined are disjoint.

For two retained neighbors, ordinary two-neighbor Schur complementation
changes their diagonals, loads and mutual coupling; there is no boundary
group to move. With one retained neighbor and no inactive neighbor, the
operation removes a settled leaf response. In all cases there are at most
two distinct neighbors, so no two inactive vertices receive a fill edge.
Save the removed scalar equation for recovery.
\end{proof}

When $j\in B$ is admitted, its current diagonal, transformed load
$-\lambda d_j+\sum_c b_{cj}$, and couplings $a_{cj}$ form its exact
bordered core equation. Retire all groups ending at $j$, then scan its
original row. A newly revealed edge to an inactive vertex starts a direct
group with $(a,b,k,\ell)=(\gamma,0,1,0)$. An original-edge dictionary
prevents repeating an old incidence. Queue the new core vertex for peeling;
after a removal, only its retained neighbors need another eligibility test.
Thus no global scan is needed to find eligible removals.

\subsection{Events, inverse restriction and recovery}

Use a scalar minimum heap at each retained coordinate, with group threshold
\begin{equation}
 \tau_{cj}=\frac{(5/4)\ell_{cj}+k_{cj}\epsappr/16-b_{cj}}{a_{cj}}.
 \label{eq:op3-grouped-flux-threshold}
\end{equation}
After updating the core, settle all these heaps. If $u_c>\tau_{cj}$,
publish the actual group flux and update the shared sum
$L_j=\sum_c\ell_{cj}$. Queue $j$ once this sum exceeds
$(11/20)\epsappr d_j$. Equality stays inactive.

Moving and coalescing groups preserves the published sum at their common
target. In particular, a moved group's lower publication is not reset to
zero. A published group remains a lower bound after every legal admission
by monotonicity of the restricted solution. A settled group satisfies
\[
 \ell_{cj}\leq F_{cj}\leq\tfrac54\ell_{cj}
                  +k_{cj}\epsappr/16.
\]
Summing over the disjoint physical incidences gives the same original
residual certificate as \cref{eq:op3-arm-terminal-residual}. Every queued
candidate stays legal, including through representation changes.

Maintain the retained inverse $\bm P=\bm H_C^{-1}$ and retained values.
An admission uses the bordered update in
\cref{eq:op3-core-rank-update}. Removing a retained vertex leaves the
restricted solution unchanged, and the inverse of its new Schur matrix is
exactly the principal submatrix of $\bm P$ on the remaining retained labels.
The implementation stores dense inverse rows keyed by original vertex
labels. Deleting one row and column costs $O(|C|)$ dictionary entries;
there is no global change of the labels used by old group or recovery records.

Each original positive vertex is either still retained or has exactly one
saved equation. Its coefficients name original remaining-neighbor labels,
including any inactive neighbor at that time. In reverse elimination order,
every later eliminated neighbor has already been recovered, and a neighbor
that never activated contributes zero. This recovers all original positive
coordinates once. A vertex's temporary core birth is not a second output
record.

\begin{theorem}[A bound in maximum live-core size]
\label{thm:op3-live-core-peeling}
Let $q$ be the largest simultaneous retained set, including temporary
admissions before peeling. Put $V=\vol(U)$, $n_U=|U|$, and
$L=1+\log_+(1/(\bar\alpha\epsappr))$ for the returned support $U$.
On every canonical input graph, the algorithm returns a nonnegative ACL
approximation with exact-real word work
\begin{equation}
 O\left((n_Uq^2+qV+VL)\log(2+V)\right)
 =\widetilde O((1+q^2)/\epsappr),
 \label{eq:op3-live-core-work}
\end{equation}
and working storage $O(V+q^2)$. Empty output costs $O(1)$.
Only positive-output adjacency rows are scanned, and degrees are queried
only at the seed or exposed neighbors. No topology, elimination order,
final support or bound on $q$ is supplied.
\end{theorem}
\begin{proof}
The reduced equations and \cref{lem:op3-live-core-removal} preserve the exact
restricted solution. Scalar publications certify positive gates for the
obstacle at $\lambda=\epsappr/2$. Inverse positivity therefore makes all
admissions safe and all old coordinates increase. At termination active
residuals equal $\lambda d_i$, while the summed group bands bound inactive
residuals by $(3/4)\epsappr d_i$. This yields the ACL witness, and
conservation gives $V\leq2/\epsappr$.

Each original row is scanned once. Each admitted vertex is removed from
the core at most once. A removal changes only two reduced neighbors and
moves at most one group. A direct edge starts one group and each move
starts at most one replacement group, possibly coalescing an existing
group. Thus the total number of group lifetimes is $O(V)$, including all
replacements. During a fixed lifetime, physical membership is fixed and
$0\leq F_{cj}\leq k_{cj}\gamma/\bar\alpha$. A first new publication
exceeds $k_{cj}\epsappr/16$, and successive publications grow by more
than $5/4$. Hence each lifetime has $O(L)$ publications, regardless of its
incidence count. Each retired entry is discarded once, including heaps
discarded with a removed core. Every unsuccessful heap test, core-queue
scan, ready-queue operation and peeling eligibility test is charged.

The number of core-queue visits is at most $n_Uq$. Each admission updates
at most $q^2$ inverse entries; sparse inverse-vector products cost at most
$O(qV)$ in total because candidate group counts sum to $O(V)$. Removing
inverse rows and columns costs $O(n_Uq)$ entries. Reduced-edge changes,
original-edge bookkeeping, group replacement and reverse equations cost
$O(V)$ records. Comparison dictionaries and scalar heaps give the displayed
logarithmic factor and a deterministic word bound. The audited Python code
uses dictionaries as a bookkeeping convenience. Dense inverse storage is
$O(q^2)$, and all group lifetimes, cached rows, scalar heaps and recovery
records require $O(V)$ words. No group operation copies its original
physical membership list; such lists appear only in the independent audit.
\end{proof}

\subsection{A source-valid separation between two admission orders}

\begin{proposition}[Peeling resolves the balanced-tree diagnostic under depth-first order]
\label{prop:op3-live-core-binary-order}
On the complete binary tree of height $h\geq2$, rooted at the seed, take
$\alpha=1/1009$ and $\epsappr=\alpha/(10\cdot6^h)$ as in
\cref{prop:op3-dense-core-binary-tree}. The last-in-first-out ready policy
has $q\leq h+1$ and therefore $\widetilde O(n)$ word work for this family.
The first-in-first-out policy has $q\geq2^h$ and incurs
$\Omega(n^3)$ explicit inverse writes before the first leaf is processed.
Both policies return valid ACL output on the same full support.
\end{proposition}
\begin{proof}
Every newly exposed boundary vertex becomes ready after the heaps settle.
To see this uniformly over connected root-containing faces, let $i$ be
the active parent of an inactive boundary vertex. Its depth is at most
$h-1$. A single walk term gives
$(\bm M_{UU}^{-1}\bm e_v)_i\geq2/6^h$.
Also $\bm M_{UU}\bm1\geq\bar\alpha\bm d_U$, so
$\bm M_{UU}^{-1}\bm d_U\leq\bm1/\bar\alpha$. Consequently
\[
 u_i^U\geq\frac2{6^h}-\frac\lambda{\bar\alpha}
 >\frac{39}{20\cdot6^h},\qquad
 r_j^U\geq\gamma u_i^U>\frac{39}{40\cdot6^h}
 >\tfrac34\epsappr d_j.
\]
The settled group bands then force
$L_j>(11/20)\epsappr d_j$. All old boundary vertices were already queued;
the newly queued vertices are exactly the children exposed by the latest
admission. Thus the two ready policies realize the usual breadth-first
and depth-first orders, with arbitrary sibling ties.

During depth-first exploration of a tree, all active vertices off the
current root path belong to completed subtrees. Their leaf responses have
been peeled, recursively, and their roots are eligible for further peeling.
No unfinished active subtree lies off the current search path. All retained
vertices therefore lie on that path; even including a temporary admission
there are at most $h+1$. Since $V=O(n)$ and $h=O(\log n)$, the work bound
follows from \cref{thm:op3-live-core-peeling}.

Under breadth-first order every internal vertex is admitted before any
leaf. Each nonseed internal vertex still has three distinct reduced
neighbors, so none has been peeled. The first leaf is temporarily added to
the $2^h-1$ internal retained vertices, proving the bound on $q$. The
preceding internal births alone force
$\sum_{k=1}^{2^h-2}k^2=\Omega(n^3)$ inverse writes. This is a difference
between two implementations on an easy graph family, not an OP3 lower bound.
\end{proof}

\paragraph{What remains open.}
Depth-first order alone does not bound $q$ on a general trace. A branch that
is quiet at one root value can become productive after another branch
changes the response. Many unfinished branches may therefore remain live.
The next target is a paid rule that either bounds the simultaneous core,
eliminates a higher-degree core vertex with a stronger grouped response,
or maintains its response and all due thresholds implicitly. Every change
of physical membership and every failed certificate must remain charged.

\begin{proposition}[Interior tree support blocks every first peeling step]
\label{prop:op3-live-core-interior-tree}
For any integer $h\geq1$, let the input be the complete binary tree of
height $H=3h+5$, seeded at its root, with
\[
 \alpha=\tfrac13,\qquad \gamma=\bar\alpha=\tfrac12,
 \qquad \epsappr=\frac1{12\cdot6^h},\qquad
 \lambda=\epsappr/2.
\]
Under either ready policy, and indeed under any legal ready policy in this
construction, no nonseed core vertex is peeled. Every valid ACL output
contains all vertices through depth $h$, so
$q=|U|\geq2^{h+1}-1$. The explicit inverse backend consequently requires
$\Omega(8^h)=\Omega((1/\epsappr)^{\log_6 8})$ word writes.
In particular, current-degree-two peeling plus depth-first order does not
give a uniform polylogarithmic live core or an OP3 work bound.
\end{proposition}
\begin{proof}
A source-to-vertex walk term, using original degrees at most three, gives
for every vertex at depth at most $h$
\[
 \pi_i/d_i=\bar\alpha(\bm M^{-1})_{iv}
 \geq\frac1{6\cdot6^h}=2\epsappr.
\]
The ACL semantic implication therefore forces positive output at every
such vertex.

At the other end of the tree, no original leaf can belong to the positive
obstacle support at $\lambda$. To verify this, write
\[
 \bm M^{-1}=(\bm I-\gamma\bm D^{-1}\bm A)^{-1}\bm D^{-1}.
\]
The transition matrix is row-stochastic. Its entries connecting a vertex
at depth $t$ to the seed vanish before time $t$ and are at most one
afterward. Since $d_v\bar\alpha=1$, the unregularized degree potential
at that vertex is at most $\gamma^t$. The obstacle potential is bounded
above by this unregularized solution. At the parent of an original leaf,
the possible incoming leaf flux is thus at most
\[
 \gamma^H=\frac1{32\cdot8^h}
 <\frac1{24\cdot6^h}=\lambda.
\]
A positive leaf would require $u_i=\gamma u_{p(i)}-\lambda>0$, a
contradiction. Safe admissions consequently never reach any original leaf.

Before any peeling step, every admitted nonseed vertex has its three
distinct original neighbors in the core or the exposed boundary. An
admission changes a neighbor's status but does not reduce that degree.
There can therefore be no first eligible removal. Induction gives no
removals at all, and $C=U$. Each successive admission changes every entry
of the previous positive retained inverse, giving
$\sum_{k=1}^{|U|-1}k^2=\Omega(|U|^3)$. The required depth-$h$ prefix
proves the stated bound. Since $\log_6 8>1$ and all relevant logarithms
grow only with $h$, this exceeds the OP3 scale for this backend.
\end{proof}

The last proposition is an implementation-specific negative control.
Teleportation is fixed in that family, so standard positive local-push
bounds already reach $O(1/\epsappr)$ there. It neither refutes OP3 nor
rules out a method choosing between response representations. It does show
that a small live-core assertion requires additional structure, and that
completed-subtree examples alone are insufficient evidence for it.

\section{Root-threshold transformations under legal tree admissions}
\label{sec:op3-root-threshold-groups}

\paragraph{Status and scope.}
The identities and counterexamples below are \statusproved{} proof drafts
awaiting independent review. All computations used to audit them are exact
reference computations, not an implemented fast local response producer.
They concern fixed $\lambda=\epsappr/2$ on canonical unweighted trees.
In particular, the counterexamples do not vary regularization or prescribe
external port values. General OP3 remains \statusopen.

\subsection{A resistance coordinate for the common transformations}

Root the input tree at the source $v$. Let $U$ be a connected strictly
positive restricted face, $E=U\setminus\{v\}$, and $B$ its exposed inactive
frontier. Keep the source potential as the formal parameter $t$ and set
\[
 \bm S=\bm M_{EE},\quad \bm G=\bm S^{-1},\quad
 \bm h_E=\bm G\gamma\bm A_{Ev},\quad
 \bm z_E=\bm G\bm b_E,
 \qquad h_v=1,\ z_v=0.
\]
Then $u_i(t)=h_it+z_i$ on $U$, with $h_i>0$ and $z_i\leq0$.
Extend $\bm G$ by a zero source row and column. Every $j\in B$ has a
unique parent $p(j)\in U$, and its exact gate and root threshold are
\begin{equation}
 g_j(t)=a_jt+\eta_j,\qquad
 a_j=\gamma h_{p(j)},\qquad
 \eta_j=-\lambda d_j+\gamma z_{p(j)},\qquad
 \tau_j=-\eta_j/a_j.
 \label{eq:op3-root-threshold-row}
\end{equation}
These affine formulas are algebraic continuations of the current face;
only the actual restricted source value is used to decide an admission.

Put $R_v=0$ and, for $i\in E$, define
\begin{equation}
 R_i=\sum_{(x,y)\in P(v,i)}\frac1{\gamma h_xh_y}.
 \label{eq:op3-harmonic-resistance}
\end{equation}
Here $P(v,i)$ is the path within the current positive tree.

\begin{lemma}[Green columns and common root-threshold transformations]
\label{lem:op3-root-resistance-transform}
For $p,x\in U$, let $\ell$ be their lowest common ancestor. Then
\begin{equation}
 G_{px}=h_ph_xR_\ell.
 \label{eq:op3-tree-green-resistance}
\end{equation}
Admit a frontier vertex $w$ with positive current gate, and put
\[
 x=p(w),\qquad
 \delta=d_w-\gamma^2G_{xx}>0,\qquad
 \kappa=\gamma^2h_x^2/\delta.
\]
For a surviving frontier vertex $j$, with
$\ell=\operatorname{LCA}(p(j),x)$, its new threshold satisfies
\begin{equation}
 \tau'_j=\frac{\tau_j+\kappa R_\ell\tau_w}
                    {1+\kappa R_\ell}
 =\tau_w+\frac{\tau_j-\tau_w}{1+\kappa R_\ell}.
 \label{eq:op3-root-threshold-transform}
\end{equation}
Thus every fixed-ancestor group shares one positive affine transformation
of its thresholds. Its internal order is preserved. Different groups need
not have the same transformation.
\end{lemma}
\begin{proof}
Let $\bm H=\operatorname{diag}(\bm h_E)$. The matrix
$\bm H\bm S\bm H$ is the grounded Laplacian of the current tree, with
edge conductances $\gamma h_ih_j$, including $h_v=1$ at the ground.
Indeed the harmonic equation for $h_i$ makes its diagonal equal the sum
of these incident conductances. Injecting one unit at $x$ sends all current
along the unique $x$--$v$ path. The resulting potential at $p$ is exactly
the resistance of the shared root path, namely $R_\ell$. Inverting the
diagonal transformation proves \cref{eq:op3-tree-green-resistance}; the
zero source row and column cover $p=v$ or $x=v$.

Eliminating the newly admitted $w$ at fixed source value gives
$u_w(t)=(a_wt+\eta_w)/\delta$. The change of an old potential is
$\gamma G_{p(j),x}u_w(t)$, and hence
\[
 (a'_j,\eta'_j)=(a_j,\eta_j)
       +\frac{\gamma^2G_{p(j),x}}\delta(a_w,\eta_w).
\]
Divide the new intercept by its positive slope and use
$a_j=\gamma h_{p(j)}$, $a_w=\gamma h_x$, and
\cref{eq:op3-tree-green-resistance}. This gives
\cref{eq:op3-root-threshold-transform}. Its slope is positive, so a
common group map preserves order. The pivot is a positive Schur pivot of
an SPD principal matrix. A legal admission also ensures that the actual
restricted solution remains strictly positive and monotone.
\end{proof}

\begin{corollary}[Updating the harmonic and resistance coordinates]
\label{cor:op3-root-coordinate-update}
For an old active vertex $p$, write
$\ell=\operatorname{LCA}(p,x)$ and $m_\ell=1+\kappa R_\ell$. Then
\begin{align}
 h'_p&=h_pm_\ell,&
 z'_p&=z_p+\gamma G_{px}\eta_w/\delta,\label{eq:op3-root-harmonic-update}\\
 R'_p&=\frac{R_\ell}{m_\ell}
          +\frac{R_p-R_\ell}{m_\ell^2}.
 \label{eq:op3-root-resistance-update}
\end{align}
On the path from the source to $x$, the latter map is
$R\mapsto R/(1+\kappa R)$.
\end{corollary}
\begin{proof}
The fixed-root Schur update and
\cref{eq:op3-tree-green-resistance} give the first two identities.
On an edge of the root--$x$ path, the updated resistance increment is
\[
 \frac{R_y-R_z}{(1+\kappa R_z)(1+\kappa R_y)}
 =\frac{R_y}{1+\kappa R_y}-\frac{R_z}{1+\kappa R_z}.
\]
These terms telescope as far as $\ell$. Below $\ell$ toward $p$, every
harmonic coordinate is multiplied by the same $m_\ell$, so every further
resistance is divided by $m_\ell^2$. This proves the formula.
\end{proof}

The corollary is an algebraic interface, not a data-structure theorem.
It does not pay for identifying the groups, changing their membership,
searching their transformed minima or maintaining their coordinates.
The existing ancestor-walk tree reporter in the companion AESP--CD note
already pays for exact threshold propagation. A faster implementation must
remove that path cost without assuming the groups are free to access.

\subsection{Source-valid failures of a stale global threshold order}

\begin{proposition}[A nine-vertex fixed-parameter ordering witness]
\label{prop:op3-tree-fixed-threshold-order}
Consider the tree with edges
\[
 01,\ 05,\ 12,\ 23,\ 34,\ 56,\ 58,\ 67,
\]
seed $0$, $\alpha=1/7$ and $\epsappr=1/20$, so
$\gamma=3/4$ and $\lambda=1/40$. FIFO exact-positive admission has prefix
$0,1,5,2,6,8$. Just before admitting $8$, the only positive frontier gate
is at $8$. The other two thresholds have order $\tau_3<\tau_7$.
After admitting $8$, their order reverses, vertex $3$ remains quiet, and
vertex $7$ becomes positive. Therefore recomputing only the old minimum
surviving threshold does not certify exact obstacle quietness.
\end{proposition}
\begin{proof}
Exact restricted elimination gives the following rows before and after
admitting $8$:
\[
\begin{array}{c|c|c|c|c}
 j&a_j&\eta_j&\tau_j&\tau'_j\\\hline
 3&27/220&-2/25&88/135&88/135\\
 7&9/116&-31/580&31/45&43/72
\end{array}
\]
The new row at $7$ is $(a'_7,\eta'_7)=(9/92,-43/736)$.
The actual source values are
$t=28523/46760$ and $t'=11143/17860$. Before admission,
\[
 g_3(t)=-961/187040,
 \quad g_7(t)=-229/37408,
 \quad g_8(t)=3523/46760>0.
\]
Afterward $g'_3(t')=-49/14288<0$, whereas
$g'_7(t')=373/142880>0$. All prefix admissions are strictly positive;
the exact FIFO audit verifies each gate on its actual preceding face.
The two group slopes in \cref{eq:op3-root-threshold-transform} are $1$
and $23/29$, corresponding to ancestors $0$ and $5$.
\end{proof}

\begin{proposition}[Minimum-threshold order can miss exact quietness]
\label{prop:op3-tree-minimum-policy-order}
Use the tree with edges
\[
 01,06,12,14,15,23,67,78,
\]
seed $0$, $\alpha=1/7$, and
$\lambda=1/31$ ($\epsappr=2/31$). The exact minimum-root-threshold policy,
with vertex-label ties, follows the prefix $0,6,1,7,4,5$.
Before $5$ activates it is the only positive frontier vertex, and the
minimum surviving threshold belongs to $8$. Afterward $8$ stays quiet
but $2$ is positive. This shortcut therefore fails exact quietness even on
a prefix generated by the minimum-threshold policy.
\end{proposition}
\begin{proof}
The five selected thresholds are respectively
$8/93,16/93,88/279,112/279,112/279$. Direct restricted elimination verifies
that each is the current minimum. Each is below the initial source value
$29/62$, so monotonicity also verifies all strict admissions. Just before
$5$, the two quiet thresholds are
$\tau_8=484/837<\tau_2=167/279$.
Afterward $\tau'_8=484/837$ and $\tau'_2=158/279$, reversing their order.
The actual source values are $1472/2573$ and $66704/115847$, giving
\[
 g'_8=-35/115847<0,\qquad g'_2=215/115847>0.
\]
No regularization parameter changes during this trace.
\end{proof}

\paragraph{Exact parameter-interval audit.}
The minimum policy can be explored without choosing a grid of tolerances.
For a fixed face write
$\bm u_U=\bm q_U-\lambda\bm y_U$, where
$\bm q_U=\bm M_{UU}^{-1}\bm e_v$ and
$\bm y_U=\bm M_{UU}^{-1}\bm d_U$.
A frontier vertex $j$ is positive exactly when
\[
 \lambda<c_j^U:=\frac{\gamma q_{p(j)}}{d_j+\gamma y_{p(j)}}.
\]
Since $\tau_j=\lambda(q_v/c_j^U-y_v)$, the next minimum-threshold label
is the maximizer of $c_j^U$, independently of $\lambda$ until stopping.
Intersect the strict upper bounds from previous admissions with the desired
quiet and positive inequalities to obtain an exact interval of valid fixed
parameters. The audit reconstructs each saved witness with separate dense
face solves and verifies every selected minimum.
Through eleven vertices, three teleportation choices and all seeds, it
checks 13,179 complete parameter-independent sequences and 122,388 symbolic
admissions. It finds 33 strict quiet-pair reorderings and 39 missed-positive
intervals. It finds no ACL miss under this policy in that finite family;
this is \textbf{Measured} evidence, not an ACL robustness theorem.

Both nine-vertex stopping witnesses concern exact stopping: their missed residuals are
still within the looser ACL tolerance. The next construction separates
that distinction explicitly.

\begin{proposition}[A stale minimum can also miss an ACL violation]
\label{prop:op3-tree-stale-minimum-acl}
Let a tree have the two arms $0$--$1$--$2$--$3$--$8$ and
$0$--$4$--$5$--$6$--$9$, the extra leaf edge $4$--$7$, four leaves
$10,11,12,13$ at $8$, and one leaf $14$ at $9$.
Set $\alpha=1/1009$ and $\epsappr=4/51$, so
$\gamma=504/505$ and $\lambda=2/51$.
The positive-admission prefix $0,1,2,3,4,5,6,7$ is legal. Before $7$
activates it is the only positive frontier vertex, and the old minimum of
$\{\tau_8,\tau_9\}$ is $\tau_8$. Afterward $8$ remains quiet but the
original physical residual at $9$ exceeds $\epsappr d_9$.
Thus exact validation of just that old minimum cannot certify ACL quietness.
\end{proposition}
\begin{proof}
An independent exact restricted solve gives
\[
 g'_8=-\frac{84478293610659282106}{9305780219492155967283}<0,
 \qquad
 g'_9=\frac{823081029939108794444}{9305780219492155967283}.
\]
Since $d_9=2$, the physical residual is $r'_9=g'_9+4/51$.
Its exact excess above $\epsappr d_9=8/51$ is
$93215914684822051912/9305780219492155967283>0$; the ratio is approximately
$1.0638583$. Before admission, the threshold difference is positive:
\[
 \tau_9-\tau_8
 =\frac{89389333405}{68556886272}
  -\frac{260536855325}{205670658816}>0.
\]
The preceding gates at $8,9$ are negative and the gate at $7$ is positive.
Every earlier gate on the stated prefix is positive as well; the durable
witness records each exact rational value. These claims can also be checked
without a dense solve by eliminating each three-vertex arm. If its first
diagonal is $a$ and its first load is $-\lambda b$, its last potential is
\[
 \frac{\gamma^3t-\lambda\{a(4+2\gamma)+\gamma^2(b-2)\}}
 {a(4-\gamma^2)-2\gamma^2}.
\]
Use $(a,b)=(2,2)$ on the left, $(3,3)$ on the right before admission,
and $(3-\gamma^2,3+\gamma)$ on the right afterward, then impose the source
row. This reproduces the displayed exact inequalities.
\end{proof}

The proposition rejects this specific stale-order shortcut on a legal
single-source trace. It is not a failure of the verified physical-flux
queues, and it is not a lower bound for local algorithms.

\subsection{Why one operation per ancestor group is still a real cost}

\begin{proposition}[A legal admission can have many distinct group maps]
\label{prop:op3-tree-many-threshold-groups}
For every $m\geq2$, there is a canonical tree and a forced legal backbone
trace with positive support volume $3m$ such that admission number $k$
has $k$ distinct transformations of surviving frontier groups, including
$k-1$ nonidentity transformations. Thus an implementation explicitly
applying a map to every such group performs $\Omega(m^2)$ group operations
on that trace. This demonstrates a growing number of distinct maps; it is
not a lower bound against an OP3 algorithm or a polylogarithmic
amortization in the accuracy.
\end{proposition}
\begin{proof}
Use the backbone $0$--$1$--$\cdots$--$m$, seeded at $0$. At every
$i<m$ attach a report vertex with $D-1$ private leaves, giving report
degree $D$. Choose
\[
 \alpha=1/1009,\quad \bar\alpha=1/505,\quad
 \lambda=\frac{\bar\alpha}{12\cdot6^m},\quad
 \epsappr=2\lambda,\quad D=12\cdot6^m\cdot505^2.
\]
The reports and their leaves remain inactive: the unregularized potential
is at most $1/\bar\alpha$, so incoming report flux is at most
$\gamma/\bar\alpha<\lambda D=1/\bar\alpha$.

Every positive backbone prefix has positive next gate. A single walk term
on a prefix gives source response at least $1/(3\cdot6^m)$ at each of its
vertices. As before, the load correction is at most
$\lambda/\bar\alpha=1/(12\cdot6^m)$. Thus its last potential is at least
$1/(4\cdot6^m)$. Since backbone degrees are at most three and
$\gamma>3\bar\alpha/2$, the next gate is strictly positive. The same bound
exceeds the finite-band readiness threshold
$(3/4)\epsappr d$, so the trace is also compatible with that interface.
There is exactly one possible next positive vertex. The resulting support
is the whole backbone and has volume $2+3(m-1)+1=3m$.

Before admitting $k$, the parent is $k-1$, and the surviving report at
$i<k$ has lowest common ancestor $i$. The resistances
$R_0<R_1<\cdots<R_{k-1}$ are strictly increasing, so the slopes
$1/(1+\kappa R_i)$ are distinct. Only $R_0=0$ gives the identity map.
Summing $k-1$ over $k=1,\ldots,m$ proves the stated explicit-group cost.
The chosen $1/\epsappr$ grows exponentially with $m$, so this example
must not be presented as work exceeding the OP3 target. Moreover its large
report degrees permit other quietness certificates that may avoid exact
threshold updates altogether. Nor does this family rule out an
$O(\vol(U)\log(1/\epsappr))$ charge for the explicit group operations.
\end{proof}

\paragraph{The narrower open interface.}
The current formulas suggest a rooted hierarchy supporting an insertion at
$x$, the LCA-dependent updates
\cref{eq:op3-root-threshold-transform,eq:op3-root-harmonic-update,eq:op3-root-resistance-update},
and a complete search for due original-residual events. Group discovery,
online hierarchy changes, failed queries and reconstruction must all be
charged. The formula for a known group is insufficient. The supplied-tree
persistent-curve solver does not supply local discovery, and the previous
ancestor iterator retains its radius factor. Ordinary independently updated
planar hulls also do not automatically implement the required changing
coordinates. This interface remains \statusopen.

\section[Ordered path clusters and projective reporters]{Ordered path clusters with persistent projective reporters}
\label{sec:op3-ordered-path-clusters}

\paragraph{Scope.}
The following \statusproved{} statements are proof drafts awaiting independent
review. They give an implemented algebraic backend for supplied rooted tree
clusters. The audit driver rebuilds supplied hierarchies; it does not implement
an online hierarchy or establish a graph-uniform OP3 result. We use exact real
word arithmetic, counting comparisons, unsuccessful searches, dictionary
operations, copies, retained records, graph access and final recovery.
Rational bit complexity and numerical stability are outside this model.
Put $\gamma=(1-\alpha)/(1+\alpha)$, $\lambda=\epsappr/2$,
$\boldsymbol M=\boldsymbol D-\gamma\boldsymbol A$ and
$\boldsymbol b=\boldsymbol e_v-\lambda\boldsymbol d$.
The current positive face $U$ is a connected rooted subtree containing $v$.

\subsection{A finite chart for both endpoint strata}

A two-port cluster has ordered ports $(L,R)$, with $L$ closer to the seed.
Eliminating its interior gives each owned boundary gate the form
\[
 g_j(\boldsymbol u)=a_j u_L+b_j u_R+c_j,
 \qquad a_j,b_j\geq0,\quad a_j+b_j>0,\quad c_j<0.
\]
The source is never an eliminated interior vertex. Consequently the interior
load is negative, and inverse positivity gives the asserted signs.
Use the homogeneous point and its finite chart
\begin{equation}
 (X_j,Y_j,W_j)=(a_j,c_j,a_j+b_j),\qquad
 (x_j,y_j)=(X_j/W_j,Y_j/W_j).
 \label{eq:op3-compact-projective-chart}
\end{equation}
Both $a_j=0$ and $b_j=0$ are finite: $0\leq x_j\leq1$ and $y_j<0$.
The normalized gate is $x_j u_L+(1-x_j)u_R+y_j$. Its maximum is attained
on the upper convex chain of these points. The normalization preserves signs,
so a positive maximizing row is an actual positive original gate, while a
nonpositive maximum certifies every owned row. This does not claim to preserve
the ordering of unnormalized gate magnitudes.

\begin{lemma}[Positive projective pullback; \statusproved]
\label{lem:op3-compact-projective-pullback}
Suppose old ports equal $\boldsymbol P$ times new ports plus
$\boldsymbol t$, where
$\boldsymbol P=\left(\begin{smallmatrix}a&b\\c&d\end{smallmatrix}\right)$
is entrywise nonnegative with positive determinant and
$\boldsymbol t=(t_L,t_R)^\top\leq0$.
The complete row pullback is
\begin{equation}
 \begin{pmatrix}X'\\Y'\\W'\end{pmatrix}=
 \begin{pmatrix}
 a-c&0&c\\ t_L-t_R&1&t_R\\a+b-c-d&0&c+d
 \end{pmatrix}
 \begin{pmatrix}X\\Y\\W\end{pmatrix}.
 \label{eq:op3-compact-projective-matrix}
\end{equation}
It has positive denominator on the whole compact chart, preserves increasing
$x$ order, and maps the old upper convex chain to the new upper convex chain.
\end{lemma}
\begin{proof}
The new coefficients are $aX+c(W-X)$ and $bX+d(W-X)$; their sum is
$W'>0$, including both endpoints since both row sums of $\boldsymbol P$
are positive. Differentiation of $x'=X'/W'$ gives the positive numerator
$ad-bc$. The homogeneous determinant in
\eqref{eq:op3-compact-projective-matrix} is also $ad-bc$.
Every orientation sign is therefore preserved after division by positive
homogeneous coordinates. At fixed $x$, the new $y$ coordinate increases
strictly with the old $y$ coordinate. Thus upper and lower sides cannot
exchange, including equal-$x$ endpoint strata. Negative offsets preserve
$Y'<0$. These facts establish the upper-chain assertion without scanning rows.
\end{proof}

\begin{lemma}[Implemented persistent chain; \statusproved]
\label{lem:op3-persistent-projective-chain}
Let $N$ bound the number of owned rows and put $\ell=\log(2+N)$.
An immutable AVL chain supports a whole-chain pullback in $O(1)$ work and one
copied root; the upper hull of two chains with separated $x$ intervals in
$O(\ell^3)$ work and $O(\ell^2)$ newly allocated nodes; and a gate-sign query
or positive one-port substitution in $O(\ell^2)$ work. Earlier versions remain
queryable. All allocated nodes, including subsequently discarded ones, are
charged to these bounds.
\end{lemma}
\begin{proof}
Store a constant-size homogeneous matrix tag at every node, along with its
point, two children, height and subtree size. A whole-chain map transforms the
root point and composes its tag. A point access carries the product of tags
along one root-to-leaf path in $O(\ell)$ work. Pushing a tag when a structural
operation needs its children copies at most two roots. Persistent AVL split,
join and rotations therefore have at most $O(\ell^2)$ copied nodes and work
under a conservative accounting. The separated-bridge searches from
\cref{lem:op3-avl-upper-hull} use $O(\ell^2)$ point accesses, including failed
comparisons. Homogeneous orientation and endpoint comparisons cost constant
word work per accessed point. After the bridge is found, persistent splits
retain the two surviving subchains and concatenate them. Equal endpoint
abscissae keep the higher point, with a fixed label tie rule; collinear
interior points may be omitted. This proves the stated conservative bounds.

For a one-port substitution
$(u_L,u_R)=(p,q)t+(\xi,\omega)$ with $p,q>0$ and $\xi,\omega\leq0$, a row threshold is
\[
 -\frac{Y+(\xi-\omega)X+\omega W}{(p-q)X+qW}.
\]
The denominator is positive on $[0,1]$. The fractional transformation
$x\mapsto x/((p-q)x+q)$ is increasing, and transforms the upper chain into
an upper chain. The corresponding ordinate, the negative threshold, is thus
unimodal along its ordered vertices. Binary adjacent-value comparisons find
its maximum with $O(\ell)$ point accesses. A flat adjacent maximum is a
supporting edge, so ties cannot conceal a later larger value. Ordinary gate
queries use the same argument with constant denominator. Nodes are immutable;
reclamation, if performed, can be charged once per allocated node. Keeping
all historical versions instead costs at most the total allocation count.
\end{proof}

\subsection{Compression, raking, and one-port reduction}

Cluster edges retain the original diagonal and load at each endpoint.
For a cluster, its matrix uses its own edges and the original diagonal and
load once per vertex. These matrices are positive definite M-matrices:
original degrees dominate the cluster degrees, with strict diagonal dominance
because $0<\gamma<1$. When two edge-disjoint clusters share a port, subtract
one copy of that port's original diagonal and load before eliminating it.
This convention avoids distributing a vertex's load over a changing number
of incident edges.

\begin{lemma}[Adjacent compression has separated chains; \statusproved]
\label{lem:op3-adjacent-chain-separation}
Let two clusters have ports $(L,M)$ and $(M,R)$, with $M\ne v$.
Write their Schur states as $(\boldsymbol H^l,\boldsymbol f^l)$ and
$(\boldsymbol H^r,\boldsymbol f^r)$. Set
\[
 \delta=H^l_{MM}+H^r_{MM}-d_M>0,\quad
 \beta=f^l_M+f^r_M-b_M<0,\quad
 w_L=-H^l_{LM}>0,\quad w_R=-H^r_{MR}>0.
\]
Then
\begin{equation}
 u_M=\eta u_L+\zeta u_R+\sigma,
 \qquad(\eta,\zeta,\sigma)=(w_L,w_R,\beta)/\delta.
 \label{eq:op3-adjacent-recovery}
\end{equation}
The left and right reporter pullbacks have respective matrices
$\left(\begin{smallmatrix}1&0\\\eta&\zeta\end{smallmatrix}\right)$ and
$\left(\begin{smallmatrix}\eta&\zeta\\0&1\end{smallmatrix}\right)$,
and offsets $(0,\sigma)$ and $(\sigma,0)$. Their new compact abscissae lie in
\[
 \text{right child: }[0,\eta/(\eta+\zeta)],\qquad
 \text{left child: }[\eta/(\eta+\zeta),1].
\]
Hence their exact union reporter uses one separated merge, with the right
child first. The new Schur matrix, load and recovery record need $O(1)$ words.
\end{lemma}
\begin{proof}
Schur elimination gives \eqref{eq:op3-adjacent-recovery}. Positivity follows
from the M-matrix property; $\beta<0$ follows because every eliminated vertex
is a nonseed with negative original load. Each displayed port matrix is
nonnegative with positive determinant, so
\cref{lem:op3-compact-projective-pullback} applies. Evaluate the monotone
abscissa maps at $0$ and $1$ to obtain the intervals. They meet only at the
shared-port endpoint. The new entries are
\[
 H_{LL}=H^l_{LL}-w_L\eta,\quad
 H_{RR}=H^r_{RR}-w_R\zeta,\quad
 H_{LR}=-w_L\zeta=-w_R\eta,
\]
with loads $f_L=f^l_L+w_L\sigma$ and $f_R=f^r_R+w_R\sigma$.
\end{proof}

\paragraph{Raking and forgetting.}
A one-port side cluster at $i$ is represented by its scalar Schur state and
its smallest boundary threshold $\theta$. Raking it into a path cluster adds
its diagonal and load, subtracting the original $d_i,b_i$ once. Conditional
on the unchanged path ports, the old path interior and reporter do not
change. All side gates reduce exactly to the single inequality $u_i>\theta$,
so its reporter contributes one endpoint point, merged at $x=0$ or $x=1$.
Raking two one-port clusters takes the minimum threshold and the same scalar
state sum. If there are no owned gates, omit the threshold.

To forget the farther port $R\ne v$ of a two-port cluster, substitute
\[
 u_R=(-H_{LR}/H_{RR})u_L+f_R/H_{RR}.
\]
Its multiplier is positive and offset negative. Apply the one-port query of
\cref{lem:op3-persistent-projective-chain}, update the scalar Schur state,
and save the affine recovery record. This proves closure of all these
operations with $O(\ell^3)$ work and $O(\ell^2)$ allocated hull nodes per
operation, independently of how many original rows the cluster represents.

\paragraph{Unique row ownership and recovery.}
Every inactive boundary vertex $j$ of a connected active subtree has exactly
one active parent $i$. At $i$, maintain the minimum of the local thresholds
$\lambda d_j/\gamma$ in a heap. All these gates have the same coefficient
of $u_i$, so this single minimum is an exact local certificate. Home the
resulting single row on $i$'s permanent parent edge; home the seed's row on
its first admitted edge. An admission changes the old parent's one home-edge
payload and creates the new vertex's payload. Copying that row to every
incident edge would be unnecessary and could cost quadratically on a star.
The singleton face is handled directly before the first edge exists.
Original adjacency entries and neighbor-degree queries are charged when a
vertex is admitted, with cached records and heap construction also charged.

The current hierarchy has a linear number of cluster records. A final
traversal uses the saved affine recovery equations, passes port values to
children, and writes each active vertex once. A dictionary verifies or
coalesces repeated shared ports; balanced dictionaries give $O(|U|\log(2+|U|))$
work, or a preassigned active index gives linear word work. Historical roots
need not be traversed. The returned PageRank vector is
$\bar\alpha\boldsymbol D\boldsymbol u$, with original residual
$\boldsymbol e_v-\boldsymbol M\boldsymbol u$.

\subsection{Precise remaining hierarchy contract}

\paragraph{Source.}
Alstrup et al.\ \cite{alstrup2005toptrees}, journal Section 2 and
Theorem 2.1 (pp.\ 245--248), maintain edge-induced clusters with at
most two boundary vertices. A link, cut or expose uses $O(\log n)$ joins and
splits, with logarithmic height and linear structural space. Section 6.1
(pp.\ 259--260) explains exposure; Section 6.2 (p.\ 260) handles arbitrary
degrees within the implementation. Link and cut reset external boundaries.
This source supplies a hierarchy theorem, not the PageRank summaries above.

\paragraph{Conditional.}
If an online hierarchy maintains the stated rooted cluster summaries with
$O(\log(2+|U|))$ changed records per admission and home-edge update, then the
backend yields total word work
\[
 O\bigl(\operatorname{cvol}(U)\log^4(2+\operatorname{cvol}(U))\bigr)
\]
and at most $O(\operatorname{cvol}(U)\log^3(2+\operatorname{cvol}(U)))$
retained state, even keeping every old application version. Temporary
query matrices are discarded as the search advances; their allocations
are charged to word work, not accumulated as persistent hull nodes. The root's one-port
Schur equation gives $u_v$ exactly, and its minimum threshold gives either a
strict positive original gate or a complete quiet certificate. Every strict
admission increases the positive face, whose volume is at most
$1/\lambda=2/\epsappr$. Thus this contract would give
$\widetilde O(1/\epsappr)$ local ACL work on arbitrary trees. This statement
charges failed searches, the final quiet query, local heaps, graph discovery,
all factor and hull records, dictionaries and output.

\paragraph{Integration obligations at the primitive stage.}
The code currently supplies and reparents static hierarchies. Before removing
the preceding condition, reconcile every join shape, singleton initialization,
source exposure during updates, changing base-edge payloads and dynamic
vertex allocation with the source interface. In particular, a temporary
cluster with an interior seed cannot use the negative-load argument. One
candidate is to mark such a summary unavailable until seed exposure splits
it; another is the source's composite-update interface. The unavailable-summary option is proved and fully integrated in
\cref{sec:op3-local-top-trees}; it is not silently assumed by this primitive. The source's internal high-degree representation is
structural and must not become extra PageRank vertices. No ambient graph
preprocessing is permitted.

\paragraph{Measured.}
The implemented immutable chain passes 516 independent materialized-chain
comparisons, 360 exact extreme queries, 360 one-port queries and 96 adjacent
separation checks, with 130 retained versions rechecked. The cluster audit
uses every nonisomorphic tree through seven vertices, every seed and three
parameter pairs, following exact positive faces. It checks independent dense
Schur complements, every raw boundary inequality, recovered coordinates,
three path parenthesizations and both orders of one-port rakes. Five implicit
unweighted path/private-star families through 128 active vertices keep all
boundary rows on the upper chain; their inactive adjacency is never scanned.
These are correctness and representation audits, not an online hierarchy
benchmark. Exact counts and source hashes are in
\texttt{PROJECTIVE\_HULL\_ROPE\_AUDIT.json} and
\texttt{PATH\_CLUSTER\_REPORTER\_AUDIT.json}.

\section[Complete local tree construction]{A complete local tree construction from the top-tree interface}
\label{sec:op3-local-top-trees}

\paragraph{Status and scope.}
This section discharges the hierarchy condition in
\cref{sec:op3-ordered-path-clusters} for arbitrary finite unweighted trees.
It is \statusproved{} as a proof draft awaiting independent review. The
published balancing algorithm is a \statussource{} import; the application
callbacks and exact audits are implemented here. The exhaustive audit driver
is not a fast online balancing implementation. No conclusion for arbitrary
cyclic graphs follows from this tree result.
Use the exact-real word model and the local degree/adjacency access model of
\cref{sec:scope}. The physical seed is $v$,
$\boldsymbol M=\boldsymbol D-\gamma\boldsymbol A$, and
$\boldsymbol b=\boldsymbol e_v-\lambda\boldsymbol d$, with
$\lambda=\epsappr/2$. Every checkpoint exposes $v$ as the single external
boundary of the current active tree. Depths and parent edges refer to that
permanently rooted mathematical tree, not to the changing cluster hierarchy.

\begin{lemma}[Rooted ports and unavailable summaries; \statusproved]
\label{lem:op3-source-interior-adapter}
During a leaf insertion and subsequent exposure of $v$, a cluster can store
its port set, a Boolean indicating whether it contains $v$, and either its
exact summary from \cref{sec:op3-ordered-path-clusters} or an unavailable
marker. Mark it unavailable exactly when $v$ is in its interior.
A valid cluster has one or two ports, its rootward vertex is a port, and its
two ports, if present, are comparable in the tree rooted at $v$.
A valid parent never needs an unavailable child. After exposing $v$, every
cluster is valid. Marking, splitting and propagating these flags cost
constant work per structural operation.
\end{lemma}
\begin{proof}
The maintained forest has one component containing edges, namely the active
tree containing $v$, and at most one newly allocated isolated vertex.
A zero-port edge cluster must be the whole component, hence contains $v$
in its interior and is unavailable. Consider any other cluster $C$ whose
rootward vertex is $r$. If $v\notin C$, the edge from $r$ toward $v$ leaves
$C$, making $r$ a boundary. If $v\in C$ and $C$ is valid, then $r=v$ is a
boundary by validity. Thus the rootward vertex is always a port. With at
most two ports, they are comparable: otherwise their least common ancestor,
or the exit from it toward $v$, would give a third boundary.

A vertex interior to a child cannot be a boundary of its parent: every edge
incident to it already belongs to the child, and it is not externally
exposed. In particular, an unavailable child forces its parent to be
unavailable. Conversely, a valid parent has only valid children and may use
their algebraic summaries. These observations apply between updates.
The source discipline in Section 2.1 of \cite{alstrup2005toptrees} first
splits every cluster whose edges or boundary set will change. Thus an old
interior vertex cannot silently become a new boundary in a retained cluster.
After exposure, every cluster containing $v$ has it as a boundary, so every
unavailable cluster has been split or its base edge recreated.
Containment is the Boolean OR of the children's flags; no membership scan
is needed. An unavailable node retains its structural child pointers so a
later split needs no algebraic reconstruction.
\end{proof}

\begin{lemma}[Complete callback closure; \statusproved]
\label{lem:op3-top-tree-callback-closure}
With the convention of \cref{lem:op3-source-interior-adapter}, every base
creation and legal top-tree join is implemented by a constant number of the
edge, compress, rake and forget operations of
\cref{sec:op3-ordered-path-clusters}. It costs
$O(\log^3(2+N))$ exact word work and allocates
$O(\log^2(2+N))$ hull nodes, where $N$ bounds the number of owned rows.
No source load is eliminated in a valid summary.
\end{lemma}
\begin{proof}
At a base edge, sort its endpoints by stored depth. With two ports, use the
edge summary; with one port, retain its rootward endpoint and forget its
nonseed farther endpoint. A source-interior base edge is unavailable and
requires no algebra.

In a join, the two child clusters share exactly one vertex, which is a port
of each child. Every other child port remains a parent port: an edge leaving
the child there cannot belong to the other child, since their only common
vertex is the shared port. There are therefore exactly the following cases.
\begin{enumerate}
\item Two two-port children produce a two-port parent and eliminate their
shared nonseed middle port. Their depth ordering identifies left and right,
so \cref{lem:op3-adjacent-chain-separation} gives a separated compression.
\item One one-port and one two-port child produce either a two-port parent,
by raking, or a one-port parent, by raking and forgetting the shared port.
In the latter case the shared port is farther from $v$; otherwise the
parent's rootward vertex would be interior, contrary to
\cref{lem:op3-source-interior-adapter}.
\item Two one-port children share their sole port. If it remains a parent
port, rake their scalar states and take the minimum threshold. Otherwise
the parent has no ports and is unavailable.
\end{enumerate}
An unavailable parent is constructed as a structural record, independently
of its children's algebraic status. For every valid parent all children
are valid by \cref{lem:op3-source-interior-adapter}. The list exhausts the
source's five join shapes. At a valid elimination the eliminated vertex is
not $v$, and the eliminated interior is source-free. Its negative load
invariant and positive Schur pivots therefore hold. The primitive bounds in
\cref{lem:op3-persistent-projective-chain} finish the accounting.
\end{proof}

\begin{lemma}[One edge payload and online allocation; \statusproved]
\label{lem:op3-top-tree-payload}
A changed owned boundary row on one base edge can be maintained in
$O(\log^4(2+n))$ word work and $O(\log^3(2+n))$ allocated hull nodes in a
height-bounded top tree on $n$ active vertices. The local application needs
no preprocessing indexed by the ambient graph size.
\end{lemma}
\begin{proof}
Use the parent/child and edge-leaf pointers in the source representation,
and walk from the changed leaf to the current hierarchy root. Recompute the
base summary and every ancestor from its current children, preserving its
existing port set and structural join shape. Even unavailable ancestors
receive fresh child pointers. The walk has logarithmic length by Theorem 2.1 (journal p.\ 247)
of \cite{alstrup2005toptrees}; the callback bound proves the claim. This
operation changes application data only, so it does not rebalance the
structural hierarchy. Later joins, including joins that undo a previous
exposure, always derive fresh summaries from current children. They do not
reactivate stale numerical snapshots from before the payload change.

The underlying forest contains only admitted vertices and one newly created
isolated vertex. Source vertex and edge identifiers can therefore be local
indices allocated on demand. Store a comparison dictionary from discovered
original labels to local records, and keep each active incidence list as
local edge records. Resizable arrays can instead be allocated by doubling,
charging each copy to the growing number of exposed words. The source's
internal treatment of large degrees concerns these active structural
incidences; it neither accesses inactive adjacency nor adds vertices to the
PageRank matrix. Initialization of a singleton's empty top tree needs
constant state. The dictionary, cache, allocation and occasional copies all
cost at most a logarithmic factor per discovered word.

If a concrete source implementation requests a fixed vertex capacity,
start with constant capacity and double it when exhausted. Rebuild only
the currently active structural tree, using its cached edges and current
home rows. A size-$m$ rebuild uses $O(m\log^4(2+m))$ application work by
inserting its edges in parent-before-child order through the same source
interface, and allocates
$O(m\log^3(2+m))$ persistent hull state. The geometric sum of the doubled
capacities is $O(n)$, so these bounds remain $O(n\log^4(2+n))$ work and
$O(n\log^3(2+n))$ retained state in total. This alternative never reserves
capacity indexed by the original graph's unknown size.
\end{proof}

\begin{theorem}[Local exact obstacle and ACL solves on trees; \statusproved]
\label{thm:op3-local-top-tree}
On an arbitrary finite, simple, connected, unweighted tree with a supplied
seed vertex and local degree/adjacency queries, there exists a deterministic
algorithm returning the exact obstacle solution for
$\boldsymbol M=\boldsymbol D-\gamma\boldsymbol A$ and
$\boldsymbol b=\boldsymbol e_v-\lambda\boldsymbol d$.
If $U$ is its returned positive support, its total exact-real word work is
\begin{equation}
 O\!\left(1+\operatorname{cvol}(U)
     \log^4(2+\operatorname{cvol}(U))\right),
 \label{eq:op3-local-top-tree-work}
\end{equation}
and its space, even retaining every historical application version, is
$O(1+\operatorname{cvol}(U)\log^3(2+\operatorname{cvol}(U)))$ words.
For $0<\epsappr<1$ and $\lambda=\epsappr/2$, the vector
$\widehat{\boldsymbol\pi}=\bar\alpha\boldsymbol D\boldsymbol u$
has original source residual
$0\leq\boldsymbol e_v-\boldsymbol M\boldsymbol u
 \leq\lambda\boldsymbol d\leq\epsappr\boldsymbol d$.
Consequently its fully charged local ACL work is
$\widetilde O(1/\epsappr)$, with no polynomial dependence on $1/\alpha$.
\end{theorem}
\begin{proof}
First query $d_v$. If $\lambda d_v\geq1$, return zero, which is the exact
obstacle solution. Otherwise the singleton face has
$u_v=(1-\lambda d_v)/d_v>0$. Scan its row, query and cache neighbor degrees,
and build its heap of thresholds $\lambda d_j/\gamma$. The singleton's
minimum threshold gives the exact next gate or a quiet stop directly.
When the first child is admitted, create the first active edge and home
both endpoint rows there.

For every later positive face, retain the source-exposed top tree with its
complete root summary. The scalar Schur equation gives $u_v$ in constant
work. The root's minimum threshold either names an actual strict positive
original boundary gate or certifies all gates nonpositive, since one-port
normalizations preserve signs. A named candidate has exactly one active
parent: two active parents, together with the path between them in the
connected active set, would create a cycle in the original tree.

Admit the named candidate $j$. Remove it from its parent's local heap and
refresh that parent's one home-edge payload. Scan $j$'s original row once;
its sole previously active neighbor is its parent. Cache its other
neighbors and their degree replies, build its own heap and one home row,
and allocate its new active vertex and edge records. Invoke one source
\texttt{link}, then \texttt{expose} on $v$. Original degrees and loads in all
existing numerical records remain fixed. Changes to active incidence and
port sets are handled by the source's split/create/join discipline, while
\cref{lem:op3-source-interior-adapter} handles temporary source-interior
clusters. Only the heap minimum is homed on an edge; the other original
candidate thresholds remain in the local heap for later admissions.

Every admitted gate is evaluated at the exact restricted solution. A
positive Schur complement correction gives $u_j>0$ and increases all old
active coordinates, so the face stays positive and grows monotonically.
Summing the original residual over the entire graph gives
\[
 \boldsymbol1^{\mathsf T}
 (\boldsymbol e_v-\boldsymbol M\boldsymbol u)
 =1-\bar\alpha\boldsymbol d^{\mathsf T}\boldsymbol u.
\]
Every inactive residual is nonnegative, and each active residual equals
$\lambda d_i$. Therefore $\lambda\vol(U)<1$ at every positive checkpoint.
There are finitely many strict admissions. At a quiet checkpoint the active
equations and all boundary inequalities are exactly the obstacle KKT
conditions; other inactive rows have zero residual. Uniqueness gives the
exact obstacle solution and the displayed ACL certificate. This proof
identity does not require an algorithmic scan of the whole graph.

Let $V=\operatorname{cvol}(U)$ for the final nonempty support. Scanning every
admitted row once costs at most $V$. The number of discovered candidate
records and degree replies is at most $V$; dictionary accesses, heap
construction/deletion and all unsuccessful heap or final quiet checks cost
$O(V\log(2+V))$. Each admission uses one home-edge update, one link and one
expose, hence $O(\log(2+V))$ recomputed cluster records by the source theorem
and \cref{lem:op3-top-tree-payload}. Every callback costs at most
$O(\log^3(2+V))$ word work and allocates $O(\log^2(2+V))$ hull nodes.
Thus all comparisons, searches, factor updates, copies and historical
records satisfy \eqref{eq:op3-local-top-tree-work} and the claimed space
bound. Structural source state is linear in active size and already paid.
The optional capacity-doubling rebuilds in
\cref{lem:op3-top-tree-payload} satisfy these same total bounds.

Finally solve the root's scalar equation once and traverse only the current
cluster hierarchy with its saved affine recovery records. Each represented
edge or helper record is visited a constant number of times. Active indices
allow each vertex to be emitted once; comparison dictionaries would add only
an already allowed logarithm. Multiply by $\bar\alpha d_i$ for the sparse
PageRank output. Neither the past hierarchies nor old full solution vectors
are traversed. Since $V\leq2\vol(U)<2/\lambda$, the stated ACL bound follows.
\end{proof}

\paragraph{An exact RPPR consequence.}
In the nonzero regime $0<\rho<1/d_v$, taking $\lambda=\rho$ and returning
$\boldsymbol x=\bar\alpha\boldsymbol D^{1/2}\boldsymbol u$ gives the exact
shared RPPR minimizer on trees, by the mapping in \cref{sec:mapping}.
Its word work is $\widetilde O(1/\rho)$ and it has zero objective gap in this
exact-real model. This structural statement does not convert ACL accuracy
into the semantic PPR tolerance of OP1 and does not assert numerical stability.

\paragraph{Exact implementation evidence.}
\texttt{top\_tree\_callback\_adapter.py} implements the application layer.
Its final audit covers 650 canonical positive faces: every nonisomorphic
core tree through eight vertices, every seed, and two teleportation choices.
Each core vertex has one inactive private-star report. The original degrees
and tiny positive $\lambda$ give a legal positive core face with every report
quiet; the stars need not be materialized. All 167,384 legal binary cluster
joins under exposed and unexposed seed boundaries are checked, covering all
five join shapes. There are 26,130 independent Schur oracles, 168,294 valid
summary comparisons, 104,794 owned rows, 270,792 exact two-port queries,
78,030 one-port checks and 833,718 recovered coordinates. The audit also
checks 650 complete exposure transitions and 1,296 intermediate edge-payload
refreshes, then rechecks the old roots. Payload-only perturbations test the
algebraic interface; they are not represented as legal graph admissions.

The driver enumerates clusters, constructs supplied hierarchies and performs
dense reference solves. Its 240.431-second runtime does not measure the
source hierarchy's online bound. The asymptotic result imports that source
algorithm and proves every additional application cost above; it does not
infer complexity from finite timing. Audit counts, graph/parameter rules,
source and backend hashes are in \texttt{TOP\_TREE\_CALLBACK\_AUDIT.json}.

\section[Uniform shift and geometric root]{A uniform shift separates the source from the geometric root}
\label{sec:op3-shifted-tree-clusters}

\paragraph{Status.}
The algebra below is \statusproved{} as a draft awaiting independent review.
Its application to an online unicyclic solver is completed in
\cref{sec:op3-local-unicyclic}, also as a proof draft.
The exact word model is essential: subtracting two large nearly equal real
numbers here is not a finite-precision stability guarantee.

\begin{lemma}[Negative loads with an arbitrary physical seed; \statusproved]
\label{lem:op3-uniform-negative-shift}
Set $C=1/\bar\alpha$ and write $u_i=w_i+C$ on a represented active tree.
For an edge-induced cluster $B$, let $d_B(i)$ be its own edge degree, while
$d_i$ is the original graph degree. Its shifted load is
\begin{equation}
 b'_B(i)=\mathbf1_{i=v}-\lambda d_i-C\bigl(d_i-\gamma d_B(i)\bigr)<0.
 \label{eq:op3-shifted-cluster-load}
\end{equation}
It is represented locally by the fixed vertex load
$\beta_i=\mathbf1_{i=v}-\lambda d_i-Cd_i$ and a contribution $\gamma C$
to each endpoint of every cluster edge. Joins subtract one duplicate
$\beta_i$ at a shared port. No vertex-wide load redistribution is required.
This remains true when the physical seed is an eliminated interior vertex.
\end{lemma}
\begin{proof}
The identity is $\boldsymbol b'_B=\boldsymbol b_B-C\boldsymbol M_B\boldsymbol1$.
Because $d_B(i)\leq d_i$ and $C(1-\gamma)=1$,
\[
 b'_B(i)\leq\mathbf1_{i=v}-\lambda d_i-d_i<0.
\]
Each incident cluster edge supplies exactly one $\gamma C$ term. Combining
two clusters counts the fixed vertex load twice at their shared port, so
subtracting one $\beta_i$ gives the desired union load. All interior loads
are negative, independently of the source's location. Positive inverse
entries then preserve the negative Schur-load and affine-offset invariants.
The diagonal and off-diagonal matrix entries are unchanged.
\end{proof}

\begin{lemma}[Reporter equivalence after shifting; \statusproved]
\label{lem:op3-shifted-reporters}
The compact homogeneous reporter and its persistent chain operations remain
correct with arbitrary real intercepts $Y$, including $Y>0$.
In the shifted coordinates, a local gate at parent $i$ has threshold
$\lambda d_j/\gamma-C$. If its original normalized two-port point is $(x,y)$,
its shifted point is $(x,y+C)$. Thus every upper-chain membership, ordering
and strict gate sign is preserved by this common vertical translation.
One-port thresholds decrease by $C$.
\end{lemma}
\begin{proof}
For a row $g=a u_L+b u_R+c$, substitution gives
$g=a w_L+b w_R+c+C(a+b)$, proving the point and threshold identities.
The upper-chain and bridge proofs use positive $W=a+b$, orientation signs
and nonnegative port coefficients, not the sign of $Y=c$. The fractional
one-port query uses a positive denominator and a positive coefficient of
$Y$ in its numerator, independently of the numerator's value. Therefore
negative thresholds and arbitrary intercepts do not change the algorithms
or their work bounds. In particular, \cref{lem:op3-uniform-negative-shift}
keeps the existing nonpositive-offset pullback contract intact, even when
the source has been eliminated. If a conditional interior response was
$u_i=z_i+\boldsymbol a_i^{\mathsf T}\boldsymbol u_P$, it becomes
\[
 w_i=z_i+C(\boldsymbol a_i^{\mathsf T}\boldsymbol1-1)
       +\boldsymbol a_i^{\mathsf T}\boldsymbol w_P.
\]
This also gives an independent identity for checking every recovered value.
\end{proof}

\paragraph{A separate geometric anchor.}
The unavailable-summary argument can now refer to any chosen geometric
anchor $r$, retained as an external vertex, rather than to the physical seed
$v$. Root the represented tree at $r$ so ordered path ports remain
comparable. The physical seed may lie in a valid cluster's interior because
\eqref{eq:op3-shifted-cluster-load} is still negative. The only modification
in \texttt{ShiftedBackend} is to add $\gamma C$ to the two base-edge loads;
all compress/rake/forget operations use their existing negative-load paths.
Changing a geometric root by rebuilding the currently exposed tree costs
work and must be charged explicitly; the shift does not make repeated
rebuilding free.

\paragraph{One extra edge: an exact algebraic reduction.}
Suppose an active connected graph is represented by a spanning tree plus
one extra edge $\{r,t\}$. Use the original graph degrees, retain $r,t$ as
ports, and let $(\boldsymbol H,\boldsymbol h')$ be the shifted tree's root
Schur state. The full active matrix is restored by solving
\begin{equation}
 \left[\boldsymbol H-\gamma
   \begin{pmatrix}0&1\\1&0\end{pmatrix}\right]
 \begin{pmatrix}w_r\\w_t\end{pmatrix}
 =\boldsymbol h'+\gamma C\begin{pmatrix}1\\1\end{pmatrix}.
 \label{eq:op3-one-edge-restoration}
\end{equation}
The two added load terms are required by the per-edge rule in
\cref{lem:op3-uniform-negative-shift}. The corrected matrix is positive
definite, being a Schur complement of the full original principal matrix.
All interior conditional responses and reporters are unchanged by the extra
edge, so recovery followed by adding $C$ gives the exact full face solution.
\Cref{sec:op3-local-unicyclic} supplies the separately charged end-to-end
local construction using this identity.

\paragraph{A source-valid pitfall.}
On the triangle with physical seed $0$, $\alpha=1/3$ and $\lambda=1/20$,
the legal active prefix $(0,1)$ has values $(7/15,1/15)$. The gate for vertex
$2$ with both active parents is $1/6>0$, while its gate using only parent
$1$ is $-1/15$. After admission, the full triangle solution is
$(1/2,1/10,1/10)$. However, deleting edge $\{0,2\}$ and solving just the
path's unconstrained system with the original degrees gives
$(13/28,2/35,-1/28)$.
A cycle-restoration algorithm must therefore retain the exact signed
spanning-tree response. It cannot replace that response by the spanning
tree's obstacle solution or assume the uncorrected spanning-tree face is
positive. Equation \eqref{eq:op3-one-edge-restoration} restores the full
positive solution exactly. A separate four-vertex triangle-with-leaf case
checks the physical seed strictly inside a one-port side cluster while the
cycle endpoints are retained.

\paragraph{Measured.}
\texttt{shifted\_tree\_clusters.py} checks every core tree through seven
vertices, every physical seed, geometric anchors at the seed and at a
farthest core vertex, and two teleportation choices. Each vertex has an
inactive private-star report; small and large report degrees give unfinished
positive faces and both signs of reporter intercept. The final sweep covers
564 faces, 33,550 conditional shift identities, 9,496 independent Schur
oracles and 28,606 valid summary comparisons. It includes 1,972 valid clusters
with the physical seed interior and 14,012 positive-intercept rows. Both
cycle-restoration witnesses use independent dense full-matrix references.
These audits establish implementation evidence for the identities. The
separate local cyclic work proof is in \cref{sec:op3-local-unicyclic}.
\texttt{UNICYCLE\_TOP\_TREE\_PROBE.md} preserves the original obligations
and records their resolution.

\section[Local unicyclic construction]{Local unicyclic discovery with two fixed cycle ports}
\label{sec:op3-local-unicyclic}

\paragraph{Status and model.}
The construction below is \statusproved{} as a proof draft awaiting
independent review. It extends \cref{thm:op3-local-top-tree} to arbitrary
finite simple connected unweighted unicyclic graphs, with an arbitrary
physical seed and arbitrarily large tree attachments. Neither the cycle
nor any attachment is supplied. The exact-real word and local graph-access
models are unchanged. The height-bounded top-tree balancing algorithm is
the same explicit \statussource{} import. The local state machine,
application algebra and named queries are implemented, while the audit
rebuilds reference hierarchies and counts those rebuilds separately.
The result does not settle arbitrary-graph OP3.

\begin{lemma}[Named coordinates through current cluster records; \statusproved]
\label{lem:op3-cluster-point-query}
Given a current height-$h$ cluster hierarchy, solved values at its at most
two root ports, and a stored incident-edge locator for an active vertex
$i$, its exact coordinate can be queried in $O(1+h)$ word operations.
The query uses a constant-size affine row and current parent pointers; it
does not recover or scan any other active coordinate. It therefore costs
$O(\log(2+n))$ on the height-bounded source hierarchy.
\end{lemma}
\begin{proof}
Start with the basis row selecting $i$ from its incident base edge's two
endpoints. A one-port base edge has one forget record, which immediately
pulls this row back to the structural leaf's port. Each structural join
introduces at most two application records. Compression recovers the shared
port as $w_m=\eta w_L+\zeta w_R+\sigma$; forgetting recovers
$w_b=\eta w_a+\sigma$; raking preserves each child port among the parent
ports. Hence a child's at most two port coordinates are affine functions
of the parent's at most two ports, computable through these at most two
new records in constant work. Pull the current affine row through that map
and follow the structural parent pointer. At the root, evaluate the row at
the solved port values and add $C$ if physical $u_i$ is requested.

The published source stores edge-leaf and structural parent pointers
(Section 2.4 of \cite{alstrup2005toptrees}). Store one current incident home
edge per vertex as in \cref{lem:op3-top-tree-payload}. An endpoint lookup,
including the root's first edge, therefore needs no search. Current child
to parent affine maps can be computed from the saved helper records or
cached with each freshly created parent in constant additional work.
Changed payload ancestors recompute these maps from current children.
No old numerical parent version is reactivated after an exposure update.
Equivalently, the expanded application hierarchy has at most twice the
structural height, and the query follows its current parent links.
The audit's whole-hierarchy index construction is a reference convenience;
it is not needed by this source-pointer implementation.
\end{proof}

\begin{lemma}[The unique exceptional boundary candidate; \statusproved]
\label{lem:op3-unicyclic-exception}
Let $U$ be a connected active set in a graph of cycle rank at most one.
If $G[U]$ is a tree, every inactive boundary vertex has one active neighbor,
except possibly one vertex having two active neighbors. If $G[U]$ is
already unicyclic, every inactive boundary vertex has exactly one active
neighbor. The exceptional vertex, if present, can be discovered from
cached incidences while scanning newly admitted rows.
\end{lemma}
\begin{proof}
Adjoin any collection $B$ of inactive boundary vertices and only their edges
to $U$. The resulting connected subgraph has cycle rank
\[
 |E(U)|-|U|+1+\sum_{j\in B}\bigl(|N(j)\cap U|-1\bigr)\leq1.
\]
Every summand is nonnegative. The assertions follow immediately, including
the impossibility of a third active neighbor. Since every active row has
already been scanned, a boundary record contains all its current active
parents. A newly scanned incidence either creates its first parent or
reveals the sole possible second parent. Thus no inactive row or unknown
cycle information is required.
\end{proof}

\paragraph{State before closure.}
Use the uniform shift of \cref{lem:op3-uniform-negative-shift} from the
start. Root the active spanning tree at the physical seed $v$, exposing
only $v$ at each checkpoint. Every ordinary candidate $j$ has one active
parent $i$, so its heap key is $\kappa_j=\lambda d_j/\gamma-C$.
Home only the minimum live key from each active vertex's heap. This gives
the root's exact ordinary sign certificate by
\cref{lem:op3-shifted-reporters}.

When a second active parent is discovered for $j$, remove its entry from
the first parent's ordinary heap and exclude it from the new parent's
heap. Refresh the first parent's home edge. Store the exceptional triple
$(j,p,q)$ separately. Lazy heap deletion is also valid: a live-entry test
checks the cached parent tuple and active status; every stale entry is
removed at most once. At every checkpoint, including unsuccessful checks,
use \cref{lem:op3-cluster-point-query} twice to evaluate
\[
 g_j=\gamma(u_p+u_q)-\lambda d_j.
\]
Either an ordinary reporter or this exceptional gate may supply the next
strict positive admission. Stop before closure only if both are quiet.
The exceptional record is created at most once: it persists until its
vertex is admitted, which closes the sole cycle, or until termination.

\paragraph{The single closure and subsequent growth.}
Admitting an ordinary candidate adds one leaf to the stored spanning tree.
Admitting the exceptional candidate $j$ adds two edges to the active graph.
Keep either new edge in the spanning tree and designate the other as the
extra edge $\{c,j\}$. Re-root and rebuild the cached active spanning tree
once at geometric anchor $c$, re-homing its current minimum rows. Insert
its edges through the source interface in parent-before-child order, then
expose $c,j$. The uniform negative-load shift allows the physical seed to
be interior, and \eqref{eq:op3-one-edge-restoration} gives the exact full
active solution at the two retained ports. The underlying uncorrected
spanning-tree response remains a signed affine object; no positivity or
obstacle truncation is imposed on that response.

By \cref{lem:op3-unicyclic-exception}, every subsequent inactive candidate
has one active parent. Its admission adds a leaf to the stored spanning
tree and changes its parent's minimum row. Expose the same two fixed cycle
ports after each link. The root hull now gives a complete sign certificate
for all remaining ordinary boundary candidates, using the corrected root
values. Link operations may temporarily reset external vertices: apply the
unavailable-anchor convention of \cref{lem:op3-source-interior-adapter}
and make no gate query until the required exposure has been restored.
The source splits all clusters whose boundary set changes before an update,
so both one-port and two-port exposures preserve the callback contract.

\begin{theorem}[Local exact obstacle and ACL solves on unicyclic graphs; \statusproved]
\label{thm:op3-local-unicyclic}
For a finite simple connected unweighted graph with at least two vertices
and cycle rank at most one, the preceding deterministic local algorithm
returns the exact obstacle solution for
$\boldsymbol M=\boldsymbol D-\gamma\boldsymbol A$ and
$\boldsymbol b=\boldsymbol e_v-\lambda\boldsymbol d$.
For $0<\alpha<1$, $\lambda>0$, and returned positive support $U$, its total
exact-real word work is
\[
 O\!\left(1+\operatorname{cvol}(U)
       \log^4(2+\operatorname{cvol}(U))\right).
\]
Its space, even retaining all historical application versions, is
$O(1+\operatorname{cvol}(U)\log^3(2+\operatorname{cvol}(U)))$ words.
For $\lambda=\epsappr/2$ and $0<\epsappr<1$, it returns
$\widehat{\boldsymbol\pi}=\bar\alpha\boldsymbol D\boldsymbol u$ with
original residual between zero and $\lambda\boldsymbol d$, in fully
charged local work $\widetilde O(1/\epsappr)$.
\end{theorem}
\begin{proof}
If $\lambda d_v\geq1$, return zero after the seed degree query. Otherwise
start at the positive singleton. Exact Schur states and reporters,
\cref{lem:op3-unicyclic-exception}, and the separately checked exceptional
gate ensure that every admission has a strictly positive original gate.
The usual positive inverse block update makes the new coordinate positive
and increases the old coordinates. Thus every checkpoint is a positive
face, all admissions lie in the final obstacle support, and termination
occurs after at most $|U|$ admissions. At termination every active equation
is exact and every inactive gate is nonpositive. These are the obstacle KKT
conditions, proving uniqueness and correctness. Final recovery visits only
the current application records and adds $C$ once per emitted coordinate.
The original residual is nonnegative, equals $\lambda d_i$ on $U$, and is
at most $\lambda d_i$ outside $U$.

Write $V=\operatorname{cvol}(U)$. Each admitted original row is scanned once,
for at most $V$ adjacency entries. Distinct degree replies, discovered-label
records, active-membership tests and cached boundary incidences require
$O(1+V)$ logical operations; comparison dictionaries add a logarithm.
There is one heap insertion per first boundary incidence and at most one
later deletion per entry, with $O(\log(2+V))$ work per heap operation.
Each admission changes at most two old home payloads and one new payload;
the exceptional extra old-parent refresh occurs only once. By
\cref{lem:op3-top-tree-payload}, these changes and the source's link/expose
callbacks cost $O((1+V)\log^4(2+V))$ work. At most two named-coordinate
queries per checkpoint add $O(|U|\log(2+|U|))$ work. Constant-size root
solves and the final failed ordinary and exceptional queries are included.

Closure re-roots only the cached active spanning tree, at most once.
Traversing its edges, assigning new depths and homes, and rebuilding through
the source interface costs $O(|U|\log^4(2+|U|))$, with
$O(|U|\log^3(2+|U|))$ newly allocated persistent state. This is charged to
the final $V$, even when the remaining graph is enormous. The capacity-
doubling alternative in \cref{lem:op3-top-tree-payload} adds the same bounds;
no array is initialized using the ambient vertex count. All historical
application records satisfy the stated total allocation bound, while each
current hierarchy has only linear structural state and constant-size helper
records. Temporary arithmetic matrices and query rows have constant size;
their work is charged and they need not be retained after use. The final
recovery and dictionary output cost $O((1+V)\log(2+V))$.

For nonempty $U$, the source-mass identity gives $\lambda\operatorname{vol}(U)<1$.
The positive-degree graph implies $|U|\leq\operatorname{vol}(U)$, so
$V=O(1/\lambda)$. Substituting $\lambda=\epsappr/2$ proves the ACL bound.
\end{proof}

\paragraph{RPPR consequence.}
With $\lambda=\rho$, the same construction gives the exact regularized
solution $\boldsymbol x=\bar\alpha\boldsymbol D^{1/2}\boldsymbol u$,
with objective gap zero and $\widetilde O(1/\rho)$ work in the nonzero
regime $0<\rho<1/d_v$. This removes both the cycle-seed assumption and the
polynomial attachment-size factor in the earlier bounded-attachment result.
It remains a structural exact-word theorem, not a numerical implementation
or an arbitrary-graph OP3 theorem.

\paragraph{Measured and implementation boundary.}
\texttt{cluster\_point\_queries.py} checks 58,816 exact named coordinates
on all core trees through six vertices, including arbitrary root assignments,
one-port base edges, all recovery helper shapes and retained payload versions.

\texttt{local\_unicyclic\_clusters.py} uses actual strict local admissions on
all 78 connected tree and unicyclic atlas graphs through seven vertices,
every seed, four parameter pairs and two admission/extra-edge policies.
Independent dense references check every face, every original boundary
gate, cycle correction, final obstacle solution and ACL residual. There are
3,888 atlas executions and sixteen larger explicit attachment executions,
checking 16,371 faces and 26,043 original gates. They include 319 negative
uncorrected spanning-tree faces. Four finite implicit private-star examples
have 9, 17, 15 and 16 active vertices, respectively; their original row
scans inspect only 27, 51, 45 and 48 entries despite enormous ambient size.
Inactive hub rows are never scanned. Exact totals and source hashes are
recorded in \texttt{LOCAL\_UNICYCLIC\_CLUSTER\_AUDIT.json}.
The audit's rebuilt hierarchies and full validation recoveries are explicitly
separate reference work. The complexity theorem imports the verified online
source hierarchy, rather than claiming that these reference rebuilds are
fast online updates.

\section[Locally exposed cycle rank]{A coarse system for the locally exposed cycles}
\label{sec:op3-local-cycle-rank}

\paragraph{Status and parameters.}
This section gives a \statusproved{} construction as a proof draft awaiting
independent review. It applies to arbitrary finite simple connected
unweighted graphs in the same local degree/adjacency and exact-real word
models. Its work depends on the cycle ranks of the final active support
and of the locally revealed graph. These dependencies remain explicit;
the result does not resolve graph-uniform OP3. The published height-bounded
top-tree algorithm remains an explicit \statussource{} import, and the
implemented audit rebuilds application hierarchies as reference work.

\begin{lemma}[Revealed cycle rank and exceptional incidences; \statusproved]
\label{lem:op3-revealed-cycle-rank}
For a nonempty connected active set $U$, let
\[
 r(U)=|E(G[U])|-|U|+1,\qquad
 a_j=|N(j)\cap U|\quad(j\in\partial U),
\]
and define the revealed cycle rank
\begin{equation}
 q(U)=r(U)+\sum_{j\in\partial U}(a_j-1).
 \label{eq:op3-revealed-cycle-rank}
\end{equation}
It is the cycle rank of the connected graph on $U\cup\partial U$ with
exactly the original edges incident to $U$. Both $r(U)$ and $q(U)$ are
nondecreasing under connected admissions, and
$r(U)\leq q(U)\leq\operatorname{vol}(U)$.
There are at most $q(U)-r(U)$ exceptional candidates with $a_j\geq2$,
and their total number of active parents is at most $2(q(U)-r(U))$.
All these quantities are maintained from scanned incidences only.
\end{lemma}
\begin{proof}
The revealed graph has $|U|+|\partial U|$ vertices and
$|E(G[U])|+\sum_j a_j$ edges, proving the identity. Admitting a candidate
with $a$ active parents increases $r$ by $a-1$ and removes the same amount
from the boundary sum. Each newly scanned incidence to a previously
discovered boundary vertex increases $q$ by one; an incidence introducing
a new boundary vertex changes its edge and vertex counts equally.
Thus both ranks are nondecreasing. Each exceptional summand $a_j-1$ is
at least one, and $a_j\leq2(a_j-1)$ for $a_j\geq2$. Finally the revealed
edge count is at most the number of scanned entries $\operatorname{vol}(U)$.
No edge between two still-inactive vertices is queried or included.
\end{proof}

\begin{lemma}[Two-port partition with permanent row homes; \statusproved]
\label{lem:op3-marked-tree-partition}
Let $T$ be an insertion spanning tree of $G[U]$, permanently rooted at $v$,
and let $R$ be the $r=r(U)$ remaining active edges. If $r\geq1$, there is
a retained set $P$ of at most $4r$ active vertices and a partition of
$E(T)$ into $|P|-1$ connected tree pieces, each with exactly two ports in
$P$. Distinct pieces have disjoint interiors and share only retained ports.
The total number of vertex copies over all pieces is
$|U|-1+(|P|-1)$.
The partition, original-label maps and permanent row-home component map
can be constructed in $O((|U|+r)\log(2+\operatorname{cvol}(U)))$ word
work from cached active tree edges, using comparison dictionaries.
An ordinary one-parent admission subsequently extends just one piece;
no partition change or global re-homing is required.
\end{lemma}
\begin{proof}
Mark $v$ and all endpoints of edges in $R$, giving a set $S$ of at most
$2r+1$ vertices. Let $K$ be the minimal subtree of $T$ connecting $S$.
Add every degree-at-least-three vertex of $K$ to $S$ to obtain $P$.
Every leaf of $K$ is marked, so the number of unmarked branching vertices
is at most $|S|-2$. Consequently $|P|\leq2|S|-2\leq4r$.
The maximal paths in $K$ between vertices of $P$ form a coarse tree with
$|P|-1$ edges. Since $v\in P$ and $T$ is rooted at $v$, the endpoints of
each such path are comparable; choose its rootward endpoint as its anchor.

Assign all unmarked branches attached to an interior path vertex to that
path's piece. At a retained vertex other than $v$, assign all its off-$K$
branches to its rootward path piece. At $v$, choose an incident piece for
any off-$K$ branches. The pieces are connected trees whose only shared
vertices are their two retained endpoints. Their edge sets partition $T$.
If there are $c$ pieces, summing the tree identity
$|V(B)|=|E(B)|+1$ gives total copy count $|U|-1+c$.

Home vertex $i\ne v$ on its permanent insertion-parent edge, and home $v$
on the first admitted edge. A row stays on that physical tree edge through
every repartition. A retained nonroot vertex therefore homes its row in
the rootward piece. At $v$, if its home edge is off $K$, assign that branch
to the selected home piece; otherwise its path already determines the
home piece. Each ordinary minimum row appears exactly once. Other copies
of a shared retained vertex contribute no duplicate reporter row.

To build $K$, process the cached insertion order backward, marking whether
each subtree contains a marked vertex. A tree edge is in $K$ exactly when
its farther subtree is marked. Count these incidences, identify $P$, and
traverse the resulting nonbranching paths once. Process the insertion order
forward to assign off-$K$ edges to their parent's home piece. This reads
$O(|U|+r)$ records and writes the same order of component, copy and home
records; dictionaries and deterministic sorting add the stated logarithm.
There is no ambient adjacency scan or reserved ambient vertex array.

If a new candidate has one active parent, it adds a leaf to $T$ and does
not change $R$, $S$, $K$ or $P$. Assign its new edge to its parent's home
piece. That piece contains the parent, even when the parent is shared.
The new vertex is unmarked, so both ports and every other piece remain
unchanged. Its own permanent home is its new insertion edge.
\end{proof}

\paragraph{The tree regime.}
When $r=0$, use one piece, the whole active tree, with its sole port $v$.
The ordinary reporter is the scalar tree reporter. There may still be many
exceptional inactive candidates, so the single-exception argument of the
unicyclic theorem must not be used here.

\begin{lemma}[Shared-port assembly and cycle restoration; \statusproved]
\label{lem:op3-coarse-cycle-assembly}
For every piece $B$, maintain its shifted Schur state
$(\boldsymbol H_B,\boldsymbol h'_B)$ at its one or two ports, using
original full degrees and
$\beta_i=\mathbf1_{i=v}-\lambda d_i-Cd_i$, where $C=1/\bar\alpha$.
Embed and sum these states on the retained coordinates $P$. If $i\in P$
appears in $m_i$ pieces, subtract $(m_i-1)d_i$ from its summed diagonal
and $(m_i-1)\beta_i$ from its summed load. For every extra edge
$\{i,j\}\in R$, subtract $\gamma$ from the two off-diagonal entries and
add $\gamma C$ to both endpoint loads. The resulting system is the exact
Schur system of the full active matrix and shifted load on $P$.
It is positive definite. Its solution supplies the exact port values for
every component reporter and named-coordinate query.
\end{lemma}
\begin{proof}
Before interior elimination, each piece contributes one $d_i$ and one
$\beta_i$ at every represented vertex, plus the unique contributions from
its own tree edges. The interiors are disjoint and have no edges to other
interiors, while all extra-edge endpoints are retained. Eliminating the
interiors can therefore be done independently within pieces. At shared
ports, the stated diagonal and load subtractions remove exactly the
duplicated fixed vertex terms. The extra-edge update is the per-edge
shifted rule from \cref{lem:op3-uniform-negative-shift}. Thus the assembled
system is precisely the full principal matrix's Schur complement, with its
correct shifted right-hand side. Positive definiteness follows from that
of $\boldsymbol M_{UU}$. In particular, an uncorrected piece or spanning-tree
response need not be positive in the physical coordinates; no obstacle
truncation is applied to these affine response objects.
\end{proof}

\paragraph{Local update and query algorithm.}
Cache each boundary candidate's active-parent list. A one-parent candidate
belongs to its parent's ordinary heap with key $\lambda d_j/\gamma-C$.
On acquiring a second parent, remove its ordinary row, refresh the old
parent's home edge, and put it in an exceptional dictionary. Later parents
only extend its incidence record; the candidate stays outside all ordinary
heaps. At every checkpoint, solve the coarse system and query every piece's
ordinary reporter at its resulting port values. Separately evaluate each
exceptional original gate
\[
 g_j=\gamma\sum_{i\in N(j)\cap U}u_i-\lambda d_j
\]
using \cref{lem:op3-cluster-point-query} in each named parent's home piece.
All quiet queries are included. Either type of strict positive gate can
supply the next admission; stop only when all queried signs are nonpositive.

An admission with $a\geq2$ active parents keeps one new edge in $T$ and
adds the other $a-1$ edges to $R$. Rebuild the marked partition and the
current component hierarchies from cached active edges. This occurs at
most $r(U_{\rm final})$ times. A one-parent admission instead uses one
source leaf insertion in the parent's home piece and the needed old home
payload updates. The source restores each affected piece's fixed one or
two external ports before any new gate query. Separate local structural
copies of shared ports are glued only by the coarse equation; they are not
additional vertices of the PageRank problem.

\begin{theorem}[A local bound in active and revealed cycle rank; \statusproved]
\label{thm:op3-local-cycle-rank}
The preceding deterministic algorithm returns the exact obstacle solution
on an arbitrary finite simple connected unweighted graph with at least two
vertices, for $0<\alpha<1$ and $\lambda>0$.
For a nonempty returned positive support $U$, put
$V=\operatorname{cvol}(U)$, $r=r(U)$, $q=q(U)$ and $L=\log(2+V)$.
Its total exact-real word work is
\begin{equation}
 O\!\left(1+V\bigl[(1+r)L^4+(1+r)^3+qL^2\bigr]\right),
 \label{eq:op3-cycle-rank-work}
\end{equation}
and its space, including all historical application versions, is
\begin{equation}
 O\!\left(1+(1+r)VL^3+(1+r)^2\right).
 \label{eq:op3-cycle-rank-space}
\end{equation}
An empty support is certified in constant work by $\lambda d_v\geq1$.
With $\lambda=\epsappr/2$ and $0<\epsappr<1$, it gives original ACL residual between zero and
$\lambda\boldsymbol d$ and fully charged local work
$\widetilde O((1+r^3+q)/\epsappr)$.
With $\lambda=\rho$, it gives the exact RPPR solution, of objective gap zero,
in $\widetilde O((1+r^3+q)/\rho)$ work in the nonzero regime.
\end{theorem}
\begin{proof}
The singleton and zero cases are as in \cref{thm:op3-local-unicyclic}.
By \cref{lem:op3-coarse-cycle-assembly}, every checkpoint has exact full
active-face coordinates at all retained ports. The component reporters
cover each ordinary heap minimum exactly once, while the exceptional
queries use every actual active parent. Thus every admission has a strict
positive original gate and the block update of \cref{thm:block-schur}
keeps every admitted coordinate positive and increases the old ones.
Termination gives the exact obstacle KKT conditions. All admitted rows
belong to the unique final positive support, giving at most $|U|$ iterations
and a single final recovery. The residual and mass arguments from the
unicyclic theorem apply unchanged.

There are $O(1+V)$ original row entries, degree replies, discovered-label
records and incidence insertions. Ordinary heap insertions and lazy
deletions have the same total count. A new second-parent incidence causes
at most one extra old home-payload refresh, so there are $O(1+V)$ such
refreshes in total, even if a single admitted row creates many exceptional
candidates. Source leaf insertions, exposure restoration and height-bounded
payload walks cost $O((1+V)L^4)$ work and allocate $O((1+V)L^3)$ persistent
application state. Component copies use local identifiers; dictionaries,
capacity changes and their data copies are charged as in
\cref{lem:op3-top-tree-payload}.

By \cref{lem:op3-marked-tree-partition}, a cycle-birth rebuild visits
$O(|U|+r)$ cached words and creates pieces of total size $O(|U|+r)$.
Initialize their source hierarchies in parent-before-child order. Since
$r\leq V$, one such rebuild costs $O(VL^4)$ work and $O(VL^3)$ newly
allocated application state. At most $r$ rebuilds give the first terms in
\eqref{eq:op3-cycle-rank-work} and \eqref{eq:op3-cycle-rank-space}.
Per-component capacity doubling between rebuilds has the same geometric-sum
bound. Permanent original home edges and depths do not change, but their
component map and structural leaf locators are rebuilt and charged.

Each coarse system has dimension $O(1+r)$. Zeroing its dense matrix, copying
it for elimination, assembling piece entries, correcting duplicates and
restoring extra edges cost $O((1+r)^2+(1+r)L)$ word work. A deterministic
dense positive-definite solve costs $O((1+r)^3)$ work and $O((1+r)^2)$
temporary words. There is one solve per checkpoint, including termination.
Querying all $O(1+r)$ component reporters costs $O((1+r)L^2)$ per checkpoint.
By \cref{lem:op3-revealed-cycle-rank}, all exceptional parent incidences
number at most $2q$, so their named queries and dictionary work cost
$O((1+q)L^2)$ per checkpoint, including every failed gate check and the
candidate list search. These bounds are absorbed in
\eqref{eq:op3-cycle-rank-work}, since $|U|\leq V$.
If incidence lists are copied as tuples, as in the exact driver, their total
copy work is at most $O((1+q)V)$, also covered by that bound; old temporary
tuple copies need not be retained.

Current component copies and their original-label maps use $O(V+r)$ words.
Final recovery traverses each current component once, checks equal values
at shared ports, adds $C$ and emits each original coordinate once. Its work
is $O((V+r)L)$. Historical root pointers and their solved one/two-port
values cost $O((1+r)|U|)$ words and are included in the allocated-state
bound. Temporary coarse matrices and query arithmetic are paid work but
need not be saved as historical application versions. Finally
$\lambda\operatorname{vol}(U)<1$ implies $V=O(1/\lambda)$, proving both
accuracy-namespace consequences.
\end{proof}

\paragraph{Relation to OP3.}
The support ranks are output parameters, not a supplied decomposition or an
ambient preprocessing budget. For fixed revealed cycle rank this reaches
OP3's $\widetilde O(1/\epsappr)$ ACL scale, while the general bound retains
the explicit rank factors. Separately, the sufficient structural condition
$1+r^3+q=O(\alpha^{-1/2})$ puts the exact regularized solve at OP2's
conjectured $\widetilde O(1/(\rho\sqrt\alpha))$ scale. These are distinct
accuracy namespaces; no semantic PPR conversion is inferred here.
Large active cycle rank leaves a cubic coarse-system factor, and many quiet
exceptional candidates leave repeated named-query work. Neither dependency
is proved necessary. Removing or separately paying them is \statusopen{}.

\paragraph{Measured implementation boundary.}
\texttt{local\_cycle\_rank\_clusters.py} implements local incidence discovery,
permanent home edges, cycle-birth partition changes, shared-port correction,
coarse solves, complete exceptional gates and one final output. Independent
dense Schur references check the assembled matrix and shifted load on every
audited positive face, in addition to full obstacle and original-residual
comparisons. The audit exercises multiple simultaneous exceptional candidates
and single admissions that create several cycles. Rebuilt component
hierarchies, their full parent indices and dense reference recovery are
separately counted audit work. The asymptotic theorem uses the same verified
published online hierarchy as the earlier tree and unicyclic constructions.
The sweep covers all 995 connected atlas graphs through seven vertices,
every seed, four parameter pairs and two policies: 54,240 executions.
Twenty larger explicit cases and five finite implicit private-star cases
extend the families. In total, 217,410 positive faces, 168,368 independently
assembled coarse Schur systems and 483,413 original gates are checked.
The atlas alone has 33,074 admissions creating multiple cycles and 166,405
exceptional gate checks, of which 3,441 are quiet. Retained component roots
are rechecked 198,001 times. The five implicit examples inspect only
58, 51, 48, 62 and 70 original entries, respectively, and never scan an
inactive hub row. Full parameters, counts and source hashes are in
\texttt{LOCAL\_CYCLE\_RANK\_CLUSTER\_AUDIT.json}.

\section[Selected Green entries]{Selected Green entries from the same cluster records}
\label{sec:op3-cluster-green}

\paragraph{Status.}
The query identities below are \statusproved{} as a draft awaiting
independent review and are implemented in an exact audit. Replacing the
coarse solve of \cref{thm:op3-local-cycle-rank} by a complete maintained
inverse algorithm remains \statusopen{}. A block inverse formula alone
does not pay for changing retained vertices, selecting covariance entries
or updating every stored matrix word.

\begin{lemma}[Conditional variance and selected covariance; \statusproved]
\label{lem:op3-cluster-conditional-green}
Let a valid tree-piece hierarchy have height $h$, root ports $P_B$ and
interior $I_B$. For a named vertex $i$, one incident-edge parent walk returns
its conditional affine response
\[
 w_i=z_i+\boldsymbol a_i^{\mathsf T}\boldsymbol w_{P_B}
\]
and conditional variance
$g_{ii}=(\boldsymbol M_{I_BI_B}^{-1})_{ii}$, extended by zero when $i$
is a port, in $O(1+h)$ word work and constant live arithmetic state.
Optionally retain $O(1+h)$ scalar noise coefficients from that walk.
For two vertices in the same piece, their retained traces give
$g_{ij}=(\boldsymbol M_{I_BI_B}^{-1})_{ij}$, again extended by zero at
ports, in $O((1+h)\log(2+h))$ comparison-dictionary work.
Thus the height-bounded source hierarchy gives logarithmic named variance
work and at most $O(\log^2(2+n))$ selected covariance work.
\end{lemma}
\begin{proof}
At a compress or forget operation, let $m$ be the eliminated scalar and
$\delta_m>0$ its pivot. Completing its square factors the conditional
quadratic into
\[
 \tfrac12\delta_m
 \bigl(w_m-\boldsymbol\eta_m^{\mathsf T}\boldsymbol w_{\rm parent}
              -\sigma_m\bigr)^2
\]
plus the parent Schur quadratic. Raking joins disjoint interiors and
introduces no eliminated scalar. Applying this identity recursively gives
a triangular factorization with one independent scalar residual
$\xi_m$ of variance $1/\delta_m$ per eliminated interior vertex. This is
an algebraic covariance interpretation of the inverse factorization;
the algorithm samples no random variables.

Follow the affine pullback of \cref{lem:op3-cluster-point-query}. Whenever
the current row has coefficient $c_m$ on an eliminated port, add
$c_m^2/\delta_m$ to a variance accumulator and, optionally, save
$(m,c_m,1/\delta_m)$ keyed by its current elimination record. The rest of
the row pulls back to the parent ports as before. Previously collected
residual coefficients do not change, since the factorization makes those
residuals independent of the remaining coordinates. At the root,
\[
 w_i=z_i+\boldsymbol a_i^{\mathsf T}\boldsymbol w_{P_B}
          +\sum_m c_{i,m}\xi_m,
 \qquad
 g_{ij}=\sum_m\frac{c_{i,m}c_{j,m}}{\delta_m}.
\]
Only eliminated records on the query's ancestor path can occur. Each
structural callback has at most two helper records, so there are $O(1+h)$
terms. A compress pivot is the sum of the two children's shared diagonal
entries minus the original degree; a forget pivot is the child's eliminated
diagonal entry. Both are read or computed in constant work from existing
records. Matching two saved traces in a comparison dictionary gives the
claimed covariance cost. A root port has its basis response and an empty
noise trace, as required. Different pieces have conditionally independent
interiors given all coarse ports; their conditional cross term is zero by
the partition, not merely because two arbitrary record identifiers differ.
\end{proof}

\begin{lemma}[Port promotion and an original-vertex border; \statusproved]
\label{lem:op3-green-promotion-border}
Suppose the full current active inverse restricted to retained ports is
$\boldsymbol J=(\boldsymbol M_{UU}^{-1})_{PP}$.
Extend each piece response row $\boldsymbol a_i$ by zero to all of $P$.
For old active vertices $i,j$, its full Green entry is
\begin{equation}
 (\boldsymbol M_{UU}^{-1})_{ij}
   =g_{ij}+\boldsymbol a_i^{\mathsf T}\boldsymbol J\boldsymbol a_j,
 \label{eq:op3-conditional-full-green}
\end{equation}
where $g_{ij}=0$ across different piece interiors and at retained ports.
In particular, promoting old interior vertices to retained coordinates
augments $\boldsymbol J$ with these entries without solving a new full
active system.

If a new candidate $j$ has active-parent indicator $\boldsymbol s$ on $U$,
put $\boldsymbol g=\boldsymbol M_{UU}^{-1}\boldsymbol s$ and
$\delta=d_j-\gamma^2\boldsymbol s^{\mathsf T}\boldsymbol g$.
Then $\delta\geq\bar\alpha d_j>0$ and the enlarged inverse is
\begin{equation}
 \begin{pmatrix}
  \boldsymbol M_{UU}^{-1}+\gamma^2\boldsymbol g\boldsymbol g^{\mathsf T}/\delta
       &\gamma\boldsymbol g/\delta\\
  \gamma\boldsymbol g^{\mathsf T}/\delta&1/\delta
 \end{pmatrix}.
 \label{eq:op3-original-border-inverse}
\end{equation}
For the original physical face values,
$u_j=(\gamma\boldsymbol s^{\mathsf T}\boldsymbol u_U-\lambda d_j)/\delta$
and the old values increase by $\gamma\boldsymbol g u_j$.
\end{lemma}
\begin{proof}
The inverse factorization underlying
\cref{lem:op3-cluster-conditional-green} separates conditional interior
residuals from the retained coordinates, proving
\eqref{eq:op3-conditional-full-green}. Equivalently, it is the usual block
inverse of a matrix partitioned into interiors and ports. Each response
row has at most two nonzero retained coefficients.

The enlarged original matrix has blocks
$\boldsymbol M_{UU}$, $-\gamma\boldsymbol s$ and $d_j$.
Direct block elimination proves \eqref{eq:op3-original-border-inverse}
and the physical value update. For any principal active set $W$ and vector
$\boldsymbol y$, the original full degrees give
\[
 \boldsymbol y^{\mathsf T}\boldsymbol M_{WW}\boldsymbol y
 \geq \sum_{i\in W}(d_i-\gamma d_{G[W]}(i))y_i^2
 \geq\bar\alpha\sum_{i\in W}d_i y_i^2.
\]
Minimizing the enlarged quadratic over old coordinates with the new
coordinate fixed to one gives $\delta\geq\bar\alpha d_j$.
The identities are exact; this lower bound by itself is not a floating-point
stability theorem for a long sequence of inverse updates.
\end{proof}

\paragraph{Precise remaining algorithmic obligation.}
For an ordinary one-parent admission at $i$, a named trace supplies
$g_{ii}$ and the at most two nonzero coefficients of $\boldsymbol a_i$.
Then the selected old-port Green vector is
$\boldsymbol J\boldsymbol a_i$ and the full parent variance is
$g_{ii}+\boldsymbol a_i^{\mathsf T}\boldsymbol J\boldsymbol a_i$.
Equation \eqref{eq:op3-original-border-inverse} suggests a rank-one update
of the stored port inverse in $O(|P|^2)$ work, with every inverse entry
written and charged.

At a cycle birth, first promote the old active parents and newly needed
Steiner junctions using selected covariances, then border the new vertex
with all its active edges. At most one parent, the insertion-tree parent,
may be a transient port; restricting the updated inverse afterward to the
permanent retained set preserves the original partition's port bound.
Retaining every promoted marker is an alternative with a larger constant.
This step must
be reconciled with the paid repartition and with queries from the old
current piece versions. The full update implementation, exact comparisons
and total promotion/copy bound are the next falsifiable target in
\texttt{COARSE\_INVERSE\_UPDATE\_PROBE.md}. No improved end-to-end exponent
is asserted from these identities alone.

\paragraph{Measured.}
\texttt{cluster\_green\_queries.py} compares affine rows, conditional
variances and all covariance pairs against independent exact conditional
inverse matrices. It includes one-port bases, all compress/rake/forget
helpers, arbitrary physical seeds, shifted source-interior clusters and
retained payload versions. Whole-hierarchy indexing and dense inverse
columns are separate reference costs. Reproducible counts and hashes are
recorded in \texttt{CLUSTER\_GREEN\_QUERY\_AUDIT.json}. The sweep includes
256 canonical faces, 2,664 independent conditional inverse matrices,
29,408 affine/variance checks and 123,844 covariance pairs, with no mismatch.

\section[Maintained coarse inverse]{A maintained coarse inverse with paid port promotion}
\label{sec:op3-maintained-inverse}

\paragraph{Status.}
This section completes the construction proposed after
\cref{lem:op3-green-promotion-border}. Its correctness and word-work bound
are \statusproved{} as a draft awaiting independent review. The exact
update transaction is implemented and audited. The online balanced
hierarchy remains the explicit published \statussource{} import of
\cref{thm:op3-local-top-tree}; rebuilding it at every checkpoint in the
audit is separately charged reference work. General OP3 remains
\statusopen{}.

\begin{lemma}[Monotone permanent ports and one temporary parent; \statusproved]
\label{lem:op3-monotone-coarse-ports}
Use the permanent insertion tree and marked partition of
\cref{lem:op3-marked-tree-partition}. The permanent retained sets are
nested as admissions proceed. Suppose a new vertex $j$ has $a\geq2$
active parents, and let $P^+$ be the permanent retained set after its
admission. Every active parent except possibly the chosen insertion-tree
parent belongs to $P^+$. Thus
\[
 Z=(P^+\setminus(P\cup\{j\}))\cup(N(j)\cap U\setminus P)
\]
contains at most one vertex outside $P^+$. Across the entire run, there
are at most $5r$ occurrences in these promotion sets, where $r$ is final
active cycle rank. At every such event, $|P\cup Z\cup\{j\}|\leq4r+1$.
\end{lemma}
\begin{proof}
The insertion tree grows only by attaching a new leaf; old tree paths
do not change. Endpoints of extra edges are never unmarked. Consequently
the old marked Steiner subtree is contained in the new one, and an old
degree-at-least-three Steiner vertex remains branching. The seed stays
retained, including through the initial tree regime. Hence $P\subseteq P^+$.

The new vertex and every non-tree active parent are endpoints of newly
created extra edges and therefore are marked. Only its one tree parent
can fail to be retained. Each permanent old vertex is promoted at most
once, so there are at most $4r$ such promotions. The temporary parent can
be promoted again at later events, but each event increases cycle rank by
at least one. There are at most $r$ events and therefore at most $r$
temporary occurrences. Finally $P\cup Z\cup\{j\}$ is $P^+$ plus at most
that one temporary parent, proving the dimension bound.
\end{proof}

\paragraph{Maintained state.}
At a checkpoint keep the current dense matrix
\[
 \boldsymbol J=(\boldsymbol M_{UU}^{-1})_{PP}
\]
and the exact physical values $\boldsymbol u_P$. Keep the current piece
records and their one- or two-port values
$\boldsymbol w_{P_B}=\boldsymbol u_{P_B}-C\boldsymbol1$, with
$C=1/\bar\alpha$. For a queried old vertex $i$, extend its conditional
row $\boldsymbol a_i$ by zero to the full port namespace. Then
\[
 u_i=z_i+\boldsymbol a_i^{\mathsf T}
                    (\boldsymbol u_P-C\boldsymbol1)+C.
\]
If $i\in P$, use its basis row and zero conditional variance directly.
The singleton state is $P=\{v\}$, $J_{vv}=1/d_v$ and
$u_v=(1-\lambda d_v)/d_v$; the empty case is certified before allocation.

\begin{lemma}[Ordinary leaf update without a coarse solve; \statusproved]
\label{lem:op3-ordinary-coarse-inverse}
Suppose $j$ has exactly one active parent $i$. Query $u_i$, its sparse
conditional row $\boldsymbol a_i$ and conditional variance $g_{ii}$ from
the old current piece. Put
\[
 \boldsymbol z=\boldsymbol J\boldsymbol a_i,\qquad
 \nu=g_{ii}+\boldsymbol a_i^{\mathsf T}\boldsymbol z,\qquad
 \delta=d_j-\gamma^2\nu,
 \qquad u_j^+=\frac{\gamma u_i-\lambda d_j}{\delta}.
\]
For a strictly positive original gate, $u_j^+>0$ and
$\delta\geq\bar\alpha d_j$. With the same retained set $P$, the updates
\begin{equation}
 \boldsymbol J^+=\boldsymbol J+
      \frac{\gamma^2}{\delta}\boldsymbol z\boldsymbol z^{\mathsf T},
 \qquad
 \boldsymbol u_P^+=\boldsymbol u_P+\gamma\boldsymbol z u_j^+
 \label{eq:op3-ordinary-inverse-update}
\end{equation}
are exact. Their work is $O((1+|P|)^2+L^2)$, including all dense reads,
writes and allocations and the named response query.
\end{lemma}
\begin{proof}
By \eqref{eq:op3-conditional-full-green}, $\nu$ is the full parent inverse
diagonal and $\boldsymbol z$ is its full inverse column restricted to $P$.
Restrict \eqref{eq:op3-original-border-inverse} and its physical mean
update to $P$. This proves the identities and positivity assertions.
The row $\boldsymbol a_i$ has at most two nonzero coefficients, so
$\boldsymbol z$ requires at most two columns of $\boldsymbol J$ and
$O(1+|P|)$ field work. Writing the rank-one update costs
$O((1+|P|)^2)$. Store its new matrix, release the old current inverse,
and retain no dense inverse history. The unchanged port set is valid
after the leaf is appended to its parent's home piece. Restore that
piece's source exposure before the next reporter query.
\end{proof}

\begin{lemma}[Cycle birth with old-version promotion and inverse restriction; \statusproved]
\label{lem:op3-cycle-inverse-transaction}
At a multi-parent admission, the following transaction gives the exact
new coarse inverse and physical port values without a coarse factorization.
\begin{enumerate}
 \item Preserve the old current piece records, their port values, home-piece
 map, $\boldsymbol J$ and $\boldsymbol u_P$. After the positive gate has
 been certified, the newly admitted row may be scanned and the new marked
 partition computed from cached tree edges. Queries in the following step
 still use the preserved old versions.
 \item Form $Z$ from \cref{lem:op3-monotone-coarse-ports}. For each vertex
 in $Z$, obtain its old conditional response and noise trace from its old
 home piece. On $Q=P\cup Z$, assemble
 $\boldsymbol J_Q=(\boldsymbol M_{UU}^{-1})_{QQ}$ using
 \eqref{eq:op3-conditional-full-green}, always with the same old
 $\boldsymbol J$ as baseline. Also evaluate the old physical means on $Q$.
 \item All active parents of $j$ are now in $Q$. With their indicator
 $\boldsymbol s_Q$, put
 $\boldsymbol z=\boldsymbol J_Q\boldsymbol s_Q$,
 $\delta=d_j-\gamma^2\boldsymbol s_Q^{\mathsf T}\boldsymbol z$ and
 $u_j^+=(\gamma\boldsymbol s_Q^{\mathsf T}\boldsymbol u_Q-
 \lambda d_j)/\delta$. Border $\boldsymbol J_Q$ by
 \eqref{eq:op3-original-border-inverse} and update its means. This includes
 every active edge incident to $j$.
 \item Select the principal submatrix of the bordered \emph{inverse}
 indexed by $P^+$ and select the corresponding means. Drop the possible
 temporary parent, releasing all temporary dense matrices. Complete the
 already paid rebuilding of the new piece hierarchies and their home maps.
\end{enumerate}
The resulting state is
$((\boldsymbol M_{U^+U^+}^{-1})_{P^+P^+},\boldsymbol u_{P^+})$.
\end{lemma}
\begin{proof}
All responses in step 2 refer to the same old matrix. Within a piece,
their conditional cross entries are obtained by matching old elimination
records. Across different piece interiors, those entries are zero because
the interiors are disjoint and all extra edges meet retained ports.
Entries involving old ports use basis rows and zero conditional covariance.
The fixed-baseline Green identity therefore supplies precisely the selected
old full inverse, independently of the new partition.

Since every old neighbor of $j$ is in $Q$, eliminating $U\setminus Q$
does not change the new vertex's diagonal $d_j$ or its off-diagonal row
$-\gamma\boldsymbol s_Q$. The Schur matrix on $Q\cup\{j\}$ has old
block $\boldsymbol J_Q^{-1}$ and this exact border. Thus step 3 is the
block inverse identity, and $\delta\geq\bar\alpha d_j$ as before.
Step 4 selects the desired entries of the full new inverse. Selecting
instead a principal submatrix of the coarse \emph{matrix} would omit
the elimination of the temporary parent and is not this algorithm.
All maintained physical means follow the same original block solve.
Adding the newly admitted original row never changes the original degree
or fixed load of any old vertex, so preserving old numerical records is
sufficient while the new metadata is constructed.
\end{proof}

\begin{theorem}[Quadratic rank dependence; \statusproved]
\label{thm:op3-maintained-coarse-inverse}
Use the model, support, ranks and notation of
\cref{thm:op3-local-cycle-rank}. Replacing its checkpoint coarse solves
by \cref{lem:op3-ordinary-coarse-inverse,lem:op3-cycle-inverse-transaction}
returns the same exact obstacle solution, with total word work
\begin{equation}
 O\!\left(1+V\bigl[(1+r)L^4+(1+r)^2+qL^2\bigr]\right)
 \label{eq:op3-maintained-inverse-work}
\end{equation}
and space
\begin{equation}
 O\!\left(1+(1+r)VL^3+(1+r)^2\right).
 \label{eq:op3-maintained-inverse-space}
\end{equation}
Here all historical cluster/hull versions and their one/two-port values
may be retained, but only the current coarse inverse is kept. With
$\lambda=\epsappr/2$, this gives original ACL work
$\widetilde O((1+r^2+q)/\epsappr)$; with $\lambda=\rho$, it gives the
exact RPPR solution in $\widetilde O((1+r^2+q)/\rho)$ work.
\end{theorem}
\begin{proof}
The two update lemmas inductively preserve the exact inverse and physical
means. The component ordinary reporters and the exceptional original-gate
queries are unchanged. Their complete sign certificates therefore retain
the positivity, monotonicity, support locality, exact stopping, final
recovery and residual conclusions of \cref{thm:op3-local-cycle-rank}.

There are at most $|U|$ admissions. Every ordinary inverse update costs
$O((1+r)^2+L^2)$. At a cycle birth, copying the old home map is included
in the existing $O(VL^4)$ rebuild charge. Let $t_b$ be the number of
old promotion occurrences at that birth. Their parent walks and trace
storage cost $O(t_bL^2)$, and their selected conditional pair entries cost
$O(t_b^2L^2)$. By \cref{lem:op3-monotone-coarse-ports},
$\sum_b t_b\leq5r$, so these costs total $O(r^2L^2)$, which is included
in $O((1+r)VL^4)$ because $r\leq V$.

Old-to-new inverse entries use at most two old entries per pair; new-to-new
entries use at most four, plus the conditional covariance. Copying the old
block, initializing the augmented/bordered matrices, writing the border,
and restricting/reordering the final inverse cost $O((1+r)^2)$ at each
birth. Multiplying by the parent indicator costs $O((1+r)a)$ and
$a\leq4r+1$ there. Even retaining a simpler implementation that copies a
dense matrix separately for each promotion adds at most $O(r^3)$ work,
covered by $O(V(1+r)^2)$; the implemented variant fills one augmented
matrix per birth. All comparison maps, sorts and temporary port lists
fit the existing rebuild or trace charge.

The hierarchy construction, payload changes, output and original graph
access retain their previously paid $O((1+r)VL^4)$ work. At most
$O(q)$ exceptional-parent incidences are queried at a checkpoint, giving
$O(qVL^2)$ work including quiet queries. The one/two-port reporters cost
$O((1+r)VL^2)$, already covered. This proves the work bound.

The current inverse and a constant number of temporary promotion, border
and restriction matrices occupy $O((1+r)^2)$ words. Temporary response
traces require $O(rL)$ words. The source structures and retained application
versions occupy $O((1+r)VL^3)$, as before. Old coarse inverses are released;
saving them at every checkpoint would introduce a further $O(Vr^2)$ space
term and is excluded from \eqref{eq:op3-maintained-inverse-space}.
Finally $\lambda\operatorname{vol}(U)<1$ and $V\leq2\operatorname{vol}(U)$
give the stated ACL and RPPR consequences in their separate accuracy
namespaces.
\end{proof}

\paragraph{A sharper trajectory ledger.}
Let $b$ be the number of cycle-birth admissions and $p=|P_{\rm final}|$.
The same proof gives
\[
 O\!\left(1+V\bigl[(1+b)L^4+(1+p)^2+qL^2\bigr]\right).
\]
Indeed $b\leq p-1$, because each cycle-birth vertex is permanently marked;
there are at most $p+b\leq2p$ old promotion occurrences. Their total
$O(p^2L^2)$ covariance work is covered by $O(V(1+p)^2)$, since
$\log^2(2+V)=O(V)$ for $V\geq1$. Rebuilding happens only at the $b$
births, while all dense matrices have dimension at most $p+1$.
Here $p\leq\min\{|U|,4r\}$ for $r\geq1$, and $b\leq\min\{|U|-1,r\}$.
This form can be smaller when many cycles appear in the same admission;
unlike final $r$ and $q$, the selected partition and its $p,b$ can depend
on the admission and spanning-parent policy.

\paragraph{Scope and the next question.}
Fixed revealed cycle rank again meets OP3's target scale on that structural
class. The sufficient structural condition for the separate OP2 target
is now $1+r^2+q=O(\alpha^{-1/2})$, with logarithms suppressed. Neither
statement removes rank dependence on arbitrary graphs. Dense inverse
writes and repeated quiet exceptional gates are real costs of this
representation; they are not lower bounds for all local algorithms.
An approximate or compressed coarse response must still provide a paid
original-coordinate stopping certificate. The positive pivot bound alone
does not establish finite-precision stability.

\paragraph{Measured implementation contract.}
\texttt{maintained\_coarse\_inverse.py} implements the complete update
transaction. No coarse solve or full-face recovery feeds its admission
decisions. A separate validator forms the Schur matrix $\boldsymbol K$
from the original active matrix and checks
$\boldsymbol J\boldsymbol K=\boldsymbol I$ exactly, as well as physical
port values, all original gates, final obstacle values and ACL residuals.
Reference-only hierarchy rebuilding and whole-hierarchy parent indexing
are separately labeled. Every explicit inverse initialization, read,
write and restriction copy is counted. Old component versions are
rechecked using their own saved port values; dense inverse histories are
not retained. The triangle at $\alpha=1/3$, $\lambda=1/20$ exercises
the temporary-parent case and ends with ports $(0,2)$ and inverse
$\left(\begin{smallmatrix}3/5&1/5\\1/5&3/5\end{smallmatrix}\right)$.
Complete reproducible sweep counts and source hashes are in
\texttt{MAINTAINED\_COARSE\_INVERSE\_AUDIT.json}.

\section[Coarse publication producer]{A coarse response producer without repeated shared-candidate scans}
\label{sec:op3-coarse-publications}

\paragraph{Status and accuracy.}
This section composes the geometric recipient of
\cref{thm:op3-geometric-recipient} with the maintained coarse response.
The construction and its cost are \statusproved{} as a draft awaiting
independent review. Its state machine is implemented and exactly audited;
fast online hierarchy balancing remains the explicit \statussource{}
import. This construction returns an original ACL approximation and may
stop before the exact $\lambda$-obstacle optimum. Consequently the exact
RPPR conclusion of \cref{thm:op3-maintained-coarse-inverse} does not carry
over to this different stopping rule. Graph-uniform OP3 remains
\statusopen{}.

\paragraph{Publication rows and local recipient.}
Fix $0<\alpha<1$ and $0<\epsappr<1$. Use
\[
 \lambda=\frac{\epsappr}{2},\quad
 h=\frac{\epsappr}{16\gamma},\quad
 \eta=\frac54,\quad
 \vartheta=\frac{11\epsappr}{20},\quad C=\frac1{\bar\alpha}.
\]
Here $\eta$ is the publication ratio; $q$ continues to denote revealed
cycle rank. Maintain the exact physical active-face response implicitly
and a lower publication $\ell_i$, initially zero, for every admitted vertex.
Place the single row
\begin{equation}
 u_i-(\eta\ell_i+h)
   =w_i-(\eta\ell_i+h-C)
 \label{eq:op3-producer-home-row}
\end{equation}
on $i$'s permanent insertion home edge. The singleton seed row is queried
directly before its first edge exists. These rows replace the ordinary
boundary-heap minima; they refer to active vertices, not candidates.

Query every current piece's one/two-port reporter. If a returned sign is
positive, find that active coordinate by its named point query and set
$\ell_i=u_i$. Update its one home-edge payload. Along the cached original
row of $i$, deliver $\gamma(\ell_i^{\rm new}-\ell_i^{\rm old})$ to every
still-inactive neighbor's signal
$L_j=\gamma\sum_{i\in U\cap N(j)}\ell_i$. Every incidence read is paid,
including those whose recipient is already active and is skipped. A
candidate is queued once when $L_j>\vartheta d_j$. Exhaust due publications
before popping the next queued admission. If the producer is quiet and
the queue is empty, recover the current physical face once and return
$\bar\alpha\boldsymbol D\boldsymbol u$.

\begin{lemma}[Complete due-row reporting and safe admission; \statusproved]
\label{lem:op3-coarse-publication-producer}
The piece reporters locate a due publication whenever any active vertex
satisfies $u_i>\eta\ell_i+h$. If they are all nonpositive, they certify
\[
 \ell_i\leq u_i\leq\eta\ell_i+h\qquad(i\in U).
\]
Every publication meets \cref{def:op3-publications}. Every queued admission
has a strictly positive original $\lambda$-obstacle gate, so the maintained
inverse updates remain valid. At termination,
\begin{equation}
 r_i=\lambda d_i\quad(i\in U),\qquad
 0\leq r_j\leq\frac34\epsappr d_j\quad(j\notin U).
 \label{eq:op3-produced-acl-certificate}
\end{equation}
\end{lemma}
\begin{proof}
The permanent homes put one row for each active vertex in exactly one
piece. A piece's conditional affine response is exact at the maintained
full-active port values. Its normalized hull query preserves the sign of
each row, by the reporter contract in
\cref{sec:op3-ordered-path-clusters,sec:op3-shifted-tree-clusters}.
Arbitrary signs of the scalar threshold $\eta\ell_i+h-C$ are allowed.
Thus complete piece queries find a positive row or certify all of them
nonpositive. The numerical value of a normalized reporter score is not
used as the physical coordinate; a separate named query supplies $u_i$.

The first exact publication exceeds $h$, and a later one exceeds
$\eta\ell_i^{\rm old}$. These satisfy the original $h/2$ starting value
and $11/10$ growth requirements. Active-face responses increase under
admissions and stay fixed under publication-only events, so every retained
publication remains a lower bound. Changing a reporter payload changes
neither the active matrix nor its port inverse or physical means.

On first discovery of a candidate, its lower signal is initialized to
zero. This loses no previous publication: an edge to any previously active
vertex would already have revealed the candidate when that row was first
scanned. The new active vertex starts with zero publication. Hence cached
incremental delivery maintains the stated exact sum. Candidate signals
are monotone until admission; queueing once and removing the record at
admission suffices. Since $\vartheta>\lambda$ and $L_j\leq r_j$, a queued
candidate has the strictly positive original gate required by both inverse
update lemmas. Every reached face remains positive, monotone and contained
in the unique $\lambda$-obstacle support.

At quietness the band implies
$r_j\leq\eta L_j+\gamma h d_j$. An empty ready queue gives
$L_j\leq\vartheta d_j$, and
$\eta\vartheta+\gamma h=3\epsappr/4$.
The active equations give $r_i=\lambda d_i$. This proves the claimed
original residual and is precisely the geometric recipient certificate.
\end{proof}

\paragraph{Metadata and version discipline.}
Retain every original adjacency list after its single scan. Store active
parent incidences in geometrically growing arrays: allocating capacities
$1,2,4,\ldots$ makes all initialization, copy and append words linear in
the total number of inserted incidences. Copy or reverse a list at admission
only once, charging its length to that admitted row. The previous exact-gate
audit's tuple concatenation is not used: charging its copies to $qV$ would
defeat the present improvement. No exceptional-candidate dictionary is
scanned for decisions. Maintaining a diagnostic $q$ by one increment for
each newly repeated boundary incidence has only linear total record work.

At an admission, use the same old-version promotion and original border
as \cref{lem:op3-cycle-inverse-transaction}. At a publication, change only
its one reporter payload and deliver the increment; no coarse inverse
entry or physical port mean is changed. Restore the affected piece's source
exposure before another query. Maintain only the current list of component
roots and their current port values; old full port-value snapshots need
not be kept. Historical immutable cluster/hull records may still be retained.

\begin{theorem}[An ACL publication bound; \statusproved]
\label{thm:op3-coarse-publication-acl}
The preceding deterministic algorithm returns an ACL $\epsappr$-approximation
on every finite simple connected unweighted graph with at least two
vertices, in the local degree/adjacency and exact-real word models.
For its actual nonempty final support $U$, let
\[
 V=\operatorname{cvol}(U),\quad L=\log(2+V),\quad
 K=1+\log_+\frac{32\gamma}{\bar\alpha\epsappr},
\]
let $b$ count admissions that create at least one active cycle, and let
$p$ be final permanent retained-port count. Its total word work is
\begin{equation}
 O\!\left(1+V\left[(1+b)L^4+(1+p)^2+
              K\bigl(L^4+(1+p)L^2\bigr)\right]\right).
 \label{eq:op3-coarse-publication-work}
\end{equation}
Space, retaining historical application records but only current coarse
inverse and port-value state, is
\begin{equation}
 O\!\left(1+V(1+b+K)L^3+(1+p)^2\right).
 \label{eq:op3-coarse-publication-space}
\end{equation}
Writing $r=|E(G[U])|-|U|+1$, this gives
$\widetilde O((1+r^2)/\epsappr)$ ACL work, with no revealed-cycle-rank
factor. The returned support, and hence its $r,b,p$, may depend on the
publication and admission policy.
\end{theorem}
\begin{proof}
The previous lemma gives correctness and the same support-safe mass bound
$\lambda\operatorname{vol}(U)<1$. The empty case
$\lambda d_v\geq1$ takes constant work. Otherwise there are at most $|U|$
admissions, and a final recovery reads $O(V)$ current records and writes
$|U|$ original-coordinate output values.

Each coordinate has $O(K)$ publications by
\cref{thm:op3-geometric-recipient}; hence their number is $O(|U|K)$ and
their degree-weighted delivery count is $O(VK)$. A publication refreshes
one home-edge payload in $O(L^4)$ source work and allocates $O(L^3)$ new
application records. Checking all $O(1+p)$ piece reporters costs
$O((1+p)L^2)$ per producer search, including current port-value reads or
copies and failed searches. There is one such search per publication,
plus one quiet search per reached face. Named publication coordinates
cost at most $O(L^2)$ each. These terms give the $K$ bracket in the work
bound; all dictionary and cached delivery work $O(VKL)$ fits within it.

Original graph discovery, parent buffers, local identifiers and ready
queues use $O(V)$ word operations before their comparison-dictionary
factors. The $b$ cycle-birth partitions and source rebuilds cost
$O((1+b)VL^4)$. Each admission's dense update and temporary matrix copies
cost $O((1+p)^2)$. There are at most $p+b\leq2p$ old promotion occurrences,
since every cycle-birth vertex is permanently retained and $b\leq p-1$.
Their total conditional covariance work $O(p^2L^2)$ is covered by
$O(V(1+p)^2)$, exactly as in the refined inverse ledger. These charges
account for every remaining term. No inverse update is performed at a
publication, and no repeated exceptional-candidate search remains.

The original caches, parent arrays, queues and current source structures
use $O(V)$ words. Admissions and the $b$ rebuilds allocate
$O((1+b)VL^3)$ historical application state; publications allocate
$O(VKL^3)$ more. Current roots/port values and a constant number of
temporary dense matrices fit in $O(V+(1+p)^2)$, as do promotion traces
within the displayed larger bound. Old whole port-value snapshots are
released. Keeping a fresh complete root/value snapshot at every search
would add $O(VKp)$ space and is excluded from this statement.

For $r=0$, $b=0$ and $p=1$. For $r\geq1$, $b\leq r$ and $p\leq4r$.
Combining these inequalities with $V\leq2\operatorname{vol}(U)<4/\epsappr$
gives the stated rank-dependent ACL work. All dependence on teleportation
outside the structural rank parameters is logarithmic in this exact-word
model. This is not a finite-precision stability theorem.
\end{proof}

\begin{corollary}[Tree-shaped returned support in an arbitrary graph; \statusproved]
\label{cor:op3-published-tree-support}
If the returned active support is a tree, the preceding algorithm has
$\widetilde O(1/\epsappr)$ ACL work even when the graph revealed by its
active rows has arbitrarily many cycles. More generally, fixed active
cycle rank gives that scale without a restriction on shared inactive
candidates. A sufficient policy-independent condition is bounded cycle
rank of the exact $\lambda$-obstacle positive support, since the returned
connected support is contained in it and active cycle rank is monotone.
\end{corollary}

\begin{proposition}[Shared-pair stars; \statusproved]
\label{prop:op3-shared-pair-star}
For $k\geq2$, form a star with seed center $v$ and $k$ core leaves. For
each pair of core leaves add one shared hub joined to that pair and to
$D-2$ private leaves, where
\[
 \lambda=\frac{\gamma}{4k(k+\gamma)},\quad
 \epsappr=2\lambda,\quad D=\left\lceil\frac4{\bar\alpha\lambda}\right\rceil.
\]
Both the publication policy and the previous exact-gate algorithm return
exactly the star core. Its active rank is zero, revealed rank is
$\binom{k}{2}$, scanned volume is $k(k+1)$ and charged volume is $(k+1)^2$.
The previous per-face exceptional checker makes exactly
$\binom{k+1}{3}$ exceptional gate queries. The new publication algorithm
makes none and performs only $O(k(k+1)K)$ cached incidence deliveries.
No shared-hub or private-leaf row is scanned by either algorithm.
\end{proposition}
\begin{proof}
Every core vertex has original degree $k$. Directly solving the full-core
face gives
\[
 u_v=\frac{1-\lambda k(1+\gamma)}{k-\gamma^2},\qquad
 u_{\rm leaf}=\frac{3\gamma}{4k(k-\gamma^2)}>0.
\]
Each hub's original residual is at most $2\gamma/\bar\alpha<\lambda D$;
its private leaves have zero residual. Thus this positive core is the
unique obstacle support. On the singleton face, the seed value
$(1-\lambda k)/k$ exceeds $h$ and is published. It queues every core leaf,
because
\[
 \frac{\gamma(1-\lambda k)}{k}>\frac{11}{10}\lambda k.
\]
After multiplying by $40k(k+\gamma)/\gamma$, this is
$40k+30\gamma>11k$. The queued leaves remain ready, while the hub bound
prevents every hub admission. Hence the publication policy also reaches
the entire core, independently of leaf order.

After $t$ core leaves have been admitted, exactly $\binom{t}{2}$ hubs
have two active parents. The previous checker evaluates each at every
checkpoint, including final quietness, and
$\sum_{t=0}^k\binom{t}{2}=\binom{k+1}{3}$. The new loop has no such
query, and its cached-delivery bound follows from the publication count.
This is a cost separation between these implementations, not a universal
lower bound for local algorithms.
\end{proof}

\paragraph{Measured and an exact early-stop example.}
\texttt{coarse\_publication\_solver.py} independently checks each producer
sign certificate, every boundary signal, readiness membership, all original
ACL residuals, and $\boldsymbol J\boldsymbol K=\boldsymbol I$ against
original-matrix Schur elimination at admissions. It also verifies that
publication-only events leave the inverse and physical port means unchanged.
Dynamic parent-buffer copies and every cached incidence delivery are counted.
Source hierarchy rebuilds and whole-face validators are separate reference
work. The unchanged inverse transaction is rerun in
\texttt{MAINTAINED\_COARSE\_INVERSE\_REFACTOR\_AUDIT.json}; the complete
publication sweep is recorded in \texttt{COARSE\_PUBLICATION\_AUDIT.json}.

On the path $v$--$a$--$b$--$c$, with $\alpha=1/3$ and
$\epsappr=1/10$, the publication policy stops at
\[
 \boldsymbol u=(37/35,\ 3/14,\ 0,\ 0),\qquad
 \boldsymbol r=(1/20,\ 1/10,\ 3/28,\ 0).
\]
The exact $\lambda=1/20$ obstacle response is
$(55/52,\ 14/65,\ 1/260,\ 0)$. The exterior value $r_b=3/28$ exceeds
$\lambda d_b=1/10$ but is below the readiness threshold
$\vartheta d_b=11/100$, and satisfies the original ACL certificate.
This verified distinction is why this section states an ACL result rather
than an exact-RPPR result. Removing active-rank dependence while preserving
these paid original-coordinate certificates remains the next general OP3
question.


\section{Conclusions for the project conjecture}
\label{sec:conclusion}

The repeated factor can be eliminated, but the quantifiers matter.

\begin{itemize}
\item It is eliminated completely and rigorously on endpoint paths by
\cref{thm:path-linear}.  This is a new exact active-set realization, not a
heuristic warm start.
\item It cannot be eliminated inside the literal round-materialized source
implementation; \cref{thm:path-materialization} already blocks that route on
paths.
\item On arbitrary graphs, the energy telescope is useful but incomplete.
The missing object is an implicit dynamic Schur/SDD representation coupled to
a boundary-violation dictionary.
\item The new \cref{thm:op3-persistent-tree} removes repeated full curve
materialization on a supplied tree; it does not yet give support-local
discovery. In contrast, \cref{thm:op3-local-cycle,cor:op3-local-cycle-acl}
include local discovery on cycles with arbitrary leaf attachments. Their
fixed interior records rely on an attachment-settling property absent from
general cyclic graphs. The bounded-window extension in
\cref{thm:op3-bounded-attachments} handles tree attachments with a
polynomial attachment-size factor and a cycle seed. These additions are
proof drafts awaiting review.
\item \Cref{thm:op3-two-port-frontier} removes inactive attachment-size
dependence when at most one nonseed branching vertex belongs to the positive
support. It retains two vertices and uses individually updated planar
boundary records. Its fully charged implementation remains a structural
result. \Cref{thm:op3-branch-core-flux} permits a growing branching core by
using scalar physical-flux events, with an explicit quadratic core-size
factor for its dense inverse backend.
\item \Cref{thm:op3-live-core-peeling} replaces the permanent branching
count by maximum live-core size. Its physical groups move without copying
all their original incidences. Depth-first order makes that core small on
completed binary trees, but \cref{prop:op3-live-core-interior-tree} proves
that an interior-support tree can block every first peeling step. A general
response producer must therefore handle unfinished branching cores as well.
\item \Cref{thm:op3-local-top-tree} removes the core factor on arbitrary
trees using persistent projective chains and the checked top-tree source
interface. It charges online graph discovery and gives exact obstacle/ACL
work $\widetilde O(1/\epsappr)$. Its callbacks are implemented; the
published balancing structure is an explicit source dependency.
\item \Cref{thm:op3-local-unicyclic} gives the same bound on arbitrary
unicyclic graphs. A unique exceptional two-parent gate, two named-coordinate
queries and one charged re-rooting remove the earlier seed and attachment
restrictions. The following result handles multiple cycles with explicit
rank-dependent work.
\item \Cref{thm:op3-maintained-coarse-inverse} gives a complete local algorithm on
arbitrary graphs with explicit work
$\widetilde O((1+r^2+q)/\epsappr)$. It retains only a coarse cycle system
and uses independent two-port reporters, while charging shared inactive
candidates separately. The selected Green identities in
\cref{sec:op3-cluster-green} support the fully charged maintained-inverse
transaction in \cref{sec:op3-maintained-inverse}. Removing rank dependence
on arbitrary graphs remains open.
\item \Cref{thm:op3-coarse-publication-acl} uses the coarse response as a
geometric-publication producer and removes repeated shared-candidate scans.
Its ACL bound is $\widetilde O((1+r^2)/\epsappr)$, depending on active
cycle rank only. Tree-shaped returned support therefore gives the target
scale even in an arbitrary ambient graph with many shared inactive
candidates. This different stopping policy is not an exact RPPR solver.
\item If an output-linear arbitrary-graph interface is found, the resulting
$\widetilde O(1/\rho)$ RPPR bound would beat, rather than merely match, the
project's $\widetilde O(1/(\rho\sqrt\alpha))$ conjectured scale.  This does not
contradict the project's existing lower bounds, which are tied to narrower
persistent-support or repeated-row access models.
\end{itemize}

\begin{thebibliography}{9}

\bibitem[Alstrup et~al.(2005)Alstrup, Holm, de~Lichtenberg, and Thorup]{alstrup2005toptrees}
Stephen Alstrup, Jacob Holm, Kristian de~Lichtenberg, and Mikkel Thorup.
\newblock Maintaining information in fully-dynamic trees with top trees.
\newblock \emph{ACM Transactions on Algorithms}, 1(2):243--264, 2005.
\newblock Checked: journal Theorem 2.1, pp.\ 245--248 and 259--260;
also arXiv:cs/0310065v2, Sections 2 and 6.
\newblock \url{https://arxiv.org/pdf/cs/0310065v2}.

\bibitem[Jacob and Brodal(2019)]{jacob2019dynamic}
Riko Jacob and Gerth St\o lting Brodal.
\newblock Dynamic planar convex hull.
\newblock \emph{arXiv:1902.11169v1}, 2019. Theorem 1 and Section 2.2.
\newblock \url{https://arxiv.org/pdf/1902.11169v1}.

\bibitem[Agarwal et~al.(2010)Agarwal, Phillips, and Sadri]{agarwal2010act}
Pankaj~K. Agarwal, Jeff~M. Phillips, and Bardia Sadri.
\newblock Lipschitz unimodal and isotonic regression on paths and trees.
\newblock Author manuscript, January 2, 2010. Section 3, Theorem 3.1.
\newblock \url{https://www.cs.toronto.edu/~sadri/publications/regression.pdf}.

\bibitem[Koutis et~al.(2011)Koutis, Miller, and Peng]{koutis2011nearly}
Ioannis Koutis, Gary~L. Miller, and Richard Peng.
\newblock A nearly-$m\log n$ time solver for SDD linear systems.
\newblock In \emph{FOCS}, pages 590--598, 2011.

\bibitem[Wei and Yang(2026)]{wei2026simple}
Zhewei Wei and Mingji Yang.
\newblock A simple active-set method for PageRank-based local graph clustering.
\newblock \emph{arXiv preprint arXiv:2608.16339v1}, 2026.

\bibitem[Fountoulakis et~al.(2019)Fountoulakis, Roosta-Khorasani, Shun,
Cheng, and Mahoney]{fountoulakis2019variational}
Kimon Fountoulakis, Farbod Roosta-Khorasani, Julian Shun, Xiang Cheng, and
Michael~W. Mahoney.
\newblock Variational perspective on local graph clustering.
\newblock \emph{Mathematical Programming}, 174(1--2):553--573, 2019.

\end{thebibliography}

\end{document}
